\documentclass[11pt]{article}
\usepackage{fontspec}
% Charter et Menlo (macOS) ; ailleurs, XCharter et DejaVu Sans Mono, dont ils dérivent.
\IfFontExistsTF{Charter}{\setmainfont{Charter}}{\setmainfont{XCharter}}
\IfFontExistsTF{Menlo}{\setmonofont{Menlo}[Scale=0.82]}{\setmonofont{DejaVu Sans Mono}[Scale=0.82]}
\usepackage{polyglossia}\setmainlanguage{french}\setotherlanguage{english}
\usepackage{microtype}
\usepackage[most]{tcolorbox}
\usepackage{marginnote}% \nonverifie, dans les encadrés
% Chaque partie commence sur une nouvelle page.
\usepackage{etoolbox}
% Numéros de partie en chiffres : les ordinaux en toutes lettres de polyglossia s'arrêtent avant la vingtième.
\usepackage{titlesec}
\titleformat{\part}[display]{\normalfont\bfseries}{\large Partie~\arabic{part}}{0.4em}{\huge}
\titlespacing*{\part}{0pt}{0pt}{2em}
\pretocmd{\part}{\clearpage}{}{}
% Les résultats principaux de l'annexe sont encadrés.
\newtcolorbox{encadre}{enhanced, breakable, colback=white, colframe=black!55, boxrule=0.6pt, arc=0pt, left=6pt, right=6pt, top=2pt, bottom=2pt, before skip=8pt, after skip=8pt}
% L’ : la paire de crénage de Charter colle l’apostrophe au L ; on rajoute une fine espace.
\XeTeXinterchartokenstate=1
\newXeTeXintercharclass\LetterL
\newXeTeXintercharclass\Apostrophe
\XeTeXcharclass`L=\LetterL
\XeTeXcharclass`'=\Apostrophe
\XeTeXcharclass"2019=\Apostrophe
\XeTeXinterchartoks\LetterL\Apostrophe={\kern0.07em}
\usepackage[a4paper,margin=27mm]{geometry}
\usepackage{booktabs,longtable,array}
\usepackage{pdflscape}% le tableau récapitulatif, en paysage
\usepackage{setspace}\setstretch{1.1}
\usepackage{xcolor}
\usepackage[hidelinks]{hyperref}
\usepackage{amsmath,amssymb,amsthm}
\usepackage{tikz-cd}
\usetikzlibrary{decorations.pathreplacing}


\newcommand{\q}[1]{«\,#1\,»}
% Les rédactions des archives Bourbaki (Archives Henri Poincaré) : une notice par item.
\newcommand{\ahp}[1]{\href{https://archives-bourbaki.ahp-numerique.fr/items/show/#1}{item~#1}}
\newcommand{\redac}[1]{n\textsuperscript{o}\,#1}
% Le fichier de la preuve vérifiée, une fois par section. Les fichiers de lean/Bourbaki/
% n'existent pas encore sous une étiquette publiée : les liens pointent vers main.
\newcommand{\leanblobmain}{https://github.com/Commutator-IO/grothendieck-project/blob/main/lean/Bourbaki}
\newcommand{\verifie}[2][\leanblobmain]{\footnote{Preuve vérifiée : \href{#1/#2.lean}{\texttt{lean/Bourbaki/\detokenize{#2}.lean}}.}}
\definecolor{frise}{HTML}{2F4DA8}
\definecolor{piste}{HTML}{B53D1D}
% La démonstration de la page, résumée, pour un énoncé non démontré en Lean : marquée dans la marge.
\newenvironment{demopage}[1][]{\par\smallskip\noindent\marginpar{\raggedright\footnotesize\textcolor{piste}{\ensuremath{\lozenge}}~non vérifiée en Lean}\textit{Démonstration de la page}\if\relax\detokenize{#1}\relax\else{} (#1)\fi.~}{\par\smallskip}
% L'état de la recherche bibliographique, en fin de section.
\newcommand{\litterature}[1]{\par\smallskip\noindent\textit{Littérature.} #1\par}
\theoremstyle{plain}
\newtheorem*{enonce}{Énoncé}
\newtheorem{lemme}{Lemme}[section]
\newtheorem{theoreme}[lemme]{Théorème}
\newtheorem{contrexemple}[lemme]{Contre-exemple}
\newtheorem{conjecture}[lemme]{Conjecture}
\newtheorem*{conjecture*}{Conjecture}
\newtheorem{corollaire}[lemme]{Corollaire}
\newtheorem*{corollaire*}{Corollaire}
\theoremstyle{definition}
\newtheorem{exemple}[lemme]{Exemple}
\newtheorem*{exemple*}{Exemple}
\newtheorem{methode}[lemme]{Méthode}
\newtheorem*{methode*}{Méthode}
\newtheorem*{definition*}{Définition}
\theoremstyle{plain}
\newtheorem*{theoreme*}{Théorème}
\newtheorem*{contrexemple*}{Contre-exemple}
\theoremstyle{definition}
\newtheorem*{definition}{Définition}
\theoremstyle{remark}
\newtheorem*{remarque}{Remarque}
\raggedbottom\emergencystretch=1.5em
% Long Lean names (\texttt{A.b\_c\_d}) may break after an underscore, without a hyphen.
\let\bkunderscore\_
\renewcommand{\_}{\bkunderscore\allowbreak}

\title{\Large Rédactions Bourbaki : énoncés démontrés en Lean}
\author{Michel Hua\\ \small Chercheur indépendant, Paris\\ \small \texttt{michel@commutator.io}}
\date{\small Octobre 2026}

\begin{document}
\maketitle

\begin{abstract}
\noindent Les archives Bourbaki conservent les rédactions, projets de chapitres que les membres du groupe se faisaient circuler et discutaient en congrès. Ce document lit des énoncés précis dans vingt-deux de ces rédactions, de 1950 à 1974, en analyse fonctionnelle, topologie, théorie des ensembles, algèbre, géométrie algébrique, théorie des catégories et algèbre homologique, et les démontre. Chaque démonstration suit une preuve vérifiée dans l'assistant de preuve Lean~4 \cite{lean4}, avec sa bibliothèque mathlib \cite{mathlib2020}, et elle est rédigée pour être lue sans Lean. Pour chaque énoncé, le document dit d'où il vient, ce qu'il veut dire, comment on le démontre, ce qui n'est pas formalisé et ce que l'on sait de sa présence dans la littérature. Plusieurs énoncés tapés sont faux, et des contre-exemples le montrent.
\end{abstract}

\tableofcontents

\section*{Introduction}\addcontentsline{toc}{section}{Introduction}

\emph{L'archive.} Les Archives Henri Poincaré (Nancy) publient en ligne les archives de Nicolas Bourbaki, avec l'Association des collaborateurs de Nicolas Bourbaki : 891 documents, environ 42\,850 pages, sur le site \url{https://archives-bourbaki.ahp-numerique.fr}. L'essentiel est fait de rédactions numérotées, c'est-à-dire de projets de chapitres des \emph{Éléments de mathématique}, tapés par un membre, ronéotés et discutés en congrès, puis repris par un autre. Chaque document a une notice, qui donne son numéro, son titre, son auteur présumé et sa date quand on les connaît. Ce document cite chaque rédaction par son numéro, son titre et son item sur le site, et reprend l'auteur et la date tels que la notice les donne. Une notice qui ne nomme personne laisse la rédaction anonyme ici.

\emph{La règle sur les transcriptions.} Les rédactions lues ici ont été transcrites en \LaTeX{} pour ce travail. Ces transcriptions restent privées : elles ne sont ni publiées ni versées au dépôt du projet tant que les Archives Henri Poincaré et l'Association des collaborateurs de Nicolas Bourbaki n'en ont pas donné l'accord. Le document ne cite donc presque rien du texte des rédactions, au plus quelques mots ; il énonce chaque résultat dans ses propres termes et renvoie à la page du fac-similé, que chacun peut consulter sur le site des archives.

\emph{Le choix des énoncés.} Deux passes ont fourni les énoncés. La première a transcrit les documents que l'onglet Bourbaki du site du projet signale : les rédactions \redac{128}, 288, 355, 380, 384, 387, 397, 453, 457, 551, 582 et~596 en viennent. Trois d'entre elles, 355, 384 et~397, n'ont fourni que des énoncés retrouvés dans la littérature ; ils sont formalisés parce que mathlib ne les a pas. La seconde, plus récente, ne lit que les rédactions des livres que Bourbaki n'a pas publiés, là où des mathématiques inédites ont le plus de chances de se trouver : environ 5\,900 pages en 105 rédactions, sur la géométrie algébrique, les catégories et l'algèbre homologique, les variétés différentielles, la topologie élémentaire, la géométrie différentielle et d'autres sujets. Chaque rédaction y est parcourue à la recherche d'énoncés précis (définitions, théorèmes, contre-exemples, questions posées). Chaque énoncé fait l'objet d'une recherche bibliographique ; s'il est retrouvé dans un texte publié, la rédaction est marquée \q{retrouvée} et l'on passe. Seuls les énoncés précis qui n'ont pas été retrouvés, ou qui sont faux tels qu'ils sont tapés, sont formalisés. Les rédactions \redac{197}, 205, 239, 321, 402, 438, 448, 461, 469 et~476 viennent de cette seconde passe. Le relevé marque deux d'entre elles, \redac{321} et~402, comme retrouvées dans la littérature ; leurs sections le disent, avec ce que la formalisation y ajoute.

\emph{Les garanties de Lean.} Chaque section cite en note de son titre le fichier Lean qui en contient les preuves, dans la bibliothèque \texttt{lean/Bourbaki/} du projet, un fichier par rédaction. Lean vérifie mécaniquement qu'un énoncé tient sous les hypothèses écrites. Il ne garantit pas que l'énoncé rend fidèlement la page, ni qu'il est nouveau. Quand la démonstration rédigée suit une autre voie que la rédaction, la section le dit. Une démonstration que l'on résume d'après la page sans la vérifier en Lean est marquée dans la marge (\textcolor{piste}{\ensuremath{\lozenge}}~\emph{non vérifiée en Lean}).

\emph{L'organisation.} Les sections sont rangées en quatre parties, par thème, et dans chaque partie par date. Chacune présente la rédaction, explique les notions dont l'énoncé a besoin, donne l'énoncé avec son intuition et sa démonstration, puis la liste de ce qui n'est pas démontré et un paragraphe \emph{Littérature}. Ce dernier distingue les énoncés retrouvés dans la littérature et les pistes candidates, que la recherche n'a pas retrouvées dans les sources lues. Une piste candidate n'est pas un résultat nouveau établi : c'est un énoncé dont l'absence de la littérature reste à confirmer, et le paragraphe dit ce qui a été cherché et ce qui reste à lire. Les résultats principaux sont encadrés. Un énoncé que la page fait porte le nom de Bourbaki et le numéro de la rédaction ; une forme que l'édition a ajoutée (généralisation, correction, reformulation) se dit énoncé de l'édition.

\section*{Tableau récapitulatif}\addcontentsline{toc}{section}{Tableau récapitulatif}

Une ligne par rédaction, dans l'ordre des parties : le numéro, l'énoncé démontré en quelques mots, le verdict, le fichier Lean et la section. \emph{Confirmé tel que lu} : l'énoncé tient sous les hypothèses de la page ; \emph{plus général} : il tient sous des hypothèses plus faibles ; \emph{forme de l'édition} : l'édition démontre une reformulation équivalente ; \emph{partiel} : une partie seulement est démontrée, ou dans un modèle ; \emph{faux sur la page} : un contre-exemple réfute une affirmation tapée ; \emph{question tranchée} : la page pose une question, et la réponse est démontrée ; \emph{déjà dans mathlib} : le fichier ne fait que rapporter à la page des énoncés que mathlib contient.

\begin{center}
{\small\setlength{\tabcolsep}{4pt}
\begin{longtable}{@{}>{\raggedright\arraybackslash}p{1.2cm}>{\raggedright\arraybackslash}p{6.3cm}>{\raggedright\arraybackslash}p{3.6cm}>{\raggedright\arraybackslash}p{2.6cm}r@{}}
\toprule
N\textsuperscript{o} & Énoncé & Verdict & Fichier Lean & \S\\
\midrule
\endhead
\multicolumn{5}{@{}l}{\textit{Analyse et topologie}}\\[1pt]
128 & Espaces bornologiques et tonnelés & déjà dans mathlib & \href{\leanblobmain/Redaction128.lean}{\texttt{Redaction128}} & \ref{sec:bourbaki-128}\\
239 & Un faisceau mou qui n'est pas fin & question tranchée (non, contre-exemple) & \href{\leanblobmain/Redaction239.lean}{\texttt{Redaction239}} & \ref{sec:bourbaki-239}\\
288 & La loi exponentielle & plus général & \href{\leanblobmain/Redaction288.lean}{\texttt{Redaction288}} & \ref{sec:bourbaki-288}\\
402 & Distances localisables & confirmé tel que lu ; plus général & \href{\leanblobmain/Redaction402.lean}{\texttt{Redaction402}} & \ref{sec:bourbaki-402}\\
205 & Un critère de revêtement & faux sur la page (contre-exemple) & \href{\leanblobmain/Redaction205.lean}{\texttt{Redaction205}} & \ref{sec:bourbaki-205}\\
380 & Pôles de la résolvante, spectre sans unité & faux sur la page (contre-exemples) & \href{\leanblobmain/Redaction380.lean}{\texttt{Redaction380}} & \ref{sec:bourbaki-380}\\
384 & Le covolume n'est pas continu & retrouvé ; exemple de l'exercice & \href{\leanblobmain/Redaction384.lean}{\texttt{Redaction384}} & \ref{sec:bourbaki-384}\\
551 & Pseudo-covariances ; théorème~6 & faux sur la page (contre-exemple ; relevé par la rédaction) ; preuve directe de b) $\Rightarrow$ c) & \href{\leanblobmain/Redaction551.lean}{\texttt{Redaction551}} & \ref{sec:bourbaki-551}\\
\multicolumn{5}{@{}l}{\textit{Catégories et algèbre homologique}}\\[1pt]
321 & Carrés cartésiens dans une catégorie abélienne & confirmé tel que lu ; faux sur la page (coquille, contre-exemple) & \href{\leanblobmain/Redaction321.lean}{\texttt{Redaction321}} & \ref{sec:bourbaki-321}\\
438 & Une identité triangulaire suffit & confirmé tel que lu & \href{\leanblobmain/Redaction438.lean}{\texttt{Redaction438}} & \ref{sec:bourbaki-438}\\
448 & Objets libres sur un diagramme & question tranchée (non, contre-exemple) & \href{\leanblobmain/Redaction448.lean}{\texttt{Redaction448}} & \ref{sec:bourbaki-448}\\
461 & Carrés semi-cartésiens ; deux énoncés tapés & forme de l'édition ; plus général ; retrouvé pour l'essentiel (Raoult) ; faux sur la page (contre-exemples) & \href{\leanblobmain/Redaction461.lean}{\texttt{Redaction461}} & \ref{sec:bourbaki-461}\\
469 & Le foncteur constant crée les limites & confirmé tel que lu ; partiel & \href{\leanblobmain/Redaction469.lean}{\texttt{Redaction469}} & \ref{sec:bourbaki-469}\\
476 & Deux exercices sur les graphes & faux sur la page (contre-exemple) ; partiel & \href{\leanblobmain/Redaction476.lean}{\texttt{Redaction476}} & \ref{sec:bourbaki-476}\\
\multicolumn{5}{@{}l}{\textit{Théorie des ensembles}}\\[1pt]
453 & Les univers forment une chaîne & partiel & \href{\leanblobmain/Redaction453.lean}{\texttt{Redaction453}} & \ref{sec:bourbaki-453}\\
\multicolumn{5}{@{}l}{\textit{Algèbre}}\\[1pt]
197 & La courbe $y^2 = P(x)$ en caractéristique~$2$ & faux sur la page (contre-exemple) & \href{\leanblobmain/Redaction197.lean}{\texttt{Redaction197}} & \ref{sec:bourbaki-197}\\
197 & Coordonnées de Chow d'un produit de cycles de points & question tranchée (non sur $\mathbf Z$, contre-exemple) ; partiel & \href{\leanblobmain/Redaction197Chow.lean}{\texttt{Redaction197}\allowbreak\texttt{Chow}} & \ref{sec:bourbaki-197-chow}\\
457 & Degré et nombre de points & faux sur la page (contre-exemple) ; confirmé tel que lu & \href{\leanblobmain/Redaction457.lean}{\texttt{Redaction457}} & \ref{sec:bourbaki-457}\\
582 & Les indices de Schur sur $\mathbf R$ ; annexes & faux sur la page (contre-exemples) ; partiel & \href{\leanblobmain/Redaction582.lean}{\texttt{Redaction582}} & \ref{sec:bourbaki-582}\\
596 & Toute classe est une différence & plus général ; partiel & \href{\leanblobmain/Redaction596.lean}{\texttt{Redaction596}} & \ref{sec:bourbaki-596}\\
596 & Trois hypothèses manquantes (§§~7 et~12) & faux sur la page (contre-exemples) & \href{\leanblobmain/Redaction596Characters.lean}{\texttt{Redaction596Characters}} & \ref{sec:bourbaki-596-hypotheses}\\
387 & L'indice des puissances $n$-ièmes des unités & faux sur la page (contre-exemple) & \href{\leanblobmain/Redaction387.lean}{\texttt{Redaction387}} & \ref{sec:bourbaki-387}\\
397 & Divisible équivaut à injectif & confirmé tel que lu ; plus général & \href{\leanblobmain/Redaction397.lean}{\texttt{Redaction397}} & \ref{sec:bourbaki-397}\\
355 & Premiers associés et extension plate & retrouvé ; sans hypothèse noethérienne & \href{\leanblobmain/Redaction355.lean}{\texttt{Redaction355}} & \ref{sec:bourbaki-355}\\
\bottomrule
\end{longtable}}
\end{center}

\part{Analyse et topologie}
\section[Rédaction Bourbaki n\textsuperscript{o}\,128 : espaces bornologiques et tonnelés]{Rédaction Bourbaki n\textsuperscript{o}\,128 : espaces bornologiques et espaces tonnelés\protect\verifie{Redaction128}}\label{sec:bourbaki-128}

La rédaction n\textsuperscript{o}\,128 des archives Bourbaki (Archives Henri Poincaré, \href{https://archives-bourbaki.ahp-numerique.fr/items/show/537}{item 537}) est l'état~4 du chapitre~IV des \emph{Espaces vectoriels topologiques}, consacré aux espaces localement convexes métrisables. Une mention d'archive sur la première page du chapitre la donne pour l'exemplaire de Laurent Schwartz, daté de mars 1950. Le chapitre est suivi de compléments sur deux théorèmes, l'un de Grothendieck, sur la complétion du dual d'un espace localement convexe, l'autre de Banach, puis de \q{Commentaires}. Ces commentaires sont écrits à la première personne et ne sont pas signés ; rien sur l'exemplaire ne nomme leur auteur, et on ne les attribue ici à personne. Sa transcription reste locale pour l'instant ; on n'en cite ici aucun passage.

\emph{Borné et continu.} Pour une application linéaire $u$ entre espaces normés, être continue et envoyer la boule unité sur une partie bornée sont une seule et même condition. Dans un espace localement convexe quelconque, une partie $B$ est \emph{bornée} si tout voisinage $V$ de $0$ l'absorbe, c'est-à-dire si $B \subset \lambda V$ pour $\lambda$ assez grand. Une application linéaire continue envoie toujours les bornés sur des bornés, mais la réciproque peut échouer. Sur l'espace de Hilbert $\ell^2$, les parties faiblement bornées sont bornées en norme (Banach--Steinhaus), de sorte que l'identité de $\ell^2$ muni de la topologie faible dans $\ell^2$ muni de la norme envoie les bornés sur des bornés ; elle n'est pourtant pas continue. La topologie faible a trop peu de voisinages pour ses bornés.

Les Commentaires isolent la propriété qui rétablit l'équivalence. Un espace localement convexe y est dit \q{bornographique} quand toute partie convexe qui absorbe tous les bornés est un voisinage de~$0$ ; on dit aujourd'hui \emph{bornologique}. La proposition~A affirme qu'un espace localement convexe métrisable est bornologique. Sa démonstration est un argument de suite : si $V$ absorbe les bornés sans être un voisinage de~$0$, on prend dans le $n$-ième voisinage $V_n$ d'une base décroissante un point $x_n$ de $V_n/n$ hors de $V$. La suite $(nx_n)$ tend vers~$0$, donc est bornée, et $V$ devrait l'absorber : pour $\lambda$ fixé et $n \geq \lambda$, on aurait $x_n \in (\lambda/n)V \subset V$, puisque $V$ est convexe et contient~$0$. La proposition~A$'$ en tire la conséquence : sur un espace bornologique, une application linéaire qui envoie les bornés sur des bornés est continue, car l'image réciproque d'un voisinage convexe de~$0$ est une partie convexe qui absorbe les bornés.

\emph{Les tonneaux.} Un \emph{tonneau} est une partie convexe, cerclée, fermée et absorbante. Dans un espace normé, la boule unité fermée en est un. Les Commentaires appellent \q{fort} un espace localement convexe où tout tonneau est un voisinage de~$0$ ; on dit aujourd'hui \emph{tonnelé}. La proposition~B affirme qu'un espace localement convexe de Baire est tonnelé. Si $V$ est un tonneau, il absorbe chaque point, donc $E$ est la réunion des fermés $nV$ ; par le théorème de Baire, l'un d'eux, et donc $V$, a un point intérieur $x$. Comme $V$ est cerclé, $-x$ est aussi intérieur, et par convexité $0 = \tfrac12 x + \tfrac12(-x)$ l'est. Les espaces tonnelés sont le cadre naturel du théorème de Banach--Steinhaus, objet de la proposition~B$'$ : une famille d'applications linéaires continues simplement bornée y est équicontinue.

\begin{encadre}
\begin{enonce}[Bourbaki, rédaction n\textsuperscript{o}\,128, Commentaires, propositions~A, A$'$ et~B, dans la forme de mathlib]
Soient $E$ et $F$ des espaces vectoriels topologiques, sur un corps valué non discret $\mathbb K$ pour $E$ et un corps valué $\mathbb K'$ pour $F$, et $\sigma : \mathbb K \to \mathbb K'$ un morphisme isométrique.
\begin{enumerate}
\item[a)] Si $E$ admet en $0$ une base dénombrable de voisinages, en particulier s'il est métrisable, toute application $\sigma$-semi-linéaire $E \to F$ qui envoie les parties bornées sur des parties bornées est continue.
\item[b)] Si $E$ est un espace de Baire, il est tonnelé au sens de mathlib : toute semi-norme semi-continue inférieurement sur $E$ est continue.
\end{enumerate}
\end{enonce}
\end{encadre}

\noindent\emph{Vérification en Lean.} Les deux énoncés sont déjà dans mathlib, et le fichier ne fait que les rapporter à la page. Le premier est \texttt{LinearMap.\allowbreak{}continuous\_of\_locally\_bounded} (\texttt{propA'}). Mathlib n'a pas la notion d'espace bornologique : les propositions~A et~A$'$ y sont réunies en un seul énoncé, sans la propriété intermédiaire, et sans hypothèse de convexité locale. Le second est l'instance \texttt{BaireSpace.\allowbreak{}instBarrelledSpace} (\texttt{propB}). Mathlib définit les espaces tonnelés par les semi-normes. Pour un espace localement convexe réel ou complexe, cette définition équivaut à celle des Commentaires, car les tonneaux sont exactement les boules unités fermées des semi-normes semi-continues inférieurement ; cette équivalence n'est pas démontrée dans le fichier. Les noms sont ceux de \href{\leanblobmain/Redaction128.lean}{\texttt{Redaction128.\allowbreak{}lean}}.

\begin{remarque}
Les Commentaires notent que les espaces (LF), limites inductives strictes de suites d'espaces de Fréchet, sont bornologiques et tonnelés sans être métrisables ni de Baire. Les deux notions dépassent donc les cas que traitent les propositions~A et~B.
\end{remarque}

Ne sont pas démontrés ici : les remarques sur les espaces (LF) ; la proposition~B$'$, que mathlib démontre pour les espaces tonnelés (\texttt{WithSeminorms.\allowbreak{}banach\_steinhaus}) mais que le fichier ne reprend pas ; l'espace $E^f$ des Commentaires et les propositions~C, C$'$ et~C$''$ sur les parties bornées ; les compléments sur les théorèmes de Grothendieck et de Banach qui précèdent les Commentaires.

\section[Rédaction Bourbaki \redac{239} : un faisceau mou qui n'est pas fin]{Rédaction Bourbaki \redac{239} : un faisceau mou qui n'est pas fin\protect\piste\protect\verifie{Redaction239}}\label{sec:bourbaki-239}

La rédaction \redac{239}, \q{Livre~VII. Topologie préhomologique. Chapitre~I. Faisceaux (état~2)}, des archives Bourbaki (Archives Henri Poincaré, \ahp{777}) est datée de juin 1956 par sa notice, qui la donne à Jean Dieudonné ; elle compte 102 pages. Elle définit les faisceaux mous et les faisceaux fins, démontre qu'un faisceau fin est mou, et le rédacteur écrit à la page~86 qu'il ne sait pas démontrer la réciproque, qu'il en doute, et qu'il n'a pas de contre-exemple. Sa transcription reste locale pour l'instant ; on n'en cite ici aucun passage.

\emph{Mou et fin.} Un faisceau $\mathcal F$ sur un espace $X$ associe à chaque ouvert $U$ un ensemble de sections $\mathcal F(U)$, avec des restrictions, de sorte que des sections qui se recollent sur un recouvrement proviennent d'une seule. Le faisceau est \emph{mou} (§~4, n\textsuperscript{o}\,3, p.~78) si une section définie au voisinage d'un fermé coïncide, près de ce fermé, avec une section globale : les sections se prolongent. Un faisceau de modules est \emph{fin} (§~4, n\textsuperscript{o}\,6, p.~85) s'il a assez d'endomorphismes pour découper : pour un fermé $F$ et un voisinage ouvert $V$ de $F$, il existe un endomorphisme du faisceau qui est l'identité près de $F$ et nul près du complémentaire de $V$. L'exemple type est le faisceau des fonctions continues réelles, où l'on multiplie par une fonction plateau ; la proposition~4 (p.~85--86) montre qu'un faisceau fin est mou, en découpant une section avec ces endomorphismes.

\emph{L'idée du contre-exemple.} Être fin demande des endomorphismes qui valent l'identité à un endroit et zéro à un autre, et un endomorphisme de groupes doit respecter la torsion. Le faisceau des applications continues à valeurs dans le cercle $\mathbf T$ a une section d'ordre~2, la constante $-1$. Son image par un endomorphisme est une application de carré~$1$, à valeurs dans $\{\pm 1\}$, donc constante sur un espace connexe : elle ne peut valoir $-1$ à un endroit et $1$ à un autre. Le même faisceau, sur $[0, 1]$, est pourtant mou : une application dans le cercle se prolonge, parce qu'on peut suivre son argument par petits morceaux.

\begin{encadre}
\begin{contrexemple}[réponse à la question de la rédaction \redac{239}, p.~86]
Sur $X = [0, 1]$, le faisceau $U \mapsto C(U, \mathbf T)$ des applications continues à valeurs dans le cercle, faisceau de groupes abéliens (c'est-à-dire de $\mathbf Z$-modules), est mou et n'est pas fin. La partie \q{pas fin} vaut sur tout espace connexe.
\end{contrexemple}
\end{encadre}

\begin{proof}
\emph{Pas fin.} Soit $(\varphi_U)$ une famille d'endomorphismes des groupes $C(U, \mathbf T)$, compatible aux restrictions, qui soit l'identité sur un voisinage ouvert d'un point $p$ et triviale sur un voisinage ouvert d'un point $q$. L'image $g = \varphi_X(-1)$ vérifie $g^2 = \varphi_X(1) = 1$ ; elle est donc à valeurs dans $\{\pm 1\}$, et constante puisque $X$ est connexe. Près de $p$, elle vaut $-1$ ; près de $q$, elle vaut $1$. C'est absurde. Pour $X = [0, 1]$, il suffit de prendre $F = \{0\}$ et $V = [0, 1[$.

\emph{Mou.} Soient $K$ un fermé de $[0, 1]$, $U$ un ouvert qui le contient et $f : [0, 1] \to \mathbf T$ une application continue sur $U$ (ses valeurs ailleurs sont quelconques). On découpe $[0, 1]$ en $M$ intervalles de longueur $1/M$, assez fins pour que tout intervalle qui rencontre le $1/M$-voisinage de $K$ soit contenu dans $U$ et que $f$ y varie de moins de $1$ ; on appelle ces intervalles les bons. Sur un bon intervalle $[x_i, x_{i+1}]$, le quotient $f(t)/f(x_i)$ reste dans le demi-plan droit, et son argument $\Delta_i(t)$ est une fonction continue ; on la prolonge à tout $[0, 1]$ en bloquant $t$ aux bornes de l'intervalle. Sur un mauvais intervalle, on fait tourner $\Delta_i$ linéairement de $0$ à l'argument de $f(x_{i+1})/f(x_i)$. L'application globale est $g = f(0)\, e^{i\sum_i \Delta_i}$. Le produit se télescope, et sur chaque bon intervalle il redonne $f$ ; comme les bons intervalles couvrent un voisinage de $K$, $g$ coïncide avec $f$ près de $K$ (figure~\ref{fig:bourbaki-239}). La construction n'utilise pas la théorie des revêtements.
\end{proof}

\begin{figure}[htbp]
\centering
\begin{tikzpicture}[every node/.style={font=\small}]
\fill[black!7] (0.77,-1.45) rectangle (3.63,1.95);
\fill[black!7] (6.27,-1.45) rectangle (8.58,1.95);
\node[black!60, above] at (2.2,1.95) {$U$};
\node[black!60, above] at (7.42,1.95) {$U$};
\draw[black!55] (0,-1.45) -- (11,-1.45);
\foreach \j in {0,...,20} { \draw[black!55] ({0.55*\j},-1.52) -- ({0.55*\j},-1.38); }
\draw[piste, line width=2.4pt] (1.65,-1.45) -- (2.75,-1.45);
\draw[piste, line width=2.4pt] (7.15,-1.45) -- (7.70,-1.45);
\node[piste, below] at (2.2,-1.52) {$K$};
\node[piste, below] at (7.42,-1.52) {$K$};
\node[below] at (0,-1.52) {$0$};
\node[below] at (11,-1.52) {$1$};
\draw[frise!55, line width=2.6pt, line join=round] (0.000,-0.330) -- (0.550,0.495);
\draw[frise!55, line width=2.6pt, line join=round] (0.550,0.495) -- (1.100,0.647);
\draw[frise!55, line width=2.6pt, line join=round] (1.100,0.647) -- (1.169,0.656) -- (1.238,0.664) -- (1.306,0.670) -- (1.375,0.675) -- (1.444,0.678) -- (1.513,0.680) -- (1.581,0.682) -- (1.650,0.682);
\draw[frise!55, line width=2.6pt, line join=round] (1.650,0.682) -- (1.719,0.681) -- (1.788,0.680) -- (1.856,0.678) -- (1.925,0.676) -- (1.994,0.673) -- (2.062,0.670) -- (2.131,0.666) -- (2.200,0.662);
\draw[frise!55, line width=2.6pt, line join=round] (2.200,0.662) -- (2.269,0.657) -- (2.338,0.653) -- (2.406,0.648) -- (2.475,0.643) -- (2.544,0.638) -- (2.612,0.632) -- (2.681,0.626) -- (2.750,0.620);
\draw[frise!55, line width=2.6pt, line join=round] (2.750,0.620) -- (2.819,0.613) -- (2.888,0.606) -- (2.956,0.598) -- (3.025,0.590) -- (3.094,0.580) -- (3.162,0.571) -- (3.231,0.560) -- (3.300,0.548);
\draw[frise!55, line width=2.6pt, line join=round] (3.300,0.548) -- (3.850,0.660);
\draw[frise!55, line width=2.6pt, line join=round] (3.850,0.660) -- (4.400,1.650);
\draw[frise!55, line width=2.6pt, line join=round] (4.400,1.650) -- (4.950,0.770);
\draw[frise!55, line width=2.6pt, line join=round] (4.950,0.770) -- (5.500,-0.550);
\draw[frise!55, line width=2.6pt, line join=round] (5.500,-0.550) -- (6.050,0.110);
\draw[frise!55, line width=2.6pt, line join=round] (6.050,0.110) -- (6.600,-1.077);
\draw[frise!55, line width=2.6pt, line join=round] (6.600,-1.077) -- (6.669,-1.080) -- (6.738,-1.078) -- (6.806,-1.073) -- (6.875,-1.064) -- (6.944,-1.050) -- (7.012,-1.033) -- (7.081,-1.011) -- (7.150,-0.986);
\draw[frise!55, line width=2.6pt, line join=round] (7.150,-0.986) -- (7.219,-0.957) -- (7.287,-0.924) -- (7.356,-0.888) -- (7.425,-0.848) -- (7.494,-0.805) -- (7.562,-0.760) -- (7.631,-0.711) -- (7.700,-0.660);
\draw[frise!55, line width=2.6pt, line join=round] (7.700,-0.660) -- (7.769,-0.607) -- (7.837,-0.551) -- (7.906,-0.494) -- (7.975,-0.435) -- (8.044,-0.375) -- (8.113,-0.314) -- (8.181,-0.253) -- (8.250,-0.191);
\draw[frise!55, line width=2.6pt, line join=round] (8.250,-0.191) -- (8.800,-0.880);
\draw[frise!55, line width=2.6pt, line join=round] (8.800,-0.880) -- (9.350,0.440);
\draw[frise!55, line width=2.6pt, line join=round] (9.350,0.440) -- (9.900,-1.045);
\draw[frise!55, line width=2.6pt, line join=round] (9.900,-1.045) -- (10.450,-0.220);
\draw[frise!55, line width=2.6pt, line join=round] (10.450,-0.220) -- (11.000,1.045);
\draw[black!45, thin] (0.000,-0.330) -- (0.275,-0.503) -- (0.550,0.495) -- (0.660,0.406) -- (0.770,0.577);
\draw[black!45, thin] (3.630,0.475) -- (3.740,-0.111) -- (3.850,0.660) -- (4.125,0.584) -- (4.400,1.650) -- (4.675,0.648) -- (4.950,0.770) -- (5.225,-1.015) -- (5.500,-0.550) -- (5.775,-1.171) -- (6.050,0.110) -- (6.160,1.277) -- (6.270,-1.011);
\draw[black!45, thin] (8.580,0.104) -- (8.690,-0.440) -- (8.800,-0.880) -- (9.075,-0.514) -- (9.350,0.440) -- (9.625,1.747) -- (9.900,-1.045) -- (10.175,0.187) -- (10.450,-0.220) -- (10.725,1.274) -- (11.000,1.045);
\draw[black, semithick] (0.770,0.577) -- (0.842,0.596) -- (0.913,0.613) -- (0.985,0.627) -- (1.056,0.640) -- (1.128,0.651) -- (1.199,0.660) -- (1.270,0.667) -- (1.342,0.672) -- (1.413,0.677) -- (1.485,0.680) -- (1.557,0.681) -- (1.628,0.682) -- (1.700,0.682) -- (1.771,0.681) -- (1.843,0.679) -- (1.914,0.676) -- (1.986,0.673) -- (2.057,0.670) -- (2.128,0.666) -- (2.200,0.662) -- (2.272,0.657) -- (2.343,0.653) -- (2.415,0.647) -- (2.486,0.642) -- (2.558,0.636) -- (2.629,0.631) -- (2.701,0.624) -- (2.772,0.618) -- (2.844,0.610) -- (2.915,0.603) -- (2.987,0.594) -- (3.058,0.585) -- (3.129,0.575) -- (3.201,0.565) -- (3.272,0.553) -- (3.344,0.540) -- (3.415,0.526) -- (3.487,0.510) -- (3.558,0.493) -- (3.630,0.475);
\draw[black, semithick] (6.270,-1.011) -- (6.328,-1.029) -- (6.385,-1.044) -- (6.443,-1.056) -- (6.501,-1.066) -- (6.559,-1.074) -- (6.616,-1.078) -- (6.674,-1.080) -- (6.732,-1.079) -- (6.790,-1.075) -- (6.847,-1.068) -- (6.905,-1.058) -- (6.963,-1.046) -- (7.021,-1.030) -- (7.079,-1.012) -- (7.136,-0.991) -- (7.194,-0.968) -- (7.252,-0.942) -- (7.309,-0.913) -- (7.367,-0.882) -- (7.425,-0.848) -- (7.483,-0.813) -- (7.540,-0.775) -- (7.598,-0.735) -- (7.656,-0.693) -- (7.714,-0.650) -- (7.772,-0.605) -- (7.829,-0.558) -- (7.887,-0.510) -- (7.945,-0.461) -- (8.003,-0.411) -- (8.060,-0.361) -- (8.118,-0.309) -- (8.176,-0.258) -- (8.233,-0.205) -- (8.291,-0.153) -- (8.349,-0.101) -- (8.407,-0.049) -- (8.465,0.003) -- (8.522,0.054) -- (8.580,0.104);
\fill[black!70] (0.000,-0.330) circle (1.5pt);
\fill[black!70] (0.550,0.495) circle (1.5pt);
\fill[black!70] (3.850,0.660) circle (1.5pt);
\fill[black!70] (4.400,1.650) circle (1.5pt);
\fill[black!70] (4.950,0.770) circle (1.5pt);
\fill[black!70] (5.500,-0.550) circle (1.5pt);
\fill[black!70] (6.050,0.110) circle (1.5pt);
\fill[black!70] (8.800,-0.880) circle (1.5pt);
\fill[black!70] (9.350,0.440) circle (1.5pt);
\fill[black!70] (9.900,-1.045) circle (1.5pt);
\fill[black!70] (10.450,-0.220) circle (1.5pt);
\fill[black!70] (11.000,1.045) circle (1.5pt);
\end{tikzpicture}
\caption{La preuve de \q{mou} sur un exemple, avec $M = 20$. En abscisse, $[0, 1]$ et ses points de subdivision ; en ordonnée, un argument, c'est-à-dire un relevé dans $\mathbf R$ des valeurs dans le cercle. Le fermé $K$ est en rouge, l'ouvert $U$ en gris. L'application $f$ est continue sur $U$ (trait noir) et quelconque ailleurs (trait gris) ; les points gris sont ses valeurs aux points de subdivision hors de $U$. L'application globale $g$ (trait bleu large) coïncide avec $f$ sur les bons intervalles, ceux qui rencontrent le $1/M$-voisinage de $K$, et tourne à vitesse constante sur les autres, d'une valeur de $f$ à la suivante par le plus court arc.}
\label{fig:bourbaki-239}
\end{figure}

\begin{sloppypar}
\noindent\emph{Vérification en Lean.} \q{Pas fin} est \texttt{not\_exists\_endomorphism}, sur tout espace préconnexe, avec la restriction \texttt{Opens.\allowbreak{}res}, et \texttt{not\_fine\_unitInterval} pour $[0, 1]$, $F = \{0\}$ et $V = [0, 1[$ ; \q{mou} est \texttt{exists\_eqOn\_nhds} et \texttt{exists\_eqOn\_nhds'}, dont l'application globale est \texttt{globalMap}, dans \href{\leanblobmain/Redaction239.lean}{\texttt{Redaction239.\allowbreak{}lean}}.
\end{sloppypar}

Ne sont pas démontrés ici : la construction du faisceau $U \mapsto C(U, \mathbf T)$ comme faisceau de mathlib, et les définitions de la rédaction (avec ses bases d'ouverts paracompacts et ses familles de fermés) ; les deux théorèmes vérifiés sont le contenu concret des deux propriétés pour ce faisceau. La proposition~4, \q{fin implique mou}, n'est pas formalisée.

\litterature{piste candidate. Le congrès avait noté la question comme ouverte (\emph{La Tribu}, n\textsuperscript{o}\,41). La question 423444 de MathOverflow (2022) demande un exemple de faisceau mou qui ne soit pas fin, et n'avait toujours pas de réponse le 1\textsuperscript{er} novembre 2026 \cite{mo423444}. Un contre-exemple imprimé n'a pas été trouvé : la recherche dans zbMATH Open (résumés et comptes rendus) et dans les titres et résumés d'arXiv n'a rien donné. R.~Godement, \emph{Topologie algébrique et théorie des faisceaux} (1958), II.3.7 et II.9, et G.~E.~Bredon, \emph{Sheaf Theory} (2\textsuperscript{e} éd., 1997), II.9 et ses exercices, n'ont pas pu être consultés légalement en ligne : la question reste à vérifier sur ces deux livres.}

\section[Rédaction Bourbaki n\textsuperscript{o}\,288 : la loi exponentielle]{Rédaction Bourbaki n\textsuperscript{o}\,288 : la loi exponentielle pour la convergence compacte\protect\verifie{Redaction288}}\label{sec:bourbaki-288}

La rédaction n\textsuperscript{o}\,288 des archives Bourbaki, datée de 1958 (Archives Henri Poincaré, \href{https://archives-bourbaki.ahp-numerique.fr/items/show/689}{item 689}), est un contre-projet pour le chapitre~X de \emph{Topologie générale}, qui traite des espaces d'applications continues. Sa transcription reste locale pour l'instant ; on n'en cite ici que quelques mots.

\emph{La loi exponentielle.} Une fonction de deux variables $f(x, y)$ peut se lire comme une famille de fonctions d'une variable : à chaque $x$, on associe la fonction partielle $y \mapsto f(x, y)$. Pour les ensembles, ce passage est une bijection entre les applications $E \times F \to G$ et les applications de $E$ dans l'ensemble des applications $F \to G$. Pour les cardinaux, il s'écrit $|G|^{|E|\cdot|F|} = \bigl(|G|^{|F|}\bigr)^{|E|}$, d'où le nom de loi exponentielle ; les informaticiens parlent de curryfication. En topologie, on demande davantage. On note $C(E, G)$ l'espace des applications continues de $E$ dans $G$, muni de la topologie de la convergence compacte (dite aussi compacte-ouverte), dont les ouverts de base sont les ensembles $\Omega(K, U)$ des $u$ telles que $u(K) \subset U$, pour $K$ compact et $U$ ouvert. La question est de savoir si la bijection envoie $C(E \times F, G)$ sur $C(E, C(F, G))$ et si elle est un homéomorphisme. L'intérêt est concret : une homotopie $E \times [0, 1] \to G$ devient un chemin dans $C(E, G)$, et une famille continue de fonctions dépendant d'un paramètre devient une seule application continue à valeurs dans un espace de fonctions.

\emph{La page.} Le corollaire~2 (PDF, p.~32--33) suppose $F$ localement compact et $E$ séparé, et affirme que la bijection naturelle induit un homéomorphisme de $C_c(E \times F, G)$ sur $C_c(E, C_c(F, G))$. La démonstration a deux temps. La continuité dans le sens direct vient de ce que l'image réciproque de $\Omega(L, \Omega(K, U))$ est $\Omega(L \times K, U)$. Pour la réciproque, on part d'un compact $M$ de $E \times F$ et d'une $f$ qui envoie $M$ dans $U$, et l'on recouvre $M$ par un nombre fini de rectangles compacts contenus dans $f^{-1}(U)$. C'est là que sert l'hypothèse sur $E$ : les projections $K$ et $L$ de $M$ sont compactes, et $K \times L$ est un compact séparé, donc normal, où l'on sait découper. Une note dactylographiée qualifie la preuve de l'ancienne rédaction de \q{canulée} et soupçonne que le corollaire est faux si $E$ n'est pas séparé ; elle en laisse la vérification à Dieudonné, comme exercice.

\emph{Deux conventions.} Chez Bourbaki, un compact est séparé par définition, et un espace localement compact aussi. Dans mathlib, \texttt{IsCompact} ne demande que la propriété de Borel--Lebesgue, et \texttt{LocallyCompactSpace} demande seulement que chaque point ait une base de voisinages compacts, sans séparation. mathlib contient déjà l'homéomorphisme \texttt{Homeomorph.\allowbreak{}curry}, mais sous l'hypothèse que les deux facteurs $X$ et $Y$ sont localement compacts. Un espace séparé qui n'est pas localement compact, comme $\mathbf Q$, n'y entre pas, alors que la page le couvre.

\begin{definition}
On dit qu'un espace $X$ vérifie (P) si, pour tout compact $K$ de $X$, tout point $x$ de $K$ et tout voisinage $V$ de $x$ dans $X$, il existe un compact $C \subset V$ qui est un voisinage de $x$ dans $K$.
\end{definition}

Tout espace localement compact vérifie (P), en prenant pour $C$ un voisinage compact de $x$ dans $X$. Un espace séparé la vérifie aussi : un compact séparé est localement compact. Plus généralement, il suffit que $X$ soit préregulier (deux points qui n'ont pas les mêmes voisinages ont des voisinages disjoints, \texttt{R1Space}), car un compact préregulier est régulier, donc localement compact.

\begin{enonce}[Corollaire~2, dans les conventions de mathlib]
Soient $X$, $Y$, $Z$ des espaces topologiques. Si $X$ vérifie (P), l'application $f \mapsto (x \mapsto f(x, \cdot))$ de $C(X \times Y, Z)$ dans $C(X, C(Y, Z))$ est continue, sans hypothèse sur $Y$ ni sur $Z$. Si de plus $Y$ est localement compact, c'est un homéomorphisme. En particulier, c'est un homéomorphisme dès que $X$ est séparé et $Y$ localement compact.
\end{enonce}

\begin{proof}
On utilise deux fois le même principe : une propriété des parties d'un compact qui est stable par réunion finie et par passage à une partie, et qui est vraie au voisinage de chaque point, est vraie du compact entier. Pour la continuité directe, on fixe $f$, un compact $K$ de $X$ et un ouvert $W$ de $C(Y, Z)$ contenant les fonctions partielles $f(x, \cdot)$ pour $x \in K$. Il suffit de montrer que chaque point $x$ de $K$ a un voisinage $C$ dans $K$ tel que, pour $f'$ assez proche de $f$, les $f'(x', \cdot)$ avec $x' \in C$ restent dans $W$. Or, si $f(x, \cdot)$ envoie un compact $L$ dans un ouvert $U$, le lemme du tube donne un voisinage $V$ de $x$ tel que $f(V \times L) \subset U$. La propriété (P) donne un compact $C \subset V$ voisinage de $x$ dans $K$, et l'ouvert $\Omega(C \times L, U)$ contient $f$. Toute $f'$ dans cet ouvert a ses fonctions partielles en les points de $C$ dans $\Omega(L, U)$, et des intersections finies de tels ouverts suffisent pour $W$.

Pour la réciproque, on fixe $g \in C(X, C(Y, Z))$, un compact $M$ de $X \times Y$ et un ouvert $U$ avec $g(x)(y) \in U$ pour $(x, y) \in M$. En un point $(x, y)$ de $M$, la locale compacité de $Y$ donne un voisinage compact $L$ de $y$ avec $g(x)(L) \subset U$, puis la continuité de $g$ un voisinage $V$ de $x$ avec $g(V) \subset \Omega(L, U)$. On applique (P) au compact $p_1(M)$, projection de $M$ sur $X$ : il existe un compact $C \subset V$ voisinage de $x$ dans $p_1(M)$. Le rectangle $C \times L$ est un voisinage de $(x, y)$ dans $M$, et l'ouvert $\Omega(C, \Omega(L, U))$ contient $g$. Le principe de recouvrement fait le reste. Les deux applications sont inverses l'une de l'autre par définition.
\end{proof}

\begin{remarque}
Sous les hypothèses de la page, $E$ séparé et $F$ localement compact séparé, tous les compacts de $E$, de $F$ et de $E \times F$ sont séparés : les deux conventions donnent les mêmes topologies, et l'énoncé démontré contient le corollaire~2. Il le renforce de deux façons : $F$ n'a pas besoin d'être séparé, et $E$ peut être seulement préregulier, ou localement compact au sens de mathlib.
\end{remarque}

\emph{Vérification en Lean.} Lean vérifie l'énoncé ci-dessus dans les conventions de mathlib, avec les deux cas où l'on établit (P) : $X$ localement compact et $X$ préregulier. Il ne dit rien de la question posée par la note.

\emph{La question de la note.} La démonstration n'utilise $E$ qu'à travers (P), au même endroit que la page. Savoir si l'on peut s'en passer reste ouvert ici, dans les deux conventions, et la différence entre elles déplace la difficulté. Dans celle de Bourbaki, la topologie de la convergence compacte ne voit que les compacts séparés, et un compact séparé vérifie toujours (P) en lui-même ; la continuité directe vaut alors pour tout $E$. Pour la réciproque, il faut appliquer (P) à la projection sur $E$ d'un compact séparé $M$ de $E \times F$, et cette projection peut ne pas être séparée quand $E$ ne l'est pas. Sur la droite à deux origines $0'$ et $0''$, par exemple, la partie $\{(0', y_1), (0'', y_2)\}$ avec $y_1 \neq y_2$ est un compact séparé, mais sa projection $\{0', 0''\}$ ne l'est pas. Dans cet exemple, tout se passe bien, car la droite à deux origines est localement compacte au sens de mathlib. Un contre-exemple devrait donc chercher un espace $E$ qui n'est ni préregulier ni localement compact ; aucun n'a été trouvé, et le travail d'Escardó et Heckmann sur les topologies des espaces de fonctions, où la question a sa place, n'a pas été consulté.

\section[Rédaction Bourbaki \redac{402} : distances localisables]{Rédaction Bourbaki \redac{402} : distances localisables et chaînes\protect\verifie{Redaction402}}\label{sec:bourbaki-402}

La rédaction \redac{402}, \q{Géométrie différentielle. Chapitre~II : Géodésiques}, des archives Bourbaki (Archives Henri Poincaré, \ahp{953}) est datée de juin 1964 par sa notice, qui la donne à Jean-Louis Koszul ; elle compte 127 pages. Son §~5 s'ouvre (p.~10--11 du tapuscrit) sur une notion de pure topologie métrique, qui servira à dire que la distance d'une variété riemannienne se calcule localement. Sa transcription reste locale pour l'instant ; on n'en cite ici aucun passage.

\emph{Chaînes et longueur.} Soient $X$ un espace muni d'une distance $d$ et $W$ un voisinage de la diagonale de $X \times X$. Une $W$-chaîne de $a$ à $b$ est une suite finie $a = x_0, x_1, \dots, x_n = b$ dont chaque pas $(x_i, x_{i+1})$ est dans $W$ ; sa longueur est $\sum_i d(x_i, x_{i+1})$. Par l'inégalité triangulaire, une chaîne n'est jamais plus courte que $d(a, b)$. La définition~1 dit que $d$ est \emph{localisable} si, pour tout voisinage $W$ de la diagonale et tous $a$, $b$, la distance $d(a, b)$ est la borne inférieure des longueurs des $W$-chaînes de $a$ à $b$.

\emph{Un exemple.} Sur le cercle unité du plan, la distance euclidienne (la corde) n'est pas localisable. Une chaîne à petits pas suit le cercle, et la somme de ses cordes est voisine de la longueur d'arc parcourue, qui dépasse la corde de $a$ à $b$ dès que $a \neq b$. La distance géodésique, longueur du plus petit arc, est au contraire localisable : les chaînes qui suivent l'arc ont exactement sa longueur. Être localisable, c'est donc dire que la distance se mesure en mettant bout à bout de petits pas, comme une longueur de chemin. Cet exemple n'est pas formalisé.

\emph{Le choix de la définition.} Le fichier écrit la borne inférieure sous la forme suivante : pour tout $\delta > 0$, il existe une $W$-chaîne de longueur $< d(a, b) + \delta$. Cette forme garde l'existence des chaînes, qu'une borne inférieure sur un ensemble vide de réels cacherait. La rédaction quantifie sur \emph{tous} les voisinages de la diagonale, et non sur les seuls voisinages uniformes $\{d < \varepsilon\}$ ; la démonstration de la proposition~2 a besoin d'un voisinage qui n'est pas uniforme. Une note manuscrite en regard propose la forme en $\varepsilon$. Les deux ne sont pas équivalentes, comme on le verra.

\begin{encadre}
\begin{enonce}[Bourbaki, rédaction \redac{402}, §~5, propositions~1 et~2, p.~11]
Soit $X$ un espace pseudo-métrique dont la distance est localisable.
\begin{enumerate}
\item[a)] (Proposition~1.) Une seconde distance localisable, qui définit la même topologie et coïncide avec $d$ sur un voisinage de la diagonale, est égale à $d$.
\item[b)] (Proposition~2.) Si $a \in A$ et $b \notin A$, alors $d(a, b)$ est la borne inférieure de $d(a, c) + d(c, b)$ pour $c$ dans la frontière de $A$.
\item[c)] (Corollaires.) On a $d(a, F) + d(F, b) \leq d(a, b)$, où $F$ est la frontière de $A$. La frontière de la boule ouverte de centre $a$ et de rayon $r > 0$ est la sphère, et son adhérence est la boule fermée. Si la boule fermée de rayon $r$ est compacte et $d(a, b) \geq r$, il existe $c$ avec $d(a, c) = r$ et $d(c, b) = d(a, b) - r$.
\end{enumerate}
L'espace est alors connexe. La rédaction le suppose connexe ; avec la définition par tous les voisinages, cette hypothèse est une conséquence et n'est pas nécessaire.
\end{enonce}
\end{encadre}

\begin{proof}
a) Soient $d_1$ et $d_2$ les deux distances, égales sur le voisinage $W$. Pour une $W$-chaîne de $a$ à $b$, on a $d_1(a, b) \leq \sum d_1(x_i, x_{i+1}) = \sum d_2(x_i, x_{i+1})$, la première inégalité par l'inégalité triangulaire, l'égalité parce que chaque pas est dans $W$. Comme $d_2$ est localisable, le membre de droite peut être rendu $< d_2(a, b) + \delta$ pour tout $\delta > 0$. Donc $d_1 \leq d_2$, et par symétrie $d_1 = d_2$.

b) Le cœur est le suivant : pour tout $\varepsilon > 0$, il existe $c$ dans la frontière avec $d(a, c) + d(c, b) < d(a, b) + \varepsilon$. Si $a$ ou $b$ est sur la frontière, on le prend pour $c$. Sinon, $a$ est intérieur à $A$ et $b$ extérieur. Soit $V$ l'ensemble des points à distance $< \varepsilon/3$ de la frontière, et $W$ la réunion de trois carrés : intérieur de $A$ fois lui-même, $V \times V$, extérieur de $A$ fois lui-même. C'est un ouvert, et il contient la diagonale, puisque tout point est intérieur, sur la frontière ou extérieur. Une $W$-chaîne de longueur $< d(a, b) + \varepsilon/3$ part de l'intérieur et finit à l'extérieur. Un pas dans le premier carré reste à l'intérieur, et un pas ne peut passer de l'intérieur à l'extérieur que par $V \times V$ ; la chaîne a donc un point $x_i$ dans $V$, à distance $< \varepsilon/3$ d'un point $c$ de la frontière. En coupant la chaîne en $x_i$, on a $d(a, x_i) + d(x_i, b) < d(a, b) + \varepsilon/3$, et deux inégalités triangulaires par $c$ donnent la conclusion. L'inégalité inverse, $d(a, b) \leq d(a, c) + d(c, b)$, est l'inégalité triangulaire.

c) Les corollaires en découlent. Pour la boule ouverte, un point $b$ de la sphère est approché par des points $c$ de la frontière de la boule ; une telle frontière est toujours contenue dans la sphère. Pour le dernier, la frontière de la boule est un fermé du compact, et la fonction $c \mapsto d(c, b)$ y atteint son minimum, qui vaut $d(a, b) - r$ par b). Enfin, si $X$ était réunion de deux ouverts fermés disjoints non vides $A$ et son complémentaire, b) donnerait un point de la frontière de $A$, qui est vide : $X$ est connexe.
\end{proof}

\emph{La note manuscrite.} Appelons $\varepsilon$-localisable une distance qui vérifie la définition pour les seuls voisinages $\{d < \varepsilon\}$. Sur $\mathbf Q$, la distance usuelle est $\varepsilon$-localisable : une progression arithmétique de $a$ à $b$ en $n$ pas a pour longueur exactement $|b - a|$, avec des pas aussi petits qu'on veut. Mais $\mathbf Q$ n'est pas connexe, donc n'est pas localisable, et la proposition~2 y est fausse : pour $A = \{x < \sqrt 2\}$, ouvert et fermé dans $\mathbf Q$, la frontière est vide. La forme proposée en regard ne suffit donc pas à la démonstration de la page.

\begin{sloppypar}
\noindent\emph{Vérification en Lean.} La définition est \texttt{IsLocalizable} (pour une fonction $d$ sur un espace topologique) et \texttt{Localizable} ; la proposition~1 est \texttt{eq\_of\_isLocalizable} et \texttt{PseudoMetricSpace.\allowbreak{}eq\_of\_localizable} ; la proposition~2 est \texttt{Localizable.\allowbreak{}exists\_mem\_frontier} et \texttt{Localizable.\allowbreak{}dist\_eq\_iInf\_frontier} ; les corollaires sont \texttt{infDist\_add\_infDist\_le}, \texttt{frontier\_ball}, \texttt{closure\_ball} et \texttt{exists\_dist\_eq\_sub} ; la connexité est \texttt{Localizable.\allowbreak{}preconnectedSpace}, et l'exemple de $\mathbf Q$ est \texttt{epsLocalizable\_rat}, \texttt{not\_localizable\_rat} et \texttt{not\_exists\_mem\_frontier\_rat}, dans \href{\leanblobmain/Redaction402.lean}{\texttt{Redaction402.\allowbreak{}lean}} ; Lean y tire l'échec de la proposition~2 de la non-connexité de $\mathbf Q$, sans exhiber la partie $A$. Pour la proposition~1, Lean n'utilise de $d_1$ et $d_2$ que l'inégalité du polygone, et non leurs autres propriétés de distance.
\end{sloppypar}

Ne sont pas démontrés ici : l'usage du §~5 dans la suite du chapitre, à savoir que la distance d'une variété riemannienne est localisable ; l'exemple du cercle.

\litterature{le relevé des rédactions la classe parmi les textes retrouvés. Les distances localisables au sens des voisinages sont proches des distances de longueur (ou intrinsèques), que Burago, Burago et Ivanov caractérisent au §~2.4 de leur cours par l'existence de milieux approchés \cite{bbi2001} ; la forme par chaînes, pour des espaces non complets, n'a pas été comparée à leur énoncé. Mathlib n'a pas de théorie des espaces de longueur.}

\section[Rédaction Bourbaki \redac{205} : un critère de revêtement]{Rédaction Bourbaki \redac{205} : relever les carrés ne suffit pas à faire un revêtement\protect\verifie{Redaction205}}\label{sec:bourbaki-205}

La rédaction \redac{205}, \q{Topologie élémentaire. Chapitre~III. Groupe fondamental et revêtements (état~1)}, des archives Bourbaki (Archives Henri Poincaré, \ahp{702} ; l'\ahp{867} en donne le même texte) n'est pas datée, et sa notice ne nomme pas d'auteur ; elle compte 46 pages. Elle se termine par une liste d'exercices, dont le deuxième (p.~45) propose une caractérisation des revêtements. Sa transcription reste locale pour l'instant ; on n'en cite ici aucun passage.

\emph{Revêtements et relèvements.} Une application continue $p : Y \to X$ est un revêtement si chaque point de $X$ a un voisinage ouvert $U$ dont l'image réciproque est une réunion disjointe d'ouverts, les \q{feuillets}, que $p$ envoie chacun homéomorphiquement sur $U$. L'exemple type est $\mathbf R \to \mathbf S^1$, $t \mapsto e^{2i\pi t}$. Un revêtement a deux propriétés : ses fibres $p^{-1}(x)$ sont discrètes, et toute application continue d'un carré $[0, 1]^2$ dans $X$ se relève continûment à $Y$, de façon unique une fois choisi le relevé d'un point. L'exercice~2, pour une base $X$ connexe, localement connexe et soumise à une condition de la rédaction notée (RU), affirme la réciproque : ces deux propriétés feraient de $p$ un revêtement.

\emph{Pourquoi cela échoue.} Les deux propriétés ne parlent que d'applications définies sur des espaces connexes, et un espace $Y$ peut être fait de morceaux qui se rapprochent sans se toucher. Une application d'un carré, connexe, reste alors dans un seul morceau, et les relèvements sont uniques ; mais deux morceaux peuvent s'approcher l'un de l'autre au-dessus d'un même point, et $p$ n'est injective sur aucun voisinage du point limite.

\begin{encadre}
\begin{contrexemple}[rédaction \redac{205}, exercice~2, p.~45]
Soit $Y \subset [0, 1] \times \mathbf R \times \mathbf R$ la réunion du segment $\{(x, 0, 0)\}$ et, pour $n \geq 1$, des graphes $\{(x, |nx - 1|, 1/n)\}$, et soit $p : Y \to [0, 1]$ la première projection. Les fibres de $p$ sont discrètes ; toute application continue d'un espace connexe dans $[0, 1]$, en particulier du carré, se relève de façon unique une fois fixé le relevé d'un point. Mais $p$ n'est pas un homéomorphisme local, et n'est donc pas un revêtement, bien que $[0, 1]$ soit connexe, localement connexe et contractile.
\end{contrexemple}
\end{encadre}

\begin{proof}
Chaque feuillet est l'image d'une section continue $x \mapsto (x, |nx - 1|, 1/n)$ de $p$, et un point de $Y$ est déterminé par sa première et sa dernière coordonnée. La dernière coordonnée, le \emph{niveau}, est une fonction continue sur $Y$ qui ne prend qu'une infinité dénombrable de valeurs ($0$ et les $1/n$). Sur un espace connexe, une fonction réelle continue a pour image un intervalle ; un intervalle dénombrable est réduit à un point. Un relevé continu d'une application d'un espace connexe a donc un niveau constant : il reste sur un feuillet. Il existe (on compose avec la section du feuillet du point imposé) et il est unique (le niveau et la projection le déterminent). Les fibres sont discrètes, car au-dessus de $x$ les points des différents feuillets ont des niveaux ou des deuxièmes coordonnées séparés. Enfin, au-dessus de $1/n$, le feuillet $n$ passe par $(1/n, 0, 1/n)$ et le segment par $(1/n, 0, 0)$. Ces deux points distincts ont la même image, et tous deux tendent vers $(0, 0, 0)$ quand $n$ grandit : aucun voisinage de ce point n'est envoyé injectivement par $p$.
\end{proof}

\noindent L'énoncé juste demande en plus que $p$ soit un homéomorphisme local, ce que suppose d'ailleurs l'exemple que l'exercice donne entre parenthèses, ou que $Y$ soit localement connexe.

\begin{sloppypar}
\noindent\emph{Vérification en Lean.} L'échec du revêtement est \texttt{not\_isLocallyInjective\_proj}, \texttt{not\_isLocalHomeomorph\_proj} et \texttt{not\_isCoveringMap\_proj} ; les fibres discrètes sont \texttt{discreteTopology\_fiber} ; le relèvement unique est \texttt{existsUnique\_lift}, avec les cas du carré et du chemin \texttt{existsUnique\_lift\_square} et \texttt{existsUnique\_lift\_path} ; le lemme sur les fonctions à valeurs dénombrables est \texttt{level\_eq\_of\_isPreconnected}, dans \href{\leanblobmain/Redaction205.lean}{\texttt{Redaction205.\allowbreak{}lean}}.
\end{sloppypar}

Ne sont pas démontrés ici : la condition (RU) de la rédaction, qui n'est pas énoncée (l'intervalle $[0, 1]$ la vérifie sous toute lecture raisonnable, puisqu'il est contractile et a une base d'ouverts contractiles) ; la caractérisation corrigée elle-même.

\litterature{l'exercice est faux tel qu'il est écrit, et l'énoncé corrigé est un résultat de manuel (Spanier, \emph{Algebraic Topology}, chapitre~2, §~4 \cite{spanier1966} ; l'endroit exact n'a pas été vérifié). Ce n'est pas une piste.}

\section[Rédaction Bourbaki \redac{380} : deux énoncés faux tels que tapés]{Rédaction Bourbaki \redac{380} : un critère de pôle inversé et un spectre dit non compact\protect\verifie{Redaction380}}\label{sec:bourbaki-380}

La rédaction \redac{380}, \q{Analyse fonctionnelle. Chapitre I. Algèbres normées}, des archives Bourbaki (Archives Henri Poincaré, \ahp{935}) est datée d'avril 1962 ; c'est le projet de Koszul qui devint le chapitre~I des \emph{Théories spectrales} (1967). Deux énoncés y sont faux tels qu'ils sont tapés. Sa transcription reste locale pour l'instant ; on n'en cite ici aucun passage.

\emph{Le critère de pôle.} Pour un point isolé $\lambda$ du spectre de $x$, d'idempotent associé $e_\lambda$, le lemme~5 du §~2 dit que $\lambda$ est un pôle d'ordre $p$ de la résolvante quand $(x - \lambda e)^{p-1} e_\lambda \neq 0$ et $(x - \lambda e)^p e_\lambda = 0$. La proposition~6 (p.~40 du fac-similé) le reprend avec les conditions inversées : $(x - \lambda e)^k e_\lambda = 0$ pour $k < p$ et $\neq 0$ pour $k = p$. Aucun élément ne vérifie cette condition, puisque la nullité pour $k = p - 1$ entraîne la nullité pour $k = p$ ; telle que tapée, la proposition dirait que la résolvante n'a jamais de pôle.

\emph{Le spectre d'une algèbre sans unité.} La proposition~16 du §~3 (p.~54) dit que le spectre, l'espace des caractères, d'une algèbre normée sans unité n'est pas compact. Il peut être vide : si tout produit est nul, un caractère $\varphi$ vérifie $\varphi(x)^2 = \varphi(x^2) = 0$, donc $\varphi = 0$, ce qui est exclu. Il peut aussi être un point : l'algèbre $\mathbf C \times N$, où $N$ est de carré nul, n'a pas d'unité, et son seul caractère est la première projection.

\begin{encadre}
\begin{contrexemple}[aux propositions~6 et~16 de la rédaction \redac{380}, telles qu'elles sont tapées]
\begin{enumerate}
\item[a)] Dans tout anneau, pour $p \ge 1$, il n'existe pas d'éléments $a$, $e$ tels que $a^k e = 0$ pour tout $k < p$ et $a^p e \neq 0$. La condition du lemme~5, elle, est réalisée : pour $x = 0$ dans $\mathbf C$, $\lambda = 0$, $e_\lambda = 1$ et $p = 1$ (la résolvante $1/\lambda$ a un pôle simple en~$0$).
\item[b)] Soit $N$ l'espace $\mathbf C$ muni du produit nul. C'est une $\mathbf C$-algèbre de Banach commutative non nulle sans unité, et son espace de caractères est vide, donc compact. L'algèbre de Banach $\mathbf C \times N$ n'a pas d'unité non plus, et son espace de caractères est un point.
\end{enumerate}
\end{contrexemple}
\end{encadre}

\begin{proof}
a) Si $a^{p-1} e = 0$, alors $a^p e = a\,(a^{p-1} e) = 0$. b) Une unité $e$ de $N$ vérifierait $e \cdot 1 = 1$, alors que tout produit est nul ; un caractère est nul, comme on l'a vu. Une unité de $\mathbf C \times N$ aurait pour seconde composante une unité de $N$. Un caractère $\varphi$ de $\mathbf C \times N$ est nul sur $0 \times N$, dont les carrés sont nuls ; $c = \varphi(1, 0)$ vérifie $c^2 = c$, et $c \neq 0$ puisque $\varphi \neq 0$, donc $\varphi$ est la première projection.
\end{proof}

\begin{sloppypar}
\noindent\emph{Vérification en Lean.} a) est \texttt{reversed\_pole\_condition\_never} et \texttt{pole\_condition\_example} ; b) est \texttt{characterSpace\_eq\_empty\_of\_mul\_self\_eq\_zero}, pour toute algèbre où les carrés sont nuls, appliqué à l'algèbre complète \texttt{ZeroMul} (\texttt{ZeroMul.not\_exists\_one}, \texttt{ZeroMul.isCompact\_characterSpace}), puis \texttt{not\_exists\_one\_prod}, \texttt{characterSpace\_prod\_eq} et \texttt{not\_prop16}, qui nie la proposition~16 pour les $\mathbf C$-algèbres normées complètes, dans \href{\leanblobmain/Redaction380.lean}{\texttt{Redaction380.\allowbreak{}lean}}. La résolvante et ses pôles ne sont pas formalisés : a) en porte le contenu algébrique.
\end{sloppypar}

\litterature{le reste du projet est retrouvé dans les \emph{Théories spectrales} (1967) et chez Rickart (1960). Le texte imprimé n'a pas été consulté : on ne sait pas si les deux énoncés y sont corrigés.}

\section[Rédaction Bourbaki \redac{384} : le covolume n'est pas continu]{Rédaction Bourbaki \redac{384} : dans l'espace des sous-groupes fermés, le covolume n'est pas continu\protect\verifie{Redaction384}}\label{sec:bourbaki-384}

La rédaction \redac{384} des archives Bourbaki (Archives Henri Poincaré, \ahp{940}), datée de juin 1962, est le projet de Bruhat pour le §~5 du chapitre~VIII d'\emph{Intégration}, \q{L'espace des sous-groupes fermés} (1963). Elle munit l'ensemble des fermés d'un groupe localement compact $G$ d'une structure uniforme, montre que l'espace $S$ des sous-groupes fermés est compact, et étudie le covolume des sous-groupes $H$ tels que $G/H$ soit compact : il est semi-continu inférieurement, et continu là où $G$ est engendré par un compact. L'exercice~4 montre que la continuité tombe sans cette hypothèse. Sa transcription reste locale pour l'instant ; on n'en cite ici aucun passage.

\emph{L'exemple.} Soit $G = \bigoplus_{n \in \mathbf N} \mathbf Z/2$, discret, que nul compact n'engendre. Les sous-groupes $H_n = \ker p_n$, noyaux des projections sur les coordonnées, sont d'indice~$2$ : pour les mesures de comptage, leur covolume vaut~$2$. Ils tendent vers $G$, de covolume~$1$. Pour un groupe discret, la topologie de $S$ est celle de la convergence simple des fonctions indicatrices, et la convergence dit que tout $x \in G$ est dans $H_n$ pour $n$ assez grand.

\begin{encadre}
\begin{enonce}[Bourbaki, rédaction \redac{384}, exercice~4]
Dans $G = \bigoplus_{n \in \mathbf N} \mathbf Z/2$, chaque $H_n = \ker p_n$ est d'indice~$2$, et tout $x \in G$ appartient à $H_n$ pour $n$ assez grand ; les fonctions indicatrices des $H_n$ tendent simplement vers celle de $G$.
\end{enonce}
\end{encadre}

\begin{proof}
La projection $p_n$ est surjective, puisque $p_n$ envoie la suite qui vaut $c$ en $n$ et $0$ ailleurs sur $c$ ; son noyau est donc d'indice $\operatorname{Card}(\mathbf Z/2) = 2$. Un élément $x$ de $G$ n'a qu'un nombre fini de coordonnées non nulles ; au-delà de la plus grande, $p_n(x) = 0$.
\end{proof}

\begin{sloppypar}
\noindent\emph{Vérification en Lean.} $G$ est le groupe des suites à support fini \texttt{ℕ →₀ ZMod 2} ; l'indice est \texttt{index\_ker}, la convergence \texttt{eventually\_mem\_ker} et \texttt{tendsto\_indicator}, dans \href{\leanblobmain/Redaction384.lean}{\texttt{Redaction384.\allowbreak{}lean}}. La topologie de Chabauty, absente de mathlib, n'est pas formalisée, ni la compacité de $S$.
\end{sloppypar}

\litterature{retrouvé pour le reste de la rédaction : Chabauty (\emph{Bull. SMF} 78, 1950), Fell (\emph{Proc. AMS} 13, 1962), Bourbaki, \emph{Intégration}, chap.~VIII, §~5 (1963), Raghunathan, \emph{Discrete Subgroups of Lie Groups} (1972), chap.~I. L'exercice figure probablement parmi ceux du volume imprimé, qui n'a pas été consulté. L'édition ne donne pas l'exemple pour nouveau.}

\section[Rédaction Bourbaki \redac{551} : une pseudo-covariance dépend de la diagonale]{Rédaction Bourbaki \redac{551} : une pseudo-covariance dépend de ses valeurs sur la diagonale\protect\verifie{Redaction551}}\label{sec:bourbaki-551}

La rédaction \redac{551} des archives Bourbaki (Archives Henri Poincaré, \ahp{1102}), de Cartier, datée de mars 1970, contient une annexe~2, \q{Noyaux de type positif}. Elle y étudie les covariances sur un espace mesuré $(T, m)$, avec $m$ $\sigma$-finie : les fonctions réelles $C$ sur $T \times T$ dont toutes les matrices $(C(t_i, t_j))_{i,j}$ sont symétriques à valeurs propres positives. La transcription reste locale pour l'instant ; on n'en cite ici aucun passage.

\emph{Les pseudo-covariances.} Au n\textsuperscript{o}\,4 (p.~40 du fac-similé), une \emph{pseudo-covariance} est une fonction mesurable $C$ telle que, pour tout $n$ et presque tout $(t_1, \dots, t_n) \in T^n$, la matrice $(C(t_i, t_j))$ soit symétrique à valeurs propres positives. La page ajoute qu'une fonction mesurable presque partout égale, sur $T \times T$, à une pseudo-covariance en est encore une. Le théorème~5 (p.~41) donne quatre conditions dites équivalentes : a) $C$ est une pseudo-covariance ; b) et c) portent sur les intégrales $\iint C(t,t')u(t)v(t')\,dm\,dm$ ; d) $C$ est presque partout une somme $\sum_n f_n(t) f_n(t')$. Les conditions b), c) et d) ne changent pas si l'on modifie $C$ sur une partie négligeable de $T \times T$ ; la condition a), si.

\emph{Le défaut.} Pour $n = 1$, la condition a) dit que $C(t, t) \ge 0$ pour presque tout $t \in T$ : c'est une condition sur la diagonale. Si $m$ n'a pas d'atome, la diagonale est négligeable dans $T \times T$, et l'on peut y changer $C$ sans toucher à b), c), d). La \q{Note de l'éditeur} ironique qui clôt l'annexe (p.~46) relève le même point.

\begin{encadre}
\begin{contrexemple}[au n\textsuperscript{o}\,4 et au théorème~5 de la rédaction \redac{551}, tels qu'ils sont tapés]
Soient $T = \mathbf R$, muni de la mesure de Lebesgue, et $C(t, t') = -1$ si $t = t'$, $0$ sinon. Alors $C$ est presque partout égale sur $T \times T$ à la covariance nulle, et vérifie donc b), c) et d) ; mais pour tout $t$, la matrice $(C(t,t)) = (-1)$ n'est pas positive, et $C$ n'est pas une pseudo-covariance.
\end{contrexemple}
\end{encadre}

\begin{proof}
La diagonale est négligeable : par Fubini, sa mesure est l'intégrale en $t$ de la mesure de $\{t\}$, qui est nulle. Hors de la diagonale, $C$ est nulle. La matrice nulle est positive, et $(-1)$ ne l'est pas.
\end{proof}

\noindent Ce qui reste vrai, d'après la vérification faite pour cette édition et non formalisée : a) entraîne b), et b), c), d) sont équivalentes ; toute fonction qui les vérifie est presque partout égale à une covariance (le corollaire~1 de la rédaction). Il suffit de définir les pseudo-covariances par b).

\begin{sloppypar}
\noindent\emph{Vérification en Lean.} La diagonale est négligeable pour toute mesure sans atome (\texttt{prod\_diagonal\_eq\_zero}) ; la fonction $C$ est presque partout nulle (\texttt{C\_ae\_eq\_zero}), la fonction nulle est une covariance (\texttt{zero\_posSemidef}), et $C$ échoue pour $n = 1$ en tout point (\texttt{C\_not\_pseudoCovariance}), dans \href{\leanblobmain/Redaction551.lean}{\texttt{Redaction551.\allowbreak{}lean}}. Le théorème~5 corrigé n'est pas formalisé.
\end{sloppypar}

\litterature{les noyaux de type positif sont classiques (Aronszajn, 1950 ; Berg, Christensen et Ressel, \emph{Harmonic Analysis on Semigroups}, 1984). Le défaut est relevé dans la rédaction même, par sa \q{Note de l'éditeur} ; l'édition le formalise sans le donner pour nouveau. Le théorème~6 (une majoration presque partout de $\langle f(t), f(t')\rangle$ entraîne celle de $\|f(t)\|^2$) n'a pas été retrouvé tel quel ; il n'est pas formalisé.}


\part{Catégories et algèbre homologique}
\section[Rédaction Bourbaki \redac{321} : carrés cartésiens dans une catégorie abélienne]{Rédaction Bourbaki \redac{321} : carrés cartésiens et cocartésiens dans une catégorie abélienne\protect\verifie{Redaction321}}\label{sec:bourbaki-321}

La rédaction \redac{321} des archives Bourbaki (Archives Henri Poincaré, \ahp{887}), intitulée \q{Pull-backs}, tient en six pages dactylographiées que la notice date de mai--juin 1959 et donne à Serge Lang. Elle fixe un carré commutatif dans une catégorie abélienne et compare deux propriétés universelles, celle du produit fibré (carré cartésien) et celle de la somme amalgamée (carré cocartésien), avec l'exactitude d'une suite construite sur le carré. Sa transcription reste locale pour l'instant ; on n'en cite ici aucun passage.

\emph{Le carré.} On note
\[
\begin{tikzcd}
P \arrow[r, "U"] \arrow[d, "V"'] & B \arrow[d, "v"] \\
A \arrow[r, "u"'] & Q
\end{tikzcd}
\]
le carré, avec $uV = vU$. Dans la catégorie des modules sur un anneau, le produit fibré de $u$ et $v$ est le sous-module $\{(a, b) \in A \oplus B : u(a) = v(b)\}$, c'est-à-dire le noyau de l'application $(a, b) \mapsto u(a) - v(b)$. Dire que le carré est cartésien, c'est dire que $p \mapsto (V(p), U(p))$ identifie $P$ à ce noyau. De même, la somme amalgamée de $V$ et $U$ est le quotient de $A \oplus B$ par les éléments $(V(p), -U(p))$, et le carré est cocartésien quand $(a, b) \mapsto u(a) + v(b)$ identifie $Q$ à ce quotient. La rédaction transporte ces deux descriptions dans une catégorie abélienne quelconque.

\begin{encadre}
\begin{enonce}[Bourbaki, rédaction \redac{321}, proposition~2, p.~4]
Le carré est cartésien si et seulement si la suite $0 \to P \to A \oplus B \to Q$, de flèches $(V, U)$ et $u - v$, est exacte. Il est cocartésien si et seulement si la suite $P \to A \oplus B \to Q \to 0$, de flèches $(V, -U)$ et $u + v$, est exacte.
\end{enonce}
\end{encadre}

\begin{proof}
Une flèche $T \to P$ et un couple $(x, y)$ de flèches $T \to A$, $T \to B$ avec $ux = vy$ sont deux façons de dire la même chose qu'une flèche $T \to A \oplus B$ dont le composé avec $u - v$ est nul. La propriété universelle du produit fibré dit alors exactement que $P$, par $(V, U)$, est le noyau de $u - v$ : c'est l'exactitude de la première suite, flèche de gauche comprise. La seconde moitié est la même chose, flèches renversées : la propriété de la somme amalgamée dit que $Q$, par $u + v$, est le conoyau de $(V, -U)$.
\end{proof}

\noindent La rédaction utilise une seule convention de signe pour les deux suites. Les deux complexes diffèrent par l'automorphisme $(a, b) \mapsto (a, -b)$ de $A \oplus B$, qui ne change pas leur exactitude ; le fichier le vérifie (\texttt{shortComplexIso}).

\emph{Quand un carré cartésien est cocartésien.} Dans les modules, l'intersection de deux sous-modules $M$ et $N$ donne un carré qui est à la fois cartésien et cocartésien, avec $M + N$ au coin opposé. Pour un carré cartésien quelconque, la flèche $A \oplus B \to Q$ n'a aucune raison d'être surjective : si $A = B = 0$ et $Q \neq 0$, le carré est cartésien avec $P = 0$, et il n'est pas cocartésien. La proposition~4 de la page~4 dit que c'est la seule obstruction.

\begin{encadre}
\begin{enonce}[Bourbaki, rédaction \redac{321}, propositions~4 et~5, pp.~4--5]
Un carré cartésien dans lequel $v$ (ou $u$) est un épimorphisme est cocartésien. Dualement, un carré cocartésien dans lequel $U$ (ou $V$) est un monomorphisme est cartésien. Dans un carré cartésien, si $v$ est un épimorphisme, $V$ en est un ; dans un carré cocartésien, si $V$ est un monomorphisme, $v$ en est un.
\end{enonce}
\end{encadre}

\begin{proof}
Si le carré est cartésien, la suite $0 \to P \to A \oplus B \to Q$ est exacte par la proposition~2. Si $v$ est un épimorphisme, la flèche $u - v$ l'est aussi, puisque sa restriction à $B$ est $-v$. La suite $P \to A \oplus B \to Q \to 0$ est donc exacte, au signe près, et le carré est cocartésien par la proposition~2. L'énoncé dual se démontre de la même façon, flèches renversées. Pour la proposition~5, on raisonne d'abord dans les modules : si $v$ est surjective et $a \in A$, on choisit $b$ avec $v(b) = u(a)$, et le couple $(a, b)$ est un élément de $P$ d'image $a$. Dans une catégorie abélienne, le même raisonnement se fait sur des flèches $T \to A$, quitte à remplacer $T$ par un objet $T'$ qui le recouvre par un épimorphisme ; c'est la technique des \q{raffinements} de mathlib, qui tient lieu d'éléments.
\end{proof}

\emph{Une coquille dans le dual.} Tel qu'il est tapé, l'énoncé dual de la proposition~4 demande que la flèche notée $u$ soit un monomorphisme. Si on lit $u$ comme la flèche du bas, $A \to Q$, l'énoncé est faux.

\begin{contrexemple}[rédaction \redac{321}, dual de la proposition~4, lu sur la flèche du bas]
Soit $X$ un objet non nul. Le carré de sommets $P = X$ et $A = B = Q = 0$ est cocartésien, et sa flèche du bas $u$, l'identité de~$0$, est un isomorphisme. Il n'est pas cartésien.
\end{contrexemple}

\begin{proof}
La somme amalgamée de $X \to 0$ et $X \to 0$ est $0$ : le carré est cocartésien. Le produit fibré de $0 \to 0$ et $0 \to 0$ est $0$, qui n'est pas isomorphe à $X$.
\end{proof}

\noindent La flèche voulue est la flèche du haut $U$, celle qui joue pour la somme amalgamée le rôle que joue $v$ pour le produit fibré ; avec elle, l'énoncé est juste, et c'est lui qu'énonce l'encadré.

\noindent\emph{Vérification en Lean.} La proposition~2 est \texttt{isPullback\_iff\_exact} et \texttt{isPushout\_iff\_exact}. Mathlib avait déjà le sens direct de la première moitié ; le fichier ajoute les réciproques. La proposition~4 et son dual sont \texttt{isPushout\_of\_isPullback\_of\_epi\_right}, \texttt{isPullback\_of\_isPushout\_of\_mono\_top} et leurs variantes symétriques, la proposition~5 est \texttt{epi\_fst\_of\_isPullback} (sa moitié sur les monomorphismes vaut dans toute catégorie et est déjà dans mathlib), et le contre-exemple est \texttt{not\_isPullback\_of\_isPushout\_of\_mono\_bottom}, dans \href{\leanblobmain/Redaction321.lean}{\texttt{Redaction321.\allowbreak{}lean}}.

Ne sont pas démontrées ici, parce que mathlib les contient : la proposition~1, existence des produits fibrés et des sommes amalgamées, barrée sur le tapuscrit ; la proposition~3, unicité des flèches vers un produit fibré, barrée aussi ; la proposition~6, recollement de deux carrés cartésiens.

\litterature{les énoncés sont classiques. La proposition~4 est l'étiquette 08N4 du Stacks Project \cite{stacks}, et le relevé des rédactions classe celle-ci parmi les textes retrouvés dans la littérature.}

\section[Rédaction Bourbaki \redac{438} : une identité triangulaire suffit]{Rédaction Bourbaki \redac{438} : une identité triangulaire et un automorphisme font une adjonction\protect\verifie{Redaction438}}\label{sec:bourbaki-438}

La rédaction \redac{438}, \q{Catégories fibrées}, des archives Bourbaki (Archives Henri Poincaré, \ahp{987}) est datée de mars 1965 par sa notice, qui la donne à Claude Chevalley ; elle compte 48 pages. Elle fait partie des textes du livre sur les catégories que Bourbaki n'a pas publié. La démonstration de sa proposition~1, sur les catégories fibrées, a besoin d'une adjonction, et la rédaction isole l'étape dans un lemme~1 (p.~35--36). Sa transcription reste locale pour l'instant ; on n'en cite ici aucun passage.

\emph{Unité, coünité et identités triangulaires.} Une adjonction entre deux foncteurs $F : \mathcal C \to \mathcal D$ et $G : \mathcal D \to \mathcal C$ peut se donner par deux transformations naturelles, une unité $\eta : \mathrm{Id} \to GF$ et une coünité $\varepsilon : FG \to \mathrm{Id}$, soumises à deux identités triangulaires :
\[
F X \xrightarrow{F\eta_X} FGF X \xrightarrow{\varepsilon_{FX}} F X \quad\text{et}\quad G Y \xrightarrow{\eta_{GY}} GFG Y \xrightarrow{G\varepsilon_Y} G Y
\]
sont les identités. L'exemple type est celui du module libre : $F$ envoie un ensemble sur le module libre qu'il engendre, $G$ oublie la structure de module, $\eta$ envoie un élément sur le vecteur de base correspondant, et $\varepsilon$ évalue une combinaison formelle d'éléments d'un module. Les deux identités disent que ces deux opérations se compensent, d'un côté comme de l'autre.

\emph{Le lemme.} Dans la rédaction, les rôles de $F$ et $G$ sont tenus par deux foncteurs notés $R$ et $S$. On suppose qu'une des deux identités triangulaires a lieu, et, de l'autre composé $G\varepsilon_Y \circ \eta_{GY}$, seulement qu'il est un automorphisme de $GY$ pour tout $Y$ (figure~\ref{fig:bourbaki-438}). Le lemme conclut que ce composé est l'identité, de sorte que $\eta$ et $\varepsilon$ sont l'unité et la coünité d'une adjonction.

\begin{figure}[htbp]
\centering
\begin{tikzcd}
FX \arrow[r, "F\eta_X"] \arrow[dr, "\mathrm{id}_{FX}"'] & FGFX \arrow[d, "\varepsilon_{FX}"] \\
& FX
\end{tikzcd}
\qquad\qquad
\begin{tikzcd}
GY \arrow[r, "\eta_{GY}"] \arrow[dr, dashed, "e_Y"'] & GFGY \arrow[d, "G\varepsilon_Y"] \\
& GY
\end{tikzcd}
\caption{Les deux identités triangulaires. Le lemme suppose que le triangle de gauche commute pour tout $X$, et seulement que le composé $e_Y = G\varepsilon_Y \circ \eta_{GY}$ du triangle de droite (en tirets) est un isomorphisme ; il en tire $e_Y \circ e_Y = e_Y$, donc $e_Y = \mathrm{id}_{GY}$, et le triangle de droite commute aussi.}
\label{fig:bourbaki-438}
\end{figure}

\begin{encadre}
\begin{enonce}[Bourbaki, rédaction \redac{438}, lemme~1, p.~35]
Soient $F : \mathcal C \to \mathcal D$, $G : \mathcal D \to \mathcal C$, $\eta : \mathrm{Id} \to GF$ et $\varepsilon : FG \to \mathrm{Id}$. Si $\varepsilon_{FX} \circ F\eta_X = \mathrm{id}_{FX}$ pour tout $X$, et si $e_Y = G\varepsilon_Y \circ \eta_{GY}$ est un isomorphisme pour tout $Y$, alors $e_Y = \mathrm{id}_{GY}$ pour tout $Y$, et $(\eta, \varepsilon)$ définit une adjonction $F \dashv G$. Il en va de même en échangeant les rôles des deux identités triangulaires.
\end{enonce}
\end{encadre}

\emph{L'idée.} Un projecteur inversible est l'identité. Il suffit donc de voir que $e_Y$ est idempotent, $e_Y \circ e_Y = e_Y$ : on compose alors par l'inverse de $e_Y$, et il reste $e_Y = \mathrm{id}$. La rédaction procède autrement, en comparant les deux applications entre ensembles de morphismes qu'induisent $\eta$ et $\varepsilon$. La preuve vérifiée est plus courte.

\begin{proof}
On fixe $Y$ et on écrit $e = G\varepsilon_Y \circ \eta_{GY}$. Dans $e \circ e = G\varepsilon_Y \circ \eta_{GY} \circ G\varepsilon_Y \circ \eta_{GY}$, on fait passer les deux flèches du milieu l'une par-dessus l'autre, par naturalité de $\eta$ relativement à $G\varepsilon_Y$ :
\[
\eta_{GY} \circ G\varepsilon_Y = GFG\varepsilon_Y \circ \eta_{GFGY}.
\]
On obtient $G\varepsilon_Y \circ GFG\varepsilon_Y \circ \eta_{GFGY} \circ \eta_{GY}$. La naturalité de $\varepsilon$ relativement à $\varepsilon_Y$ donne $\varepsilon_Y \circ FG\varepsilon_Y = \varepsilon_Y \circ \varepsilon_{FGY}$, et celle de $\eta$ relativement à $\eta_{GY}$ donne $\eta_{GFGY} \circ \eta_{GY} = GF\eta_{GY} \circ \eta_{GY}$. Il vient
\[
e \circ e = G\varepsilon_Y \circ G\bigl(\varepsilon_{FGY} \circ F\eta_{GY}\bigr) \circ \eta_{GY}.
\]
La parenthèse est l'identité triangulaire supposée, au point $X = GY$ ; elle vaut l'identité, et il reste $e \circ e = e$. Comme $e$ est inversible, $e = \mathrm{id}$, et les deux identités triangulaires ont lieu. Le cas dual se traite de la même façon, avec le composé $\varepsilon_{FX} \circ F\eta_X$.
\end{proof}

\noindent\emph{Vérification en Lean.} L'énoncé encadré est \texttt{rightTriangle\_of\_leftTriangle}, appuyé sur \texttt{rightComposite\_idem} (le composé est idempotent) et \texttt{eq\_id\_of\_isIso\_of\_idem} (un idempotent inversible est l'identité). Le dual est \texttt{leftTriangle\_of\_rightTriangle}. Les adjonctions elles-mêmes sont \texttt{adjunctionOfLeftTriangle} et \texttt{adjunctionOfRightTriangle}, construites par \texttt{Adjunction.\allowbreak{}mkOfUnitCounit} de mathlib, dans \href{\leanblobmain/Redaction438.lean}{\texttt{Redaction438.\allowbreak{}lean}}.

Ne sont pas démontrées ici : la proposition~1 de la rédaction, sur les catégories fibrées et le produit fibré canonique, dont le lemme est une étape.

\litterature{l'énoncé est probablement le lemme~A1.1.1 de l'\emph{Elephant} de Johnstone \cite{johnstone2002} ; la comparaison avec le livre n'a pas été faite.}

\section[Rédaction Bourbaki \redac{448} : objets libres sur un diagramme]{Rédaction Bourbaki \redac{448} : les objets libres sur un diagramme ne se calculent pas toujours point par point\protect\piste\protect\verifie{Redaction448}}\label{sec:bourbaki-448}

La rédaction \redac{448}, \q{Foncteurs adjoints. Limites}, des archives Bourbaki (Archives Henri Poincaré, \ahp{1026}) est datée d'août 1965 par sa notice, qui la donne à Claude Chevalley ; elle compte 107 pages. Une remarque de sa page~20 pose une question et la laisse en exercice au congrès. Sa transcription reste locale pour l'instant ; on n'en cite ici aucun passage.

\emph{Objets libres.} Soient $F : \mathcal Y \to \mathcal X$ un foncteur et $T$ un objet de $\mathcal X$. Un objet libre sur $T$ relativement à $F$ est un objet $E$ de $\mathcal Y$ muni d'une flèche $\varphi : T \to F(E)$ par laquelle toute flèche $T \to F(E')$ se factorise de façon unique, en $F(g) \circ \varphi$ pour un $g : E \to E'$. Pour $F$ le foncteur d'oubli des modules vers les ensembles, c'est le module libre de base $T$. Un objet colibre est la notion duale, avec une flèche $F(E) \to T$. Dans le langage de mathlib, un objet libre est un objet initial de la catégorie des flèches $T \to F(\cdot)$, un objet colibre un objet final de celle des flèches $F(\cdot) \to T$.

\emph{Sur un diagramme.} Soit $D : I \to \mathcal X$ un diagramme et $F_* : \mathcal Y^I \to \mathcal X^I$ la composition par $F$. Un objet libre sur $D$ relativement à $F_*$ est un diagramme $E : I \to \mathcal Y$ avec une transformation naturelle $\varphi : D \to F \circ E$, universelle. Le corollaire de la proposition~4 dit que, si chaque $D(i)$ a un objet libre relativement à $F$, les valeurs $(E(i), \varphi(i))$ sont libres sur les $D(i)$ : l'objet libre sur un diagramme se calcule point par point. La remarque de la page~20 demande si l'hypothèse sur les $D(i)$ est nécessaire.

\emph{L'intuition.} Un objet libre sur un diagramme doit satisfaire une propriété universelle pour toutes les familles compatibles à la fois. Les flèches du diagramme imposent des liens entre les composantes, et ces liens peuvent rendre unique un choix qui, en un point isolé, ne l'est pas. Un objet universel sur le diagramme peut alors exister sans qu'aucun objet universel n'existe en un point. C'est le phénomène connu des limites dans une catégorie de foncteurs, qui ne se calculent pas point par point quand la catégorie d'arrivée en manque.

\begin{encadre}
\begin{contrexemple}[réponse à la remarque de la rédaction \redac{448}, p.~20]
Soient $\mathcal Y$ la catégorie à un objet et une flèche, $\mathcal X$ la catégorie des ensembles, $F$ le foncteur qui choisit un ensemble à un point, $I$ la catégorie $0 \to 1$ et $D$ l'application du point dans $\{\mathrm{vrai}, \mathrm{faux}\}$ de valeur $\mathrm{vrai}$. Il existe un objet colibre sur $D$ relativement à $F_*$, et aucun objet colibre sur $D$ n'a en $1$ une valeur colibre sur $D(1)$ relativement à $F$ : il n'y a d'ailleurs aucun objet colibre sur $D(1)$. En renversant les flèches de la catégorie des ensembles, on obtient le même échec pour les objets libres, cas qu'énonce la remarque. La réponse à la question est donc non.
\end{contrexemple}
\end{encadre}

\begin{proof}
Comme $\mathcal Y$ n'a qu'un objet et qu'une flèche, un objet sur un ensemble $T$ relativement à $F$ n'est qu'un point de $T$, et deux tels objets ont au plus une flèche entre eux, qui existe s'ils sont le même point. Un objet colibre sur $T$ est donc un point par lequel tout point de $T$ se factorise : il en existe un si et seulement si $T$ a un seul point. L'ensemble $D(1)$ en a deux, et n'a pas d'objet colibre. Sur le diagramme $D$, un objet relativement à $F_*$ est une famille compatible de points, un en $0$, un en $1$, le second image du premier. Le point de $D(0)$ est unique, et son image est $\mathrm{vrai}$ : il y a exactement une famille compatible, qui est donc colibre (figure~\ref{fig:bourbaki-448}). Sa valeur en $1$ est un point de $D(1)$, qui ne peut être colibre puisque $D(1)$ n'en a pas.
\end{proof}

\begin{figure}[htbp]
\centering
\begin{tikzpicture}[>=stealth, every node/.style={font=\small}]
\draw[black!55] (0,0) ellipse (0.55 and 0.8);
\node at (0,-1.15) {$D(0)$};
\draw[black!55] (3.4,0) ellipse (0.95 and 1.0);
\node at (3.4,-1.35) {$D(1)$};
\fill[frise] (0,0) circle (2pt) node[left=2pt] {$\ast$};
\fill[frise] (3.2,0.4) circle (2pt) node[right=2pt] {$\mathrm{vrai}$};
\fill (3.2,-0.4) circle (2pt) node[right=2pt] {$\mathrm{faux}$};
\draw[->, frise] (0.12,0.04) to[bend left=12] (3.08,0.38);
\node[above] at (1.6,0.62) {$D(0 \to 1)$};
\end{tikzpicture}
\caption{Le diagramme $D$ de la réponse, sur la catégorie $I = \{0 \to 1\}$. Un objet relativement à $F_*$ est une famille compatible de points, un dans chaque $D(i)$ ; il n'y en a qu'une (en bleu), qui est donc colibre. Sa valeur en $1$, le point $\mathrm{vrai}$, n'est pas colibre sur $D(1)$ : un objet colibre sur $D(1)$ seul serait un point égal à tous les points de $D(1)$, qui en a deux.}
\label{fig:bourbaki-448}
\end{figure}

\noindent\emph{Vérification en Lean.} Le diagramme est \texttt{D}, l'objet colibre sur $D$ est \texttt{cofree}, et le contre-exemple est \texttt{cofree\_not\_pointwise} : il réunit \texttt{isTerminalCofree} et \texttt{isEmpty\_isTerminal\_one}. Le cas des objets libres, dans la catégorie opposée à celle des ensembles et avec la valeur en $0$, est \texttt{free\_not\_pointwise}, dans \href{\leanblobmain/Redaction448.lean}{\texttt{Redaction448.\allowbreak{}lean}}.

Ne sont pas démontrés ici : le corollaire de la proposition~4 lui-même, c'est-à-dire le calcul point par point sous l'hypothèse sur les $D(i)$.

\litterature{le relevé des rédactions tient la réponse pour une piste candidate. Le phénomène général, des limites non calculées point par point dans une catégorie de foncteurs, est connu (nLab, article \emph{functor category}) ; le contre-exemple en est un cas particulier, puisque $\mathcal Y$ y est la catégorie finale. La question elle-même, telle que la pose la rédaction pour un foncteur $F$ quelconque, n'a pas été retrouvée dans la littérature consultée.}

\section[Rédaction Bourbaki \redac{461} : carrés semi-cartésiens]{Rédaction Bourbaki \redac{461} : les carrés semi-cartésiens d'une catégorie abélienne\protect\verifie{Redaction461}}\label{sec:bourbaki-461}

La rédaction \redac{461}, \q{Catégories abéliennes}, des archives Bourbaki (Archives Henri Poincaré, \ahp{1036}) est datée de février 1966 par sa notice, qui la donne à Claude Chevalley ; elle compte 121 pages. Son §~2.6 (p.~65--77 du fac-similé) étudie une classe de carrés commutatifs plus large que celle des carrés cartésiens et cocartésiens, qu'il appelle semi-cartésiens. Sa transcription reste locale pour l'instant ; on n'en cite ici aucun passage.

\emph{Le carré.} On considère, dans une catégorie abélienne, un carré commutatif
\[
\begin{tikzcd}
a \arrow[r, "\alpha_1"] \arrow[d, "\alpha_2"'] & a_1 \arrow[d, "\beta_1"] \\
a_2 \arrow[r, "\beta_2"'] & b
\end{tikzcd}
\]
On note $\mu$ la flèche de la somme amalgamée de $(\alpha_1, \alpha_2)$ vers $b$, et $\nu$ la flèche de $a$ vers le produit fibré de $(\beta_1, \beta_2)$. Le carré est cocartésien quand $\mu$ est un isomorphisme, cartésien quand $\nu$ en est un. La rédaction affaiblit les deux conditions à la fois.

\begin{definition}[rédaction \redac{461}, définition~3, p.~66]
Le carré est \emph{semi-cartésien} si $\mu$ est un monomorphisme et $\nu$ un épimorphisme.
\end{definition}

\emph{L'intuition, dans les modules.} Dire que $\nu$ est surjective, c'est dire que tout couple $(x_1, x_2) \in a_1 \times a_2$ de même image dans $b$ provient d'un élément de $a$. Dire que $\mu$ est injective, c'est dire que les seules relations entre les images de $a_1$ et de $a_2$ dans $b$ sont celles qui viennent de $a$. Les deux conditions disent la même chose. La proposition~13 de la page les donne pour équivalentes entre elles et à l'exactitude de la suite $0 \to a_1 \times_b a_2 \to a_1 \oplus a_2 \to a_1 \sqcup_a a_2 \to 0$, du produit fibré vers la somme amalgamée, ce que reprend la remarque~1 de la page~66. L'édition en démontre une forme équivalente, plus commode à vérifier : la suite $a \to a_1 \oplus a_2 \to b$, de flèches $(\alpha_1, -\alpha_2)$ et $(\beta_1, \beta_2)$, est exacte au milieu. Le carré de flèches $\mathbf Z \to 0$, $\mathbf Z \to 0$, $0 \to \mathbf Z$, $0 \to \mathbf Z$ est semi-cartésien sans être ni cartésien (le produit fibré est nul, $a$ ne l'est pas) ni cocartésien (la somme amalgamée est nulle, $b$ ne l'est pas).

\emph{Chasser les éléments sans éléments.} Une catégorie abélienne n'a pas d'éléments. Mathlib les remplace par des flèches $T \to a$, et une égalité \q{à raffinement près} est une égalité après composition par un épimorphisme $T' \to T$. On raisonne alors comme dans les modules : \q{choisir un antécédent} devient \q{trouver un antécédent après un raffinement}. Toutes les preuves de cette section sont des chasses de ce type, qui passent par un seul lemme, \texttt{exact\_iff} : l'exactitude de $a \to a_1 \oplus a_2 \to b$ s'écrit sur les couples de flèches $x_1 : T \to a_1$, $x_2 : T \to a_2$ tels que $\beta_1 x_1 + \beta_2 x_2 = 0$, qui doivent provenir de $a$ après raffinement.

\begin{encadre}
\begin{enonce}[Bourbaki, rédaction \redac{461}, §~2.6, propositions~13 à~15 et corollaires]
\begin{enumerate}
\item[a)] (Proposition~13, p.~65, sous une forme de l'édition.) Les conditions \q{$\mu$ est un monomorphisme}, \q{$\nu$ est un épimorphisme} et \q{$a \to a_1 \oplus a_2 \to b$ est exacte} sont équivalentes. La page énonce l'équivalence des deux premières, et, pour troisième condition, l'exactitude de $0 \to a_1 \times_b a_2 \to a_1 \oplus a_2 \to a_1 \sqcup_a a_2 \to 0$. Les carrés cartésiens et les carrés cocartésiens sont semi-cartésiens.
\item[b)] (Proposition~14, p.~66--67.) Un carré semi-cartésien où $\alpha_1$ est un monomorphisme est cartésien, et $\beta_2$ est alors un monomorphisme. Dualement, si $\beta_2$ est un épimorphisme, il est cocartésien, et $\alpha_1$ est un épimorphisme.
\item[c)] (Proposition~15, p.~67, et corollaire~2, p.~70.) Un carré cocartésien suivi d'un carré cartésien est semi-cartésien. Réciproquement, tout carré semi-cartésien s'obtient ainsi.
\item[d)] (Proposition~15, p.~67, et corollaire~1, p.~69.) Si un carré semi-cartésien se décompose en un carré du haut et un carré du bas, le carré du haut est semi-cartésien dès que la flèche verticale gauche du carré du bas est un monomorphisme, et le carré du bas l'est dès que la flèche verticale droite du carré du haut est un épimorphisme. Avec les deux hypothèses, le carré du haut est cocartésien et celui du bas cartésien.
\item[e)] (Corollaire~4 de la proposition~16, p.~73--74.) Le carré est semi-cartésien si et seulement si les suites $\ker \alpha_1 \to a_2 \to b$ et $a \to a_1 \to \operatorname{coker} \beta_2$ sont exactes.
\end{enumerate}
\end{enonce}
\end{encadre}

\begin{remarque}
Pour d), la rédaction suppose le carré du bas monomorphe et celui du haut épimorphe dans la catégorie des flèches, c'est-à-dire deux monomorphismes et deux épimorphismes. Une seule des deux flèches suffit chaque fois. L'énoncé démontré est donc un peu plus général que celui de la page.
\end{remarque}

\begin{proof}
a) La somme amalgamée de $(\alpha_1, \alpha_2)$ est le conoyau de $a \to a_1 \oplus a_2$ ; la flèche $a_1 \oplus a_2 \to b$ se factorise par elle, en $\mu$. Donc $\mu$ est injective si et seulement si le noyau de $a_1 \oplus a_2 \to b$ n'est pas plus gros que l'image de $a$ : c'est l'exactitude. De même, le produit fibré de $(\beta_1, \beta_2)$ est, au signe près, le noyau de $a_1 \oplus a_2 \to b$, et $\nu$ est surjective si et seulement si $a$ le recouvre : c'est encore l'exactitude. Un carré cocartésien a une suite exacte, par la proposition~2 de la rédaction \redac{321} (section~\ref{sec:bourbaki-321}) ; pour un carré cartésien, tout couple $(x_1, x_2)$ de même image dans $b$ (au signe près) se relève par la propriété universelle, sans raffinement.

b) Si $\alpha_1$ est un monomorphisme, $\nu$ en est un, car $\alpha_1$ est le composé de $\nu$ et de la première projection. Comme $\nu$ est aussi un épimorphisme, c'est un isomorphisme, et le carré est cartésien. Si $\zeta : T \to a_2$ vérifie $\beta_2 \zeta = 0$, le couple $(0, \zeta)$ définit une flèche vers le produit fibré, qui provient d'une flèche $y : T \to a$ avec $\alpha_1 y = 0$ ; comme $\alpha_1$ est un monomorphisme, $y = 0$, et $\zeta = \alpha_2 y = 0$. Donc $\beta_2$ est un monomorphisme. La seconde moitié est duale.

c) Dans un carré cocartésien suivi d'un carré cartésien, un couple $(x_1, x_2)$ de même image dans $b$ se relève d'abord dans le produit fibré du carré du bas, puis, après raffinement, dans $a$ par l'exactitude du carré du haut. Pour la réciproque, la rédaction factorise le carré dans la catégorie des flèches en un épimorphisme suivi d'un monomorphisme, puis applique le corollaire~1. La preuve vérifiée prend une décomposition plus directe : la somme amalgamée de $(\alpha_1, \alpha_2)$ au-dessus, puis le carré formé de $a_2 \to a_1 \sqcup_a a_2$, de l'identité de $a_2$, de $\mu$ et de $\beta_2$. Ce second carré est cartésien parce que $\mu$ est un monomorphisme : deux flèches qui ont même image par $\mu$ sont égales.

d) Si la flèche verticale gauche $v_0$ du carré du bas est un monomorphisme, un couple pour le carré du haut, poussé par $v_0$, devient un couple pour le grand carré ; il se relève dans $a$, et l'on revient en simplifiant par $v_0$. Si la flèche verticale droite $u_1$ du carré du haut est un épimorphisme, on relève d'abord, après raffinement, la composante qui vit au sommet du carré du bas en une flèche vers $a_1$, puis on applique la propriété du grand carré. Avec les deux hypothèses, b) appliqué aux deux petits carrés donne le dernier point.

e) Dans les modules : un couple $(x_1, x_2)$ de même image dans $b$ donne, modulo $\beta_2$, un élément de $a_1$ nul dans $\operatorname{coker} \beta_2$, qui provient de $a$ par la seconde suite ; la différence restante est un élément de $a_2$ d'image nulle dans $b$, qui provient du noyau de $\alpha_1$ par la première. La réciproque se lit dans l'autre sens. La preuve vérifiée fait la même chasse, à raffinement près.
\end{proof}

\begin{sloppypar}
\noindent\emph{Vérification en Lean.} La définition est \texttt{IsSemiCartesian}. La proposition~13 est \texttt{mono\_fromPushout\_iff\_exact}, \texttt{epi\_toPullback\_iff\_exact} et \texttt{isSemiCartesian\_iff\_exact} ; les cas cartésien et cocartésien sont \texttt{IsPullback.\allowbreak{}isSemiCartesian} et \texttt{IsPushout.\allowbreak{}isSemiCartesian}, et la symétrie du carré est \texttt{IsSemiCartesian.\allowbreak{}flip}. La proposition~14 est \texttt{IsSemiCartesian.\allowbreak{}isPullback\_of\_mono} et \texttt{IsSemiCartesian.\allowbreak{}isPushout\_of\_epi}. La proposition~15 est \texttt{isSemiCartesian\_of\_isPushout\_of\_isPullback}, \texttt{IsSemiCartesian.\allowbreak{}left\_of\_comp} et \texttt{IsSemiCartesian.\allowbreak{}right\_of\_comp}. Les corollaires~1, 2 et~4 sont \texttt{IsSemiCartesian.\allowbreak{}isPushout\_isPullback\_of\_comp}, \texttt{isSemiCartesian\_iff\_exists\_factorisation} et \texttt{isSemiCartesian\_iff\_exact\_kernel\_cokernel}, dans \href{\leanblobmain/Redaction461.lean}{\texttt{Redaction461.\allowbreak{}lean}}. Le petit exemple de $\mathbf Z$ n'est pas formalisé.
\end{sloppypar}

Ne sont pas démontrés ici : la forme de la page de la proposition~13 et de la remarque~1, l'exactitude de la suite du produit fibré vers la somme amalgamée ; le corollaire de la proposition~14 (p.~67), qui dit que les monomorphismes et les épimorphismes sont universels ; la démonstration propre à la rédaction de la nécessité dans le corollaire~2, par la factorisation épimorphisme-monomorphisme dans la catégorie des flèches ; la proposition~16 et ses corollaires~1 (p.~72), 2 et~3, ces deux derniers étant le lemme des quatre et le lemme des cinq, déjà dans mathlib ; la proposition~17 (p.~74--76), proche du lemme des neuf, et la proposition~17~bis (p.~76--77).

\litterature{retrouvé, pour l'essentiel. J.-C.~Raoult \cite{raoult2019}, qui dit tenir les carrés semi-cartésiens d'un cours de DEA de Chevalley à l'Institut Henri Poincaré (1965--1966), en donne au \S\,2 la définition et les propriétés de base (pagination de la version arXiv) : la proposition~1 (p.~3--4) établit l'équivalence de $\nu$ épimorphisme, de $\mu$ monomorphisme, de l'exactitude de $0 \to a_1 \times_b a_2 \to a_1 \oplus a_2 \to a_1 \sqcup_a a_2 \to 0$ et de celle de $a \to a_1 \oplus a_2 \to b$ ; la forme retenue par l'édition en a) y figure donc. Les carrés cartésiens et cocartésiens y sont semi-cartésiens (p.~4) ; la proposition~2 (p.~4) est la proposition~14 de la page ; la proposition~3 (p.~4--5) contient la première assertion de la proposition~15 ; la proposition~4 (p.~5) en donne la deuxième et la troisième assertion, sous l'hypothèse que le carré que l'on simplifie est un épimorphisme, resp.\ un monomorphisme, de la catégorie des flèches, plus forte que celle de d) ; son corollaire (p.~6) est le corollaire~2. Le corollaire~4 de la proposition~16 (caractérisation par $\ker \alpha_1 \to a_2 \to b$ et $a \to a_1 \to \operatorname{coker} \beta_2$) n'y a pas été vu. Les deux traités de la même époque ne l'ont pas. P.~Freyd \cite{freyd1964} caractérise, à la proposition~2.53 (p.~52), les carrés cartésiens, cocartésiens et à la fois l'un et l'autre (ses « diagrammes de Doolittle ») par l'exactitude de $a \to a_1 \oplus a_2 \to b$ aux extrémités ; l'exactitude au milieu seule, qui est la condition semi-cartésienne, n'y apparaît ni dans le chapitre~2 ni dans ses exercices. B.~Mitchell \cite{mitchell1965} n'a que le lemme des neuf avec produit fibré et somme amalgamée (I, proposition~16.5, p.~23) ; son texte entier, parcouru par mots-clés (\q{square}, \q{pullback}, \q{pushout}), n'a pas la notion. La recherche dans les articles (Puppe, Hilton) n'est pas faite.}

\section[Rédaction Bourbaki \redac{469} : le foncteur constant crée les limites]{Rédaction Bourbaki \redac{469} : le foncteur constant crée les limites\protect\piste\protect\verifie{Redaction469}}\label{sec:bourbaki-469}

La rédaction \redac{469}, \q{Catégories abéliennes}, des archives Bourbaki (Archives Henri Poincaré, \ahp{1042}) est datée de mai--juin 1966 par sa notice, qui la donne à Pierre Deligne ; elle compte 54 pages. Son théorème~4 (p.~23) montre qu'une catégorie de foncteurs $\mathcal C \to \mathcal A$ hérite des propriétés de $\mathcal A$ : limites, colimites, objet nul, structure abélienne. Les errata de la page~52 reprennent la réciproque par un lemme sur le foncteur des diagrammes constants. Sa transcription reste locale pour l'instant ; on n'en cite ici aucun passage.

\emph{La réciproque et son piège.} On voudrait conclure dans l'autre sens : si $\mathcal C \to \mathcal A$ a des limites, $\mathcal A$ en a aussi, pourvu que $\mathcal C$ ne soit pas vide. L'idée naturelle est de plonger $\mathcal A$ dans la catégorie de foncteurs par les foncteurs constants, $\mathrm{const} : \mathcal A \to \mathcal A^{\mathcal C}$, de prendre la limite $L$ du diagramme constant $\mathrm{const} \circ F$ pour un diagramme $F : J \to \mathcal A$, puis de l'évaluer en un objet $d$ de $\mathcal C$. Mais l'évaluation ne préserve les limites que si celles de la catégorie de foncteurs se calculent point par point, ce qui suppose que $\mathcal A$ en ait. L'argument tourne en rond, et il faut une autre preuve que $L(d)$ est une limite de $F$.

\begin{encadre}
\begin{enonce}[Bourbaki, rédaction \redac{469}, lemme des errata, p.~52]
Soit $\mathcal C$ une catégorie non vide. Le foncteur $\mathrm{const} : \mathcal A \to \mathcal A^{\mathcal C}$ crée les limites et les colimites : si $L$ est une limite de $\mathrm{const} \circ F$ dans $\mathcal A^{\mathcal C}$, alors, pour tout objet $d$ de $\mathcal C$, $L(d)$ est une limite de $F$ dans $\mathcal A$, et $L$ envoie toute flèche de $\mathcal C$ sur un isomorphisme. En conséquence, $\mathcal A$ a les limites (finies, ou d'une forme donnée) et les colimites que $\mathcal A^{\mathcal C}$ possède, un objet nul si $\mathcal A^{\mathcal C}$ en a un, et, si $\mathcal A$ est préadditive et $\mathcal A^{\mathcal C}$ abélienne, $\mathcal A$ est abélienne.
\end{enonce}
\end{encadre}

\emph{L'idée.} La limite $L$ est un foncteur, et chacune de ses valeurs $L(d)$ porte un cône sur $F$. Ce cône, rendu constant, donne une flèche $u_d$ du foncteur constant de valeur $L(d)$ vers $L$. On recolle ensuite les composantes diagonales $(u_d)_d$ en un endomorphisme de $L$, compatible aux projections ; par l'unicité dans la propriété universelle de $L$, c'est l'identité. Cette identité sur la diagonale suffit à donner, en chaque $d$, l'unicité qui manquait.

\begin{figure}[htbp]
\centering
\begin{tikzcd}[column sep=large]
\mathrm{const}\, L(d) \arrow[r, "u_d"] \arrow[rr, bend right=18, "\mathrm{const}\,(\pi_j)_d"'] & L \arrow[r, "\pi_j"] & \mathrm{const}\, F(j)
\end{tikzcd}
\caption{La flèche $u_d$ de la preuve, dans la catégorie de foncteurs $\mathcal A^{\mathcal C}$. Les composantes $(\pi_j)_d : L(d) \to F(j)$ forment un cône sur $F$ ; rendu constant, il se factorise par la limite $L$ de $\mathrm{const} \circ F$. Les composantes diagonales $(u_d)_d : L(d) \to L(d)$ forment un endomorphisme de $L$ compatible aux projections $\pi_j$, qui est donc l'identité.}
\label{fig:bourbaki-469}
\end{figure}

\begin{proof}
Notons $\pi_j : L \to \mathrm{const}\, F(j)$ les projections de la limite. En un objet $d$, les composantes $(\pi_j)_d : L(d) \to F(j)$ forment un cône sur $F$. Le cône constant correspondant se factorise par $L$ : il existe $u_d : \mathrm{const}\, L(d) \to L$ avec $\pi_j \circ u_d = \mathrm{const}\,(\pi_j)_d$ (figure~\ref{fig:bourbaki-469}). Pour une flèche $f : d \to d'$, les deux flèches $u_{d'} \circ \mathrm{const}\, L(f)$ et $u_d$ ont les mêmes composés avec les projections, donc sont égales. Cette relation dit que les composantes $(u_d)_d : L(d) \to L(d)$ forment une transformation naturelle $w : L \to L$. Elle vérifie $\pi_j \circ w = \pi_j$ pour tout $j$, donc $w = \mathrm{id}_L$ : chaque $(u_d)_d$ est l'identité de $L(d)$.

Montrons que $L(d)$ est une limite de $F$. Un cône $s$ de sommet $X$ sur $F$ donne, rendu constant, une flèche $\mathrm{const}\, X \to L$, dont la composante en $d$ est une flèche $X \to L(d)$ compatible au cône. Pour l'unicité, soient $m, m' : X \to L(d)$ deux flèches compatibles. Les flèches $u_d \circ \mathrm{const}\, m$ et $u_d \circ \mathrm{const}\, m'$ de $\mathrm{const}\, X$ vers $L$ ont les mêmes composés avec les projections, donc sont égales ; leurs composantes en $d$ sont $(u_d)_d \circ m = m$ et $m'$. Ainsi $m = m'$. Comme $L(d)$ et $L(d')$ sont deux limites du même diagramme, la flèche $L(f)$, compatible aux cônes, est l'isomorphisme canonique entre elles.

Les colimites se traitent de façon duale. Pour l'objet nul, un objet nul de $\mathcal A^{\mathcal C}$ est la limite et la colimite du diagramme vide, et sa valeur en un objet $d$ est donc à la fois finale et initiale dans $\mathcal A$. Pour la structure abélienne, $\mathcal A$ a les limites et colimites finies par ce qui précède ; la flèche canonique de la coimage vers l'image d'une flèche $f$ de $\mathcal A$ devient, par $\mathrm{const}$, celle de $\mathrm{const}\, f$, qui est un isomorphisme puisque $\mathcal A^{\mathcal C}$ est abélienne. Enfin $\mathrm{const}$ reflète les isomorphismes : l'inverse de $\mathrm{const}\, f$ a pour composante, en n'importe quel objet de $\mathcal C$, un inverse de $f$.
\end{proof}

\noindent\emph{Vérification en Lean.} Le lemme est l'instance \texttt{createsLimitsOfSize\_const}, avec son dual \texttt{createsColimitsOfSize\_const} ; ses étapes sont \texttt{u\_app\_self} (les composantes diagonales sont des identités), \texttt{isLimitEvalCone} ($L(d)$ est une limite) et \texttt{isIso\_pt\_map} ($L$ envoie toute flèche sur un isomorphisme). Les conséquences sont \texttt{hasLimitsOfShape\_of\_functorCategory}, \texttt{hasFiniteLimits\_of\_functorCategory} et leurs duales, \texttt{hasZeroObject\_of\_functorCategory} et \texttt{abelianOfFunctorCategory}, dans \href{\leanblobmain/Redaction469.lean}{\texttt{Redaction469.\allowbreak{}lean}}. Mathlib avait déjà le fait que $\mathrm{const}$ préserve les limites qui existent dans $\mathcal A$ (\texttt{preservesLimits\_const}), mais pas la création.

Ne sont pas démontrés ici : la structure préadditive de $\mathcal A$, que le corollaire abélien suppose donnée au lieu de la tirer de $\mathcal A^{\mathcal C}$ (addition des flèches par les biproduits et l'objet nul, puis opposés) ; le résultat est donc partiel sur ce point. Le théorème~4 lui-même, dans le sens direct, est \texttt{functorCategoryAbelian} de mathlib.

\litterature{le lemme des errata est tenu pour une piste candidate. La recherche consignée se limite à mathlib, qui n'a que la préservation des limites par $\mathrm{const}$.}

\section[Rédaction Bourbaki \redac{476} : deux exercices sur les graphes]{Rédaction Bourbaki \redac{476} : deux exercices sur les graphes qui demandent une hypothèse\protect\verifie{Redaction476}}\label{sec:bourbaki-476}

La rédaction \redac{476}, \q{Catégories. Chapitre~I. Diagrammes et graphes}, des archives Bourbaki (Archives Henri Poincaré, \ahp{1048}) est datée d'octobre 1966 par sa notice, qui la donne à Pierre Cartier ; elle compte 69 pages. Son premier paragraphe se termine (p.~6) par des exercices sur les graphes et leurs orientations. Deux d'entre eux, tels qu'ils sont tapés, sont faux. Sa transcription reste locale pour l'instant ; on n'en cite ici aucun passage.

\emph{Orientations.} Un graphe a des sommets et des arêtes, une arête joignant deux sommets distincts. L'orienter, c'est choisir pour chaque arête un sens : l'arête $\{v, w\}$ devient une flèche $v \to w$ ou $w \to v$. Un arbre est un graphe connexe sans cycle : entre deux sommets, il y a un seul chemin sans retour en arrière. Dans un arbre fini à $n$ sommets, il y a $n - 1$ arêtes ; on ne peut donc pas l'orienter de façon que chaque sommet soit la source d'exactement une flèche, car il faudrait $n$ flèches. L'exercice~2 affirme que c'est possible pour tout arbre infini.

\emph{Le contre-exemple.} L'étoile infinie, un centre relié à une infinité dénombrable de feuilles, est un arbre infini. Une feuille n'a qu'une arête, qui doit partir d'elle ; toutes les arêtes pointent donc vers le centre, et le centre n'est la source d'aucune flèche (figure~\ref{fig:bourbaki-476}).

\begin{figure}[htbp]
\centering
\begin{tikzpicture}[>=stealth, every node/.style={font=\small}]
\begin{scope}
\fill (0,0) circle (2pt);
\foreach \k in {0,...,8} {
  \fill ({90+\k*36}:1.35) circle (1.7pt);
  \draw[->, frise, shorten >=3pt, shorten <=3pt] ({90+\k*36}:1.35) -- (0,0);
}
\foreach \k in {9} { \draw[black!45, dotted, thick] ({90+\k*36-8}:1.25) arc ({90+\k*36-8}:{90+\k*36+8}:1.25); }
\node at (0,-1.95) {l'étoile infinie};
\end{scope}
\begin{scope}[xshift=4.2cm, yshift=-0.5cm]
\node (a) at (0,0) {$a$};
\node (c) at (1.2,1.3) {$c$};
\node (b) at (2.4,0) {$b$};
\draw[->, frise] (a) -- (b);
\draw[->, frise] (a) -- (c);
\draw[->, frise] (c) -- (b);
\node at (1.2,-1.45) {le carquois de l'exercice~3};
\end{scope}
\end{tikzpicture}
\caption{À gauche, l'étoile infinie, dont on dessine neuf feuilles : l'arête d'une feuille doit partir d'elle, toutes les flèches vont au centre, et le centre n'est la source d'aucune flèche. À droite, le carquois $a \to b$, $a \to c$, $c \to b$ : il n'a pas de chemin fermé, et ses deux chemins de $a$ à $b$, de longueurs $1$ et $2$, interdisent tout morphisme vers $\mathbf Z$ ou $\mathbf Z/d$ pour $d \ne 1$.}
\label{fig:bourbaki-476}
\end{figure}

\emph{Ce qui est vrai.} L'énoncé tient pour les arbres qui ont un \emph{rayon}, c'est-à-dire un chemin infini $f(0), f(1), f(2), \dots$ qui ne repasse jamais par le même sommet. Il suffit d'orienter chaque arête vers le côté où file le rayon. Si l'on retire l'arête $\{v, w\}$, l'arbre se coupe en deux morceaux, et le rayon finit dans l'un d'eux ; on oriente $v \to w$ si c'est celui de $w$. Dans un arbre infini où chaque sommet n'a qu'un nombre fini de voisins, un rayon existe toujours, par le lemme de König. C'est sans doute le cas que l'exercice avait en vue.

\begin{encadre}
\begin{contrexemple}[rédaction \redac{476}, exercice~2, p.~6]
L'étoile infinie est un arbre infini qui n'a aucune orientation où chaque sommet est la source d'exactement une flèche.
\end{contrexemple}
\begin{enonce}[exercice~2 corrigé, un sens, énoncé de l'édition]
Un arbre qui a un rayon a une orientation où chaque sommet est la source d'exactement une flèche.
\end{enonce}
\end{encadre}

\begin{proof}
Le contre-exemple a été vu. Pour l'énoncé corrigé, on oriente comme on l'a dit, et il faut voir deux choses. D'abord, chaque arête reçoit un seul sens : quand on la retire, l'arbre a exactement deux morceaux, et la queue du rayon, à partir d'un rang où il n'emprunte plus cette arête, reste dans l'un des deux. Ensuite, chaque sommet $v$ est la source d'exactement une flèche. Le rayon ne repasse pas par $v$ ; à partir d'un certain rang, il évite $v$, et sa queue reste dans un seul morceau de l'arbre privé de $v$. Ce morceau est relié à $v$ par une seule arête, faute de quoi il y aurait un cycle, et c'est la seule arête qui part de $v$.
\end{proof}

\emph{L'exercice~3.} Un carquois est un graphe orienté, où l'on admet plusieurs flèches et des boucles. Sa période est le plus grand commun diviseur des longueurs de ses chemins fermés, et vaut $0$ s'il n'en a pas. L'exercice relie la période $d$ aux morphismes du carquois vers le cycle $\mathbf Z/d$ (vers $\mathbf Z$ si $d = 0$), c'est-à-dire aux fonctions $h$ sur les sommets telles que $h(t) = h(s) + 1$ pour toute flèche $s \to t$. Pour un carquois fortement connexe, où l'on va de tout sommet à tout autre, une telle fonction existe pour $\mathbf Z/d$. Sans cette hypothèse, que l'exercice ne pose pas, l'énoncé tombe.

\begin{encadre}
\begin{contrexemple}[rédaction \redac{476}, exercice~3, sans connexité forte]
Le carquois à trois sommets $a \to b$, $a \to c$, $c \to b$ n'a aucun chemin fermé, et sa période est $0$. Il n'a de morphisme vers $\mathbf Z/d$ que pour $d = 1$ : ni vers $\mathbf Z$, ni vers $\mathbf Z/2$. Le triangle sous-jacent n'a pas de coloriage en deux couleurs.
\end{contrexemple}
\end{encadre}

\begin{proof}
Une hauteur strictement croissante le long des flèches ($0$ en $a$, $1$ en $c$, $2$ en $b$) interdit tout chemin fermé. Les deux chemins de $a$ à $b$ ont pour longueurs $1$ et $2$ ; une fonction $h$ comme ci-dessus donnerait $h(b) = h(a) + 1 = h(a) + 2$, soit $1 = 0$ dans $\mathbf Z/d$, ce qui force $d = 1$. Enfin un triangle a un nombre impair de côtés et ne se colorie pas en deux couleurs.
\end{proof}

\noindent\emph{Vérification en Lean.} Une orientation est la structure \texttt{Orientation}, la condition sur les sources est \texttt{Orientation.\allowbreak{}OutDegreeOne}. L'étoile est un arbre infini (\texttt{star\_isTree}), le contre-exemple est \texttt{not\_exists\_outDegreeOne\_star}, et l'énoncé corrigé est \texttt{exists\_outDegreeOne\_of\_ray}. Pour l'exercice~3, l'absence de chemin fermé est \texttt{ex3\_noClosedPath}, les morphismes vers les cycles sont traités par \texttt{ex3\_no\_hom\_zmod} et \texttt{ex3\_no\_hom\_zmod2}, et le coloriage par \texttt{ex3\_not\_bipartite}, dans \href{\leanblobmain/Redaction476.lean}{\texttt{Redaction476.\allowbreak{}lean}}.

Ne sont pas démontrés ici : la réciproque de l'exercice~2 corrigé (une orientation de ce type donne un rayon, en suivant les flèches depuis un sommet quelconque) ; le lemme de König ; l'énoncé général de l'exercice~3 pour les carquois fortement connexes.

\litterature{les deux exercices ne sont pas des pistes : ce sont des erreurs d'énoncé, et aucune recherche bibliographique n'a été consignée à leur sujet.}


\part{Théorie des ensembles}
\section[Rédaction Bourbaki n\textsuperscript{o}\,453 : les univers forment une chaîne]{Rédaction Bourbaki n\textsuperscript{o}\,453 : sous l'axiome de Grothendieck, les univers forment une chaîne\protect\verifie{Redaction453}}\label{sec:bourbaki-453}

La rédaction n\textsuperscript{o}\,453 des archives Bourbaki (Archives Henri Poincaré, \href{https://archives-bourbaki.ahp-numerique.fr/items/show/1031}{item 1031}), que sa notice date de septembre 1965 et donne à Pierre Samuel, tient en une page dactylographiée et non signée sur les univers et le symbole~$\tau$. Elle rapporte un résultat d'un étudiant de troisième cycle de Samuel : l'axiome de Grothendieck sur les univers entraîne que tout ensemble est artinien. Le congrès de mars 1965, ajoute-t-elle, l'avait soupçonné. Les univers sont le fondement ensembliste de la théorie des topos : SGA~4 les introduit dans l'appendice de son exposé~I \cite{sga4-1}. Sa transcription reste locale pour l'instant ; on n'en cite ici aucun passage.

\emph{Les univers.} Un univers est un ensemble $U$ stable par les opérations usuelles : il contient les éléments de ses éléments, les paires et l'ensemble des parties de ses éléments, et les réunions de familles indexées par l'un de ses éléments. Toute construction faite à partir d'éléments de $U$ reste dans $U$. On peut alors parler de catégories \q{petites} relativement à $U$ et appliquer à la catégorie de tous les petits ensembles, qui n'est pas petite, les constructions qu'on ferait sur un ensemble. SGA~4 pose que tout ensemble appartient à un univers.

\emph{Le symbole $\tau$.} La \emph{Théorie des ensembles} de Bourbaki ne prend pas les quantificateurs comme primitifs. Pour une relation $P(x)$, elle forme un objet $\tau_x P(x)$, qui vérifie $P$ dès qu'un objet quelconque la vérifie : c'est un témoin choisi une fois pour toutes, et l'on définit $(\exists x)\,P(x)$ comme $P(\tau_x P(x))$. L'axiome de Grothendieck, noté (AG) dans la rédaction, demande que ce choix respecte les univers.

\begin{definition}
Une famille $\mathcal U$ d'ensembles vérifie (AG) si, pour toute relation $P$ et tout $U \in \mathcal U$ qui contient un objet vérifiant $P$, on a $\tau_x P(x) \in U$.
\end{definition}

\noindent Un ensemble est \emph{artinien} quand l'appartenance n'y a pas de chaîne infinie descendante $x_0 \ni x_1 \ni x_2 \ni \cdots$. Dire que tout ensemble est artinien, c'est l'axiome de fondation, que la théorie de Bourbaki ne pose pas. Le résultat rapporté est donc surprenant : une condition sur le choix des témoins entraîne la fondation.

\emph{Le pas a).} La démonstration de la rédaction a deux pas, et le premier est le cœur de l'affaire. Il dit que deux univers sont toujours comparables pour l'inclusion. Si $U \not\subset U'$ et $U' \not\subset U$, on applique (AG) à la relation \q{$x \notin U \cap U'$}. Elle a un témoin dans $U$, un élément de $U$ hors de $U'$, et un témoin dans $U'$. Le témoin choisi $t = \tau_x(x \notin U \cap U')$ est donc à la fois dans $U$ et dans $U'$. La rédaction s'arrête à cette double appartenance. La contradiction vient de ce que $t$ vérifie la relation, puisqu'elle a un témoin : $t \notin U \cap U'$.

\begin{encadre}
\begin{enonce}[Bourbaki, rédaction n\textsuperscript{o}\,453, pas a)]
Si une famille $\mathcal U$ d'ensembles vérifie (AG), deux éléments de $\mathcal U$ sont toujours comparables pour l'inclusion : $\mathcal U$ est une chaîne.
\end{enonce}
\end{encadre}

\begin{remarque}
Le pas a) n'utilise aucune propriété des univers : ni leurs stabilités, ni l'existence d'un univers contenant un ensemble donné. Il vaut pour toute famille d'ensembles qui vérifie (AG), et ne demande du symbole $\tau$ que la propriété du témoin. Dans la théorie usuelle, où l'on pose l'axiome de fondation, les univers non dénombrables sont les ensembles $V_\kappa$ pour $\kappa$ fortement inaccessible, et sont donc comparables. (AG) donne ici cette comparaison sans fondation.
\end{remarque}

\begin{demopage}[pas b)]
Soit $y$ un ensemble non artinien, et $U$ un univers qui le contient. Soit $c$ la borne supérieure des cardinaux des éléments de $U$. Un ordinal de von Neumann $z$ de cardinal $2^c$ est artinien, et, d'après la rédaction n\textsuperscript{o}\,424, les éléments artiniens d'un univers qui contient $z$ forment un univers $U'$ qui contient $z$. Alors $U \not\subset U'$, à cause de $y$, et $U' \not\subset U$, car $z$ est de cardinal $2^c > c$ et ne peut appartenir à $U$. Cela contredit le pas a).
\end{demopage}

\noindent\emph{Vérification en Lean.} Lean n'a pas d'appartenance entre ensembles au sens de Bourbaki. Les \q{ensembles} y sont les éléments d'un type $\alpha$ non vide, une famille d'univers est une famille $\mathcal U$ de parties de $\alpha$, et $\tau$ est le choix de Hilbert de Lean, \texttt{Classical.\allowbreak{}epsilon}, qui vérifie $P(\varepsilon P)$ dès qu'un objet vérifie $P$ (\texttt{Classical.\allowbreak{}epsilon\_spec}). Ce modèle suffit, puisque le pas a) ne fait intervenir que l'appartenance aux deux univers et la propriété du témoin. L'axiome (AG) est \texttt{AxiomAG}, et l'énoncé encadré est \texttt{isChain\_of\_axiomAG}, dans \href{\leanblobmain/Redaction453.lean}{\texttt{Redaction453.\allowbreak{}lean}}.

Ne sont pas démontrés ici : le pas b), qui demande les ensembles artiniens, les ordinaux de von Neumann, l'arithmétique des cardinaux d'un univers et le résultat de la rédaction n\textsuperscript{o}\,424 ; la conclusion de la rédaction, que tout ensemble est artinien.


\part{Algèbre}
\section[Rédaction Bourbaki \redac{197} : la courbe $y^2 = P(x)$ en caractéristique~$2$]{Rédaction Bourbaki \redac{197} : une courbe $y^2 = P(x)$ qui n'est pas normale en caractéristique~$2$\protect\verifie{Redaction197}}\label{sec:bourbaki-197}

La rédaction \redac{197}, \q{Rapport de Géométrie algébrique. Chapitre~I. Théorie globale élémentaire}, des archives Bourbaki (Archives Henri Poincaré, \ahp{612}) est datée de mai 1954 par sa notice, qui la donne à Pierre Samuel. Elle pose les premières notions de la géométrie algébrique globale et les illustre d'exemples de courbes planes. Sa transcription reste locale pour l'instant ; on n'en cite ici aucun passage.

\emph{Courbes normales.} Soient $k$ un corps et $P \in k[x]$ un polynôme. La courbe plane affine d'équation $y^2 = P(x)$ a pour anneau de coordonnées $A = k[x, y]/(y^2 - P(x))$. Quand $P$ n'est pas un carré, $y^2 - P$ est irréductible, $A$ est intègre et la courbe a un corps de fonctions $K = k(x)[y]$, extension de degré~$2$ de $k(x)$. La courbe est \emph{normale} si $A$ est intégralement clos dans $K$ : toute fraction de $K$ racine d'un polynôme unitaire à coefficients dans $A$ est déjà dans $A$. Pour une courbe sur un corps parfait, être normal revient à n'avoir aucun point singulier. Le cas type d'échec est la cubique cuspidale $y^2 = x^3$ : la fraction $t = y/x$ vérifie $t^2 = x$, elle est entière sur $A$ sans lui appartenir, et la courbe a un point de rebroussement à l'origine.

\emph{La page.} À la page~39, la rédaction affirme que, pour $P$ irréductible, la courbe $y^2 = P(x)$ est toujours normale, sans condition sur la caractéristique de $k$. En caractéristique différente de~$2$, c'est juste. En caractéristique~$2$, c'est faux.

\emph{Pourquoi la caractéristique~$2$ fait exception.} Le critère jacobien localise les points singuliers de la courbe $F(x, y) = y^2 - P(x) = 0$ là où les deux dérivées partielles s'annulent. On a $\partial F / \partial y = 2y$ et $\partial F / \partial x = -P'(x)$. Hors de la caractéristique~$2$, la première ne s'annule que pour $y = 0$, la seconde demande alors $P'(x) = 0$ en une racine $x$ de $P$ : une racine double, impossible pour $P$ sans facteur carré lorsque $k$ est parfait. En caractéristique~$2$, la dérivée $2y$ est identiquement nulle, et tout point de la courbe où $P'$ s'annule est singulier, que $P$ ait une racine double ou non. Sur $\mathbf F_2$, pour $P = x^3 + x + 1$, on a $P' = 3x^2 + 1 = (x + 1)^2$, qui s'annule en $x = 1$, et $P(1) = 1$ : le point $(1, 1)$ de la courbe est singulier.

\begin{encadre}
\begin{contrexemple}[énoncé de l'édition]
Sur $k = \mathbf F_2$, le polynôme $P = x^3 + x + 1$ est irréductible, l'anneau $A = \mathbf F_2[x, y]/(y^2 - P(x))$ est intègre, et il n'est pas intégralement clos : la fraction $t = (y + 1)/(x + 1)$ vérifie $t^2 = x$ et n'appartient pas à $A$. L'assertion de la page~39 est donc fausse en caractéristique~$2$.
\end{contrexemple}
\end{encadre}

\begin{proof}
Le polynôme $P$ est de degré~$3$ et n'a pas de racine dans $\mathbf F_2$ ($P(0) = P(1) = 1$) ; il est donc irréductible. Ce n'est pas un carré dans $\mathbf F_2[x]$, son degré étant impair. Le polynôme $y^2 - P$, unitaire en $y$, est alors irréductible sur $\mathbf F_2[x]$ : une racine $r$ dans le corps $\mathbf F_2(x)$ serait entière sur l'anneau factoriel, donc intégralement clos, $\mathbf F_2[x]$, donc dans $\mathbf F_2[x]$, et $r^2 = P$ serait un carré. Comme $\mathbf F_2[x, y]$ est factoriel, $y^2 - P$ est premier, et $A$ est intègre.

En caractéristique~$2$, $(y + 1)^2 = y^2 + 1 = x^3 + x = x(x + 1)^2$. Dans le corps des fractions de $A$, $x + 1$ n'est pas nul, car sa valeur au point $(0, 1)$ de la courbe vaut~$1$. La fraction $t = (y + 1)/(x + 1)$ vérifie donc $t^2 = x$ : elle est racine du polynôme unitaire $T^2 - x$, et entière sur $A$. Il reste à voir que $x + 1$ ne divise pas $y + 1$ dans $A$. On utilise l'anneau des nombres duaux $\mathbf F_2[\varepsilon]$, $\varepsilon^2 = 0$, et l'homomorphisme $\psi : A \to \mathbf F_2[\varepsilon]$ qui envoie $x$ sur $1$ et $y$ sur $1 + \varepsilon$ ; il est bien défini, car $(1 + \varepsilon)^2 = 1 + 2\varepsilon + \varepsilon^2 = 1 = P(1)$. C'est le point singulier $(1, 1)$, muni d'un vecteur tangent vertical (figure~\ref{fig:bourbaki-197}). Alors $\psi(x + 1) = 0$ et $\psi(y + 1) = \varepsilon \neq 0$ ; si $y + 1 = s(x + 1)$, on aurait $\varepsilon = \psi(s) \cdot 0 = 0$. Donc $t \notin A$.
\end{proof}

\begin{figure}[htbp]
\centering
\begin{tikzpicture}[>=stealth, every node/.style={font=\small}]
\draw[black!35] (0,0) -- (2,0) (0,0) -- (0,2) (2,0) -- (2,2) (0,2) -- (2,2);
\draw[black!70] (0,0) circle (2.2pt);
\draw[black!70] (2,0) circle (2.2pt);
\fill[frise] (0,2) circle (2.2pt);
\fill[frise] (2,2) circle (2.2pt);
\draw[piste] (2,2) circle (5pt);
\draw[->, piste, thick] (2,2.18) -- (2,3.0) node[right] {$\psi$ : $x \mapsto 1$, $y \mapsto 1 + \varepsilon$};
\node[below] at (0,-0.1) {$0$};
\node[below] at (2,-0.1) {$1$};
\node[left] at (-0.1,0) {$0$};
\node[left] at (-0.1,2) {$1$};
\node[below] at (1,-0.55) {$x$};
\node[left] at (-0.6,1) {$y$};
\node[frise, above left] at (0,2.05) {$(0, 1)$};
\node[frise, right=6pt] at (2,1.8) {$(1, 1)$, singulier};
\end{tikzpicture}
\caption{Les quatre points de $\mathbf F_2^2$ ; ceux de la courbe $y^2 = x^3 + x + 1$ sont $(0, 1)$ et $(1, 1)$ (en bleu), puisque $P(0) = P(1) = 1$. L'homomorphisme $\psi : A \to \mathbf F_2[\varepsilon]$ de la preuve est le point singulier $(1, 1)$ muni d'un vecteur tangent vertical : $\psi(x + 1) = 0$ et $\psi(y + 1) = \varepsilon$, de sorte que $x + 1$ ne peut diviser $y + 1$. Au point $(0, 1)$, au contraire, $x + 1$ vaut $1$.}
\label{fig:bourbaki-197}
\end{figure}

\emph{Vérification en Lean.} Lean vérifie le contre-exemple dans \href{\leanblobmain/Redaction197.lean}{\texttt{Redaction197.\allowbreak{}lean}}, où l'anneau $A$ est construit comme $\mathbf F_2[x][y]/(y^2 - P)$ (\texttt{AdjoinRoot} de mathlib). L'irréductibilité de $P$ est \texttt{P\_irreducible}, celle de $y^2 - P$ est \texttt{f\_irreducible}, par le critère de mathlib \texttt{Monic.\allowbreak{}irreducible\_iff\_irreducible\_map\_fraction\_map}, et l'intégrité de $A$ est l'instance \texttt{isDomain}. L'identité $(y + 1)^2 = x(x + 1)^2$ est \texttt{key} ; l'homomorphisme vers les nombres duaux est \texttt{ψ}, et la non-divisibilité \texttt{not\_dvd}. La fraction $t$ est entière (\texttt{isIntegral\_t}) et ne vient pas de $A$ (\texttt{t\_not\_mem}) ; d'où \texttt{not\_isIntegrallyClosed}, et l'énoncé final \texttt{not\_page39}, qui réunit l'irréductibilité de $P$, l'intégrité de $A$ et l'échec de la normalité.

\emph{L'énoncé juste.} Hors de la caractéristique~$2$, l'assertion de la page est juste, et vaut plus généralement pour $P$ sans facteur carré.

\begin{encadre}
\begin{theoreme*}[Hartshorne, \emph{Algebraic Geometry}, II, exercice~6.4]
Soient $k$ un corps de caractéristique différente de~$2$ et $P \in k[x]$ un polynôme non constant sans facteur carré. Alors l'anneau $k[x, y]/(y^2 - P(x))$ est intégralement clos.
\end{theoreme*}
\end{encadre}

\begin{demopage}[esquisse de l'édition]
Le corps des fractions est $k(x)(y)$, de base $1, y$ sur $k(x)$. Un élément $u = a + by$, $a, b \in k(x)$, a pour trace $2a$ et pour norme $a^2 - b^2 P$ ; il est entier sur $k[x]$ si et seulement si ces deux nombres sont dans $k[x]$. Comme $2$ est inversible, $a \in k[x]$, puis $b^2 P \in k[x]$, et $P$ sans facteur carré force $b \in k[x]$. La clôture intégrale de $k[x]$ dans $k(x)(y)$ est donc $k[x] \oplus k[x]\,y = A$.
\end{demopage}

\noindent L'esquisse montre où sert l'hypothèse : en caractéristique~$2$, la trace $2a$ est nulle et ne contraint plus $a$. Dans l'exemple, $t = (y + 1)/(x + 1) = 1/(x+1) + y/(x+1)$ a pour trace $0$ et pour norme $(1 + P)/(x + 1)^2 = x$, toutes deux polynomiales, sans que ses coordonnées le soient.

Ne sont pas démontrés ici : l'énoncé juste en caractéristique différente de~$2$ ; l'équivalence entre normalité et absence de point singulier pour une courbe, qui ne sert qu'à l'intuition ; la description de la normalisation de la courbe de l'exemple.

\litterature{l'énoncé juste est un exercice de manuel (Hartshorne, \emph{Algebraic Geometry}, chapitre~II, exercice~6.4 \cite{hartshorne1977}), qui suppose explicitement la caractéristique différente de~$2$. Le contre-exemple est élémentaire et l'édition ne le donne pas pour nouveau. Ce n'est pas une piste.}

\section[Rédaction Bourbaki n\textsuperscript{o}\,197 : coordonnées de Chow d'un produit]{Rédaction Bourbaki n\textsuperscript{o}\,197 : les coordonnées de Chow d'un produit\protect\verifie{Redaction197Chow}}\label{sec:bourbaki-197-chow}

La rédaction n\textsuperscript{o}\,197 des archives Bourbaki, de P.~Samuel, datée de mai 1954 par sa notice (Archives Henri Poincaré, \ahp{612}), traite des coordonnées de Chow des cycles de l'espace projectif. Sa transcription reste locale pour l'instant ; on n'en cite ici aucun passage.

\emph{Coordonnées de Chow.} Un cycle $X$ de dimension $r$ et de degré $d$ de l'espace projectif $\mathbf P^n$ rencontre en $d$ points, comptés avec multiplicité, un sous-espace linéaire générique de codimension $r+1$. La condition pour que $r+1$ hyperplans $u^{(0)}, \ldots, u^{(r)}$ et $X$ se rencontrent est l'annulation d'une forme $F_X(u^{(0)}, \ldots, u^{(r)})$, homogène de degré $d$ en chaque $u^{(i)}$ : c'est la forme de Chow de $X$, définie à un scalaire près, et ses coefficients sont les coordonnées de Chow de $X$. Pour un cycle de dimension nulle, c'est-à-dire une somme de points, la forme est le produit des formes linéaires $\sum_k u_k x_k$ évaluées aux points. Pour trois points $(x_j, y_j)$ du plan affine, vus dans $\mathbf P^2$, elle vaut
\[
\prod_{j=1}^3 (v_0 + v_1 x_j + v_2 y_j) = \sum_{p+q \leq 3} e_{pq}(x, y)\, v_0^{3-p-q} v_1^p v_2^q ,
\]
où $e_{pq}(x, y)$ est la somme, sur les façons de choisir $p$ des trois points pour fournir leur $x$ et $q$ autres pour fournir leur $y$, du produit des coordonnées choisies : $e_{10} = x_1 + x_2 + x_3$, $e_{11} = \sum_{i \neq j} x_i y_j$, $e_{21} = x_1x_2y_3 + x_1x_3y_2 + x_2x_3y_1$, et $e_{00} = 1$. Ce sont les fonctions multisymétriques élémentaires des trois points ; elles sont symétriques sous les permutations des points, mais n'engendrent pas toujours l'anneau de toutes les fonctions symétriques, comme on le verra.

\emph{La page.} La page~82 commence par achever une démonstration : l'équation de la projection semi-générique de $X \times Y$ s'obtient par un procédé universel à partir de celles de $X$ et de $Y$. Dans le cas irréductible, cette équation est le résultant, par rapport à une variable, de l'équation de $X$ et de celle de $Y$ décalée par un changement de variable affine, et la multiplicativité du résultant étend l'énoncé au cas général. C'est le principe « les racines s'ajoutent » qui donne, pour les points, la forme de Chow de l'adhérence projective d'un produit de deux parties de l'espace affine. Vient ensuite la question qui nous occupe. D'après la page~82, il paraît vraisemblable au rédacteur que les coordonnées de Chow d'un produit $X \times Y$ de cycles projectifs, de degrés et de dimensions donnés, s'expriment par des formules universelles à coefficients entiers en celles de $X$ et de $Y$. La page ne dit pas comment $X \times Y$ est plongé ; on la lit, comme il est naturel pour un produit d'espaces projectifs, avec le plongement de Segre, qui fait de $X \times Y$ un cycle d'un espace projectif avec ses propres coordonnées de Chow. La page ne démontre pas cette assertion. On montre ci-dessous qu'elle tient sur un corps pour les cycles de dimension nulle, et qu'elle tombe sur $\mathbf Z$, pour le plongement de Segre, dès le cas de deux cycles de trois points du plan.

\emph{Le produit de deux cycles de points.} Soient $A$ le cycle des points $w_i = (c_i, d_i)$ et $B$ celui des points $q_j = (a_j, b_j)$, $i, j = 1, 2, 3$. Le produit $A \times B$ est formé des neuf points $(w_i, q_j)$, et le plongement de Segre $\mathbf P^2 \times \mathbf P^2 \to \mathbf P^8$ envoie le point $\bigl((1 : c_i : d_i), (1 : a_j : b_j)\bigr)$ sur le point dont les neuf coordonnées sont les produits $X_k Y_l$, où $X = (1, c_i, d_i)$ et $Y = (1, a_j, b_j)$. La forme de Chow du produit est donc $\prod_{i,j} \sum_{k,l} u_{kl} X_k Y_l$, en neuf variables $u_{kl}$. Donnons la valeur~$1$ à $u_{00}$, $u_{11}$ et $u_{22}$, et la valeur~$0$ aux six autres : on obtient le polynôme
\[
T = \prod_{i, j = 1}^{3} \bigl(1 + c_i a_j + d_i b_j\bigr) \in \mathbf Z[a, b, c, d],
\]
à douze variables. Ses coefficients sont des sommes de coordonnées de Chow du produit. Si l'énoncé de la page tenait, $T$ serait donc un polynôme à coefficients entiers en les neuf fonctions $e_{pq}(c, d)$ et les neuf fonctions $e_{pq}(a, b)$, avec $1 \leq p + q \leq 3$.

\emph{Le cas d'un corps.} Notons $\mathrm{TS}_w$ l'anneau des polynômes en $c, d$ symétriques sous les permutations des trois points $w_i$, et $C_w \subseteq \mathrm{TS}_w$ la sous-algèbre engendrée par les $e_{pq}(c, d)$ ; de même $\mathrm{TS}_q$ et $C_q$ pour les points $q_j$. Pour chaque $i$, le facteur $\prod_j (1 + c_i a_j + d_i b_j)$ est la forme de Chow de $B$ évaluée en $(1, c_i, d_i)$, c'est-à-dire $\sum_{p,q} e_{pq}(a, b)\, c_i^p d_i^q$. Le produit sur $i$ est donc un polynôme en les $c_i, d_i$, symétrique sous les permutations des $w_i$, dont les coefficients sont dans $C_q$ : $T$ appartient à $\mathrm{TS}_w \otimes C_q$. En regroupant les facteurs selon $j$ plutôt que selon $i$, on obtient de même $T \in C_w \otimes \mathrm{TS}_q$. Sur un corps $k$, pour des sous-espaces $U \subseteq F$ et $V \subseteq F'$, on a $(U \otimes F') \cap (F \otimes V) = U \otimes V$ : il suffit de choisir une base de $F$ qui prolonge une base de $U$ et une base de $F'$ qui prolonge une base de $V$, et de lire les coefficients d'un tenseur dans la base produit. D'où $T \in C_w \otimes C_q$, ce qui est l'énoncé de la page pour les cycles de dimension nulle. Le même argument vaut pour des cycles de $m$ et $m'$ points de l'espace affine de dimension quelconque.

\emph{Pourquoi l'argument tombe sur $\mathbf Z$.} Sur $\mathbf Z$, dans chaque bidegré, $C_w$ et $\mathrm{TS}_w$ sont des réseaux de même rang, le premier contenu dans le second, car sur $\mathbf Q$ les fonctions élémentaires engendrent toutes les fonctions symétriques. Pour des réseaux, l'intersection ne se comporte plus comme un produit d'indices. Dans le cas le plus simple, $\mathrm{TS} = \mathbf Z$ et $C = 3\mathbf Z$ : les deux réseaux $C \otimes \mathrm{TS}$ et $\mathrm{TS} \otimes C$ valent $3\mathbf Z$, leur intersection aussi, alors que $C \otimes C = 9\mathbf Z$. L'intersection est d'indice $3$, le plus petit commun multiple des indices, et non $9$, leur produit. L'argument précédent ne place alors $T$ que dans l'intersection, qui peut être plus grande que $C_w \otimes C_q$. Pour que cet écart se produise, il faut que $C$ soit strictement plus petit que $\mathrm{TS}$. D.~Rydh a déterminé quand les fonctions multisymétriques élémentaires engendrent toutes les fonctions multisymétriques sur un anneau $A$, et, pour trois points du plan, il faut et il suffit que $3$ soit inversible dans $A$ (\cite[corollaire~8.10]{rydh2007}, qui reprend un théorème d'E.~Briand \cite{briand2004}). La remarque~8.11 du même article relie ce défaut au schéma de Chow des cycles de dimension nulle. Sur $\mathbf Z$, l'écart existe donc, et le calcul qui suit montre que $T$ tombe dedans.

\emph{Le certificat.} Écrivons un monôme en $x_1, y_1, x_2, y_2, x_3, y_3$ comme la suite de ses trois couples d'exposants, $((\alpha_1, \beta_1), (\alpha_2, \beta_2), (\alpha_3, \beta_3))$. La forme linéaire $\alpha$ sur $\mathbf Z[x, y]$ lit huit coefficients avec des poids :
\[
\begin{array}{lr@{\qquad}lr}
((1,0),(1,0),(0,2)) & 1, & ((2,0),(0,1),(0,1)) & -2,\\
((1,1),(1,0),(0,1)) & -1, & ((2,0),(0,2),(0,0)) & -1,\\
((1,1),(1,1),(0,0)) & -2, & ((2,1),(0,1),(0,0)) & -1,\\
((1,2),(1,0),(0,0)) & -1, & ((2,2),(0,0),(0,0)) & -2.
\end{array}
\]
Ces huit monômes sont un représentant de chaque orbite, sous les permutations des points, des monômes de bidegré $(2,2)$, c'est-à-dire de degré $2$ en les $x$ et $2$ en les $y$. Sur un polynôme symétrique, $\alpha$ est donc une combinaison des coefficients des huit fonctions symétriques monomiales de ce bidegré. Comme $e_{pq}$ est de bidegré $(p, q)$, un produit de fonctions élémentaires de bidegré $(2,2)$ est l'un des huit produits $e_{01}^2 e_{10}^2$, $e_{01}^2 e_{20}$, $e_{01}e_{10}e_{11}$, $e_{01}e_{21}$, $e_{02}e_{10}^2$, $e_{02}e_{20}$, $e_{10}e_{12}$ et $e_{11}^2$, et $\alpha$ y prend les valeurs
\[
-21,\quad -6,\quad -12,\quad 0,\quad -9,\quad -3,\quad -3,\quad -9,
\]
toutes divisibles par $3$. Sur un produit de fonctions élémentaires d'un autre bidegré, $\alpha$ est nulle, puisqu'elle ne lit que des monômes de bidegré $(2,2)$. Elle prend donc des valeurs divisibles par $3$ sur tout polynôme à coefficients entiers en les $e_{pq}$. Sur la fonction symétrique monomiale de $x_1x_2y_3^2$, somme des trois monômes de son orbite, elle vaut $1$ : cette fonction est symétrique et n'est pas un polynôme entier en les $e_{pq}$, ce qui rend visible le défaut de la remarque de Rydh en ce bidegré (ce dernier calcul n'est pas formalisé).

\begin{encadre}
\begin{contrexemple}[réponse à la question de la rédaction \redac{197}, p.~82, cycles de trois points, énoncé de l'édition]
Soit $T = \prod_{i,j=1}^3 (1 + c_i a_j + d_i b_j) \in \mathbf Z[a, b, c, d]$. Il n'existe aucun polynôme $G$ à coefficients entiers en dix-huit variables tel que $T = G\bigl((e_{pq}(c, d))_{1 \leq p+q \leq 3}, (e_{pq}(a, b))_{1 \leq p+q \leq 3}\bigr)$. Autrement dit, $T$ n'appartient pas au sous-anneau de $\mathbf Z[a,b,c,d]$ engendré par les fonctions multisymétriques élémentaires des deux cycles.
\end{contrexemple}
\end{encadre}

\begin{proof}
Soit $\beta = \alpha \otimes \alpha$ la forme linéaire sur $\mathbf Z[a, b, c, d]$ qui applique $\alpha$ aux variables $(c, d)$ et aux variables $(a, b)$ : sur un polynôme, elle fait la somme pondérée par $w_m w_{m'}$ des coefficients des monômes produits d'un monôme $m$ lu par $\alpha$ en $(c, d)$ et d'un monôme $m'$ lu par $\alpha$ en $(a, b)$. Sur un produit $P(c, d)\, Q(a, b)$, elle vaut $\alpha(P)\, \alpha(Q)$. Un polynôme $G$ en les deux familles de fonctions élémentaires est une combinaison entière de produits $E(c, d)\, E'(a, b)$, où $E$ et $E'$ sont des produits de fonctions élémentaires, et $\beta$ y vaut $\alpha(E)\, \alpha(E')$, multiple de $3 \cdot 3 = 9$. Donc $\beta(G) \in 9\mathbf Z$.

Calculons $\beta(T)$. En regroupant les facteurs de $T$ selon $j$, on a $T = \prod_j \sum_{p,q} e_{pq}(c, d)\, a_j^p b_j^q$, de sorte que le coefficient en $(a, b)$ du monôme $((p_1, q_1), (p_2, q_2), (p_3, q_3))$ est $e_{p_1q_1} e_{p_2q_2} e_{p_3q_3}(c, d)$. Appliquer $\alpha$ en $(a, b)$ donne donc
\[
e_{10}^2 e_{02} - e_{11} e_{10} e_{01} - 2\, e_{11}^2 - e_{12} e_{10} - 2\, e_{20} e_{01}^2 - e_{20} e_{02} - e_{21} e_{01} - 2\, e_{22},
\]
avec $e_{22} = 0$ puisque $2 + 2 > 3$. Appliquer ensuite $\alpha$ en $(c, d)$ et lire les huit valeurs du certificat donne
\[
\beta(T) = -9 + 12 + 18 + 3 + 12 + 3 - 0 - 0 = 39 .
\]
Or $39$ n'est pas divisible par $9$, et $T$ n'est pas de la forme $G(e(c, d), e(a, b))$.
\end{proof}

Le nombre $39 = 3 \cdot 13$ est divisible par $3$, comme il le doit : $T$ appartient à l'intersection $(\mathrm{TS}_w \otimes C_q) \cap (C_w \otimes \mathrm{TS}_q)$, sur laquelle $\beta$ prend des valeurs divisibles par~$3$. C'est exactement l'écart entre le plus petit commun multiple et le produit des indices décrit plus haut.

\begin{sloppypar}
\emph{Vérification en Lean.} Le fichier \href{\leanblobmain/Redaction197Chow.lean}{\texttt{Redaction197Chow.\allowbreak{}lean}} écrit l'anneau $\mathbf Z[a, b, c, d]$ comme $\mathbf Z[c, d][a, b]$, définit les $e_{pq}$ par la liste de leurs monômes et les caractérise par l'identité $\prod_i (1 + c_i U + d_i W) = \sum_{p,q} e_{pq}\, U^p W^q$ (\texttt{prod\_eq\_sum\_e}). Les coefficients sont calculés sur des listes de monômes, que le noyau de Lean multiplie et lit lui-même : les huit valeurs de $\alpha$ (\texttt{alpha\_e\_products}) et $\beta(T) = 39$ (\texttt{beta\_T}) sont ainsi vérifiées par calcul. La divisibilité par~$3$ sur tout polynôme en les $e_{pq}$ (\texttt{three\_dvd\_alpha\_aeval}) passe par les bidegrés, qui s'additionnent dans les produits de listes, et par la résolution des deux équations de bidegré en les exposants. La divisibilité par~$9$ (\texttt{nine\_dvd\_beta\_aeval}) et le contre-exemple (\texttt{not\_exists\_integral\_formula}, et \texttt{T\_not\_mem\_adjoin} pour le sous-anneau engendré) en découlent.
\end{sloppypar}

\emph{Ce qui n'est pas démontré.} Le cas des cycles de dimension positive n'est pas traité : la forme de Chow d'un produit y est plus difficile à écrire, et l'on ne sait pas ici si l'énoncé de la page tient sur un corps dans ce cas, ni s'il tombe sur $\mathbf Z$ pour d'autres raisons que celle-ci. Le cas d'un corps pour les cycles de dimension nulle est démontré plus haut en prose et n'est pas formalisé ; le passage de la forme de Chow du produit de Segre à $T$, par spécialisation des variables $u_{kl}$, ne l'est pas non plus, et $T$ est défini directement en Lean.

\litterature{piste candidate. Le défaut de génération des fonctions multisymétriques par les élémentaires sur $\mathbf Z$ est connu \cite{briand2004, rydh2007}, et Rydh en tire des conséquences pour le schéma de Chow des cycles de dimension nulle. Un contre-exemple explicite à l'expression entière des coordonnées de Chow d'un produit n'a pas été retrouvé dans les sources lues ; il reste à lire la thèse de Rydh sur les familles de cycles et le schéma de Chow, que son article cite pour ces questions, et les traitements classiques des formes de Chow des produits.}

\section[Rédaction Bourbaki n\textsuperscript{o}\,457 : degré et nombre de points]{Rédaction Bourbaki n\textsuperscript{o}\,457 : degré et nombre de points d'une algèbre séparable\protect\verifie{Redaction457}}\label{sec:bourbaki-457}

La rédaction n\textsuperscript{o}\,457 des archives Bourbaki, datée d'octobre 1965 par sa notice (Archives Henri Poincaré, \href{https://archives-bourbaki.ahp-numerique.fr/items/show/792}{item 792}), qui la donne à Grothendieck, est un projet de réédition du chapitre~V d'\emph{Algèbre}, consacré aux corps commutatifs. Elle y substitue l'étude des algèbres à celle des extensions, et classe les algèbres étales sur un corps par les ensembles finis sur lesquels opère le groupe de Galois. Sa transcription reste locale pour l'instant ; on n'en cite ici aucun passage.

\emph{Les points d'une algèbre.} Soient $k$ un corps et $\Omega$ une extension algébriquement close de $k$. Pour une $k$-algèbre commutative $A$, on note $P(A)$ l'ensemble des homomorphismes de $k$-algèbres de $A$ dans $\Omega$. En langage géométrique, ce sont les points de $\operatorname{Spec} A$ à valeurs dans $\Omega$. Pour $A = k[X]/(f)$, un homomorphisme $A \to \Omega$ est déterminé par l'image de $X$, qui doit être une racine de $f$ : les points de $A$ sont les racines de $f$ dans $\Omega$. Pour $k = \mathbf Q$ et $f = X^2 - 2$, on trouve les deux plongements de $\mathbf Q(\sqrt 2)$ dans $\mathbf C$.

\emph{En degré fini, la dimension compte les points.} Si $f$ est séparable de degré $n$, il a $n$ racines distinctes dans $\Omega$, et $A = k[X]/(f)$ a exactement $n$ points pour une dimension $n$. Plus généralement, une algèbre étale de degré $n$ devient, après extension des scalaires à $\Omega$, le produit $\Omega^n$, dont les $n$ projections sont les points. La séparabilité est indispensable : pour $k$ de caractéristique $p$ et $a \in k$ qui n'est pas une puissance $p$-ième, l'algèbre $k[X]/(X^p - a)$ est de dimension $p$ et n'a qu'un point, car $X^p - a$ n'a qu'une racine dans $\Omega$.

\emph{La page.} Le corollaire~6 de la proposition~14 (PDF, p.~32--33) considère une algèbre $A$ entière et séparable sur $k$, c'est-à-dire réunion filtrante de sous-algèbres étales de degré fini $A_i$, et affirme que son degré $[A:k]$ est égal au cardinal de $P(A)$ ; en particulier, $A$ est de degré fini si et seulement si $P(A)$ est fini. La démonstration ne traite que ce dernier point. Les $P(A_i)$ sont finis, $P(A)$ est leur limite projective, et les applications de transition sont surjectives ; un lemme sur les limites projectives d'ensembles finis non vides donne alors des surjections de $P(A)$ sur chaque $P(A_i)$, dont les cardinaux tendent vers l'infini. On a donc montré que $P(A)$ est infini lorsque $A$ est de degré infini, et rien de plus. Une note dactylographiée le reconnaît : le rédacteur y constate qu'il a donné à \q{card} la seule valeur $\infty$ pour tout ensemble infini. Un lecteur a écrit à l'encre, à droite de la formule, que l'égalité est fausse en dimension infinie, avec les mots \q{dénombrable} et \q{continu}. Le contre-exemple qui suit rend précise cette objection.

\emph{Pourquoi l'égalité échoue en degré infini.} Une base dénombrable ne voit que des combinaisons finies, tandis qu'un point de $A$ est un choix compatible à tous les étages $A_i$ à la fois. L'ensemble de Cantor $C = \{0, 1\}^{\mathbf N}$ montre l'écart. Ses parties à la fois ouvertes et fermées sont les réunions finies de cylindres, un cylindre étant l'ensemble des suites qui commencent par un mot fini donné ; il n'y en a qu'une infinité dénombrable. Ses points sont les suites infinies de $0$ et de $1$, qui ont la puissance du continu. Les fonctions localement constantes $C \to k$ sont les combinaisons finies d'indicatrices de cylindres : elles forment une algèbre de dimension dénombrable, et chaque point de $C$ en donne un point, l'évaluation en ce point. Étage par étage, les fonctions qui ne dépendent que des $n$ premiers termes forment l'algèbre étale $k^{\{0,1\}^n}$, de degré $2^n$ et à $2^n$ points ; chaque étage double les deux nombres, mais à la limite la dimension reste dénombrable et les points deviennent les suites infinies (figure~\ref{fig:457-cantor}).

\begin{figure}[htbp]
\centering
\begin{tikzpicture}[x=0.62cm, y=0.95cm, >=stealth, every node/.style={font=\small}]
% le cylindre [01] : les branches qui passent par le noeud 01
\fill[piste!12] (-2,-2) -- (-3.95,-4.75) -- (-0.05,-4.75) -- cycle;
\foreach \d in {0,...,3} {
  \pgfmathtruncatemacro{\n}{2^\d - 1}
  \foreach \i in {0,...,\n} {
    \pgfmathsetmacro{\x}{(\i + 0.5)*2^(4 - \d) - 8}
    \pgfmathsetmacro{\xl}{(2*\i + 0.5)*2^(3 - \d) - 8}
    \pgfmathsetmacro{\xr}{(2*\i + 1.5)*2^(3 - \d) - 8}
    \draw[black!55] (\x,-\d) -- (\xl,-\d-1);
    \draw[black!55] (\x,-\d) -- (\xr,-\d-1);
    \fill (\x,-\d) circle (1.6pt);
  }
}
\foreach \i in {0,...,15} { \fill (\i + 0.5 - 8,-4) circle (1.6pt); \draw[black!40, dotted] (\i + 0.5 - 8,-4) -- (\i + 0.5 - 8,-4.6); }
\node[above left=-1pt] at (-2,-0.5) {\footnotesize $0$};
\node[above right=-1pt] at (2,-0.5) {\footnotesize $1$};
% la branche 1011...
\draw[frise, very thick] (0,0) -- (4,-1) -- (2,-2) -- (3,-3) -- (3.5,-4);
\draw[frise, very thick, dotted] (3.5,-4) -- (3.5,-4.6);
\node[frise, below, align=center] at (3.5,-4.6) {\footnotesize branche $1011\ldots$,\\\footnotesize un point de $C$};
\node[piste, below, align=center] at (-2,-4.6) {\footnotesize cylindre des suites\\\footnotesize qui commencent par $01$};
\fill[piste] (-2,-2) circle (2.2pt);
\node[piste, above left] at (-2,-2) {\footnotesize $01$};
\foreach \d in {0,...,4} { \pgfmathtruncatemacro{\p}{2^\d} \node[black!65, anchor=west] at (8.4,-\d) {\footnotesize $2^{\d} = \p$}; }
\end{tikzpicture}
\caption{L'arbre binaire des mots finis, jusqu'à la longueur~$4$. Chaque n\oe ud est un mot $v$ et donne le cylindre des suites qui commencent par $v$ : il y en a $2^n$ à l'étage~$n$, une infinité dénombrable en tout, et leurs indicatrices engendrent l'algèbre $A$. Un point de $C$ est une branche infinie de l'arbre, un choix à chaque étage : il y en a la puissance du continu. La branche dessinée appartient aux cylindres $1$, $10$, $101$, $1011$, et non au cylindre $01$.}
\label{fig:457-cantor}
\end{figure}

\begin{contrexemple}[Bourbaki, rédaction n\textsuperscript{o}\,457, corollaire~6]
Soient $k$ un corps, $C = \{0, 1\}^{\mathbf N}$ et $A$ la sous-$k$-algèbre de $k^C$ formée des fonctions qui ne dépendent que d'un nombre fini de termes de la suite. Tout élément de $A$ est racine d'un polynôme unitaire séparable scindé sur $k$ ; en particulier $A$ est entière et séparable sur $k$. On a $[A:k] \leq \aleph_0$, et, pour toute $k$-algèbre non nulle $\Omega$, en particulier toute extension algébriquement close de $k$, l'ensemble des homomorphismes de $k$-algèbres $A \to \Omega$ a au moins la puissance du continu. Donc $[A:k] < \operatorname{card} P(A)$.
\end{contrexemple}

\begin{proof}
Si $f$ ne dépend que des $n$ premiers termes, elle ne prend qu'un nombre fini de valeurs $c_1, \dots, c_r$, deux à deux distinctes, et $f$ est racine de $\prod_j (X - c_j)$, qui est unitaire, scindé sur $k$ et à racines simples, donc séparable : au point $z$, le facteur $f(z) - c_j$ correspondant à la valeur $f(z)$ est nul. Pour chaque mot $v$ de longueur $n$, soit $\chi_v$ l'indicatrice du cylindre des suites qui commencent par $v$. Alors $f = \sum_v f(\bar v)\, \chi_v$, où $\bar v$ est la suite $v$ prolongée par des zéros : en un point $z$, seul le terme du mot formé des $n$ premiers termes de $z$ subsiste, et $f(\bar v) = f(z)$ puisque $\bar v$ et $z$ ont les mêmes $n$ premiers termes. Les $\chi_v$, pour tous les mots finis $v$, forment une famille dénombrable qui engendre $A$, d'où $[A:k] \leq \aleph_0$. Pour un point $x$ de $C$, l'évaluation $f \mapsto f(x)$, suivie de l'application $k \to \Omega$, est un homomorphisme de $k$-algèbres $A \to \Omega$. Si $x \neq y$, ils diffèrent en un rang $m$, et la fonction qui vaut $1$ sur les suites de terme $1$ au rang $m$ et $0$ ailleurs appartient à $A$ et prend les valeurs $1$ et $0$ en $x$ et $y$, distinctes dans $\Omega$ non nulle. L'application $x \mapsto \mathrm{ev}_x$ est donc injective, et $\operatorname{card} P(A) \geq \operatorname{card} C = 2^{\aleph_0} > \aleph_0$.
\end{proof}

\emph{Vérification en Lean.} Lean vérifie l'énoncé ci-dessus : $A$ y est définie directement comme l'algèbre des fonctions qui ne dépendent que d'un nombre fini de termes, et l'énoncé final dit que, pour toute extension algébriquement close $\Omega$ de $k$, il existe une $k$-algèbre entière dont tout élément est racine d'un polynôme séparable scindé, et dont le degré diffère du cardinal de l'ensemble de ses homomorphismes dans $\Omega$. Trois points ne sont pas formalisés : l'identification de $A$ avec les fonctions localement constantes sur l'ensemble de Cantor, les égalités $[A:k] = \aleph_0$ et $P(A) \simeq C$ (pour $\Omega$ un corps), dont la réfutation n'a pas besoin, et la dernière assertion du corollaire, qui est juste. Le cas du degré fini, où le corollaire est juste aussi, est démontré plus bas.

\emph{Le cas du degré fini.} En degré fini, le corollaire est juste, et sa démonstration tient en trois remarques. Une algèbre étale sur un corps est un produit fini $A \simeq K_1 \times \cdots \times K_r$ d'extensions finies séparables. Un homomorphisme $A \to \Omega$ vers un corps se factorise par une seule projection : les idempotents $e_i = (0, \ldots, 1, \ldots, 0)$ vérifient $e_ie_j = 0$ pour $i \neq j$ et $\sum e_i = 1$, leurs images dans $\Omega$ sont des idempotents d'un corps, donc $0$ ou $1$, et une seule vaut~$1$. Les points de $A$ sont ainsi la réunion disjointe des points des $K_i$. Enfin, une extension finie séparable $K_i$ a exactement $[K_i : k]$ plongements dans $\Omega$ algébriquement clos : c'est le degré séparable, égal au degré.

\begin{encadre}
\begin{enonce}[Bourbaki, rédaction n\textsuperscript{o}\,457, corollaire~6, cas du degré fini]
Soient $k$ un corps, $\Omega$ une extension algébriquement close de $k$ et $A$ une $k$-algèbre étale. Alors $[A : k] = \operatorname{card} P(A)$.
\end{enonce}
\end{encadre}

\begin{proof}
Écrivons $A \simeq \prod_i K_i$. Les applications $(i, f) \mapsto f \circ \mathrm{pr}_i$, pour $f : K_i \to \Omega$, sont injectives : deux d'entre elles d'indices différents diffèrent sur un idempotent, et pour un même indice la projection est surjective. Il y a donc au moins $\sum_i [K_i : k] = [A : k]$ points. Il n'y en a pas plus : des homomorphismes distincts $A \to \Omega$ sont linéairement indépendants sur $\Omega$ (lemme de Dedekind), dans l'espace des applications $k$-linéaires $A \to \Omega$, qui est de dimension $[A : k]$ sur $\Omega$.
\end{proof}

\begin{sloppypar}
\noindent Lean le vérifie dans \href{\leanblobmain/Redaction457.lean}{\texttt{Redaction457.\allowbreak{}lean}} (\texttt{finrank\_eq\_card\_points}). La décomposition en produit est \texttt{Algebra.\allowbreak{}Etale.\allowbreak{}iff\_exists\_algEquiv\_prod} de mathlib, le décompte des plongements d'un corps vient du degré séparable (\texttt{Field.\allowbreak{}finSepDegree\_eq\_of\_isAlgClosed}), et la majoration est \texttt{card\_algHom\_le\_finrank}.
\end{sloppypar}

\begin{remarque}
Le même phénomène se rencontre pour les corps. L'extension $K$ de $\mathbf Q$ engendrée par les racines carrées des nombres premiers, $\sqrt 2$, $\sqrt 3$, $\sqrt 5$, etc., est de degré dénombrable, et ses plongements dans $\mathbf C$ correspondent aux éléments de son groupe de Galois $(\mathbf Z/2\mathbf Z)^{\mathbf N}$, qui a la puissance du continu. En degré infini, le cardinal de $P(A)$ n'est d'ailleurs pas fixé par le degré : pour un ensemble infini $D$, l'algèbre des fonctions $D \to k$ constantes hors d'une partie finie a pour degré le cardinal de $D$, et ses points, les évaluations en les éléments de $D$ et la valeur constante à l'infini, sont en même nombre. Ces deux exemples ne sont pas formalisés. La forme juste du corollaire est sa dernière assertion, que la démonstration de la page établit.
\end{remarque}

\section[Rédaction Bourbaki n\textsuperscript{o}\,582 : les indices de Schur sur $\mathbf R$]{Rédaction Bourbaki n\textsuperscript{o}\,582 : les indices de Schur sur $\mathbf R$\protect\verifie{Redaction582}}\label{sec:bourbaki-582}

La rédaction n\textsuperscript{o}\,582 des archives Bourbaki (Archives Henri Poincaré, \href{https://archives-bourbaki.ahp-numerique.fr/items/show/1130}{item 1130}) est le rapport de Jean-Pierre Serre, daté de juin 1972 par sa notice, sur la réédition du chapitre~VIII d'\emph{Algèbre}, consacré aux modules et anneaux semi-simples. Son n\textsuperscript{o}\,3 demande que le chapitre définisse les groupes de Grothendieck de modules et explique les indices de Schur ; seule cette partie est transcrite (PDF, p.~10--12). La rédaction n\textsuperscript{o}\,596 du même chapitre (section~\ref{sec:bourbaki-596}) consacre un paragraphe à ces groupes. Sa transcription reste locale pour l'instant ; on n'en cite ici aucun passage.

\emph{L'indice de Schur.} Soient $k$ un corps, $A$ une $k$-algèbre et $S$ un $A$-module simple de rang fini sur $k$. Par le lemme de Schur, $D = \operatorname{End}_A(S)$ est un corps, en général non commutatif ; son centre $L$ est un corps commutatif qui contient $k$. On pose $m = [L : k]$, et l'on écrit $[D : L] = e^2$ : la dimension d'un corps gauche sur son centre est toujours un carré, et $e$ s'appelle l'indice de Schur. Si l'on étend les scalaires à une extension galoisienne $k'$ de $k$ assez grande, $k' \otimes_k S$ se décompose en $m$ modules simples deux à deux non isomorphes, un pour chaque plongement de $L$ dans $k'$, et chacun y apparaît $e$ fois. Le rapport énonce cette décomposition à sa page~11. En comptant les dimensions, $[D : k] = [D : L]\,[L : k] = m e^2$.

\emph{Le cas réel.} Pour $k = \mathbf R$ et $k' = \mathbf C$, le théorème de Frobenius ne laisse que trois corps de dimension finie sur $\mathbf R$ : $\mathbf R$, $\mathbf C$ et le corps $\mathbf H$ des quaternions. Chacun se rencontre déjà pour des représentations de groupes finis. Le groupe $\mathbf Z/3$, agissant sur $\mathbf R^2$ par la rotation d'un tiers de tour, a pour commutant les similitudes directes, c'est-à-dire $\mathbf C$ ; après complexification, la rotation se diagonalise avec les valeurs propres $\omega = e^{2i\pi/3}$ et $\bar\omega$, et le module se scinde en deux caractères conjugués, distincts. Le groupe des quaternions $\{\pm 1, \pm i, \pm j, \pm k\}$, agissant sur $\mathbf H$ par multiplication à gauche, a pour commutant les multiplications à droite, isomorphes à $\mathbf H$ ; après complexification, $\mathbf H \otimes_{\mathbf R} \mathbf C \simeq M_2(\mathbf C)$ est la somme de deux copies d'un même module simple de dimension~$2$.

\emph{La page.} L'exemple de la page~12 donne les trois types de modules simples réels avec leurs invariants. Pour $D = \mathbf R$, il donne $(m, e) = (1, 1)$, ce qui est juste. Pour $D = \mathbf C$ comme pour $D = \mathbf H$, le tapuscrit donne $(m, e) = (2, 2)$. Ces valeurs ne suivent pas les définitions de la page~11. La suite de la page s'accorde d'ailleurs avec les valeurs justes. Elle décrit les modules simples complexes de deuxième type, non isomorphes à leur conjugué, qui viennent de $\mathbf C$ : $T$ et $\bar T$ y apparaissent chacun une fois. Ceux de troisième type, isomorphes à leur conjugué sans être réalisables sur $\mathbf R$, viennent de $\mathbf H$ et y apparaissent deux fois ; c'est pourquoi le critère de réalisabilité qui suit leur demande une multiplicité paire.

\begin{encadre}
\begin{enonce}[Rédaction n\textsuperscript{o}\,582, n\textsuperscript{o}\,3, l'exemple de $\mathbf R$, valeurs corrigées par l'édition]
Pour $D = \mathbf C$, on a $m = 2$ et $e = 1$. Pour $D = \mathbf H$, on a $m = 1$ et $e = 2$. Aucun des deux ne donne $(m, e) = (2, 2)$.
\end{enonce}
\end{encadre}

\begin{proof}
Le corps $\mathbf C$ est commutatif : son centre est $\mathbf C$, d'où $m = [\mathbf C : \mathbf R] = 2$ et $e^2 = [\mathbf C : \mathbf C] = 1$. Un quaternion qui commute avec $i$ et avec $j$ est réel, de sorte que le centre de $\mathbf H$ est $\mathbf R$ ; d'où $m = 1$ et $e^2 = [\mathbf H : \mathbf R] = 4$. Les valeurs tapées sont exclues sans même connaître les centres : elles donneraient $[D : \mathbf R] = m e^2 = 8$, qui n'est la dimension ni de $\mathbf C$ ni de $\mathbf H$.
\end{proof}

\noindent\emph{Vérification en Lean.} Lean vérifie les dimensions qui tranchent : $[\mathbf C : \mathbf R] = 2 = 2 \cdot 1^2$ (\texttt{schur\_complex}, par \texttt{Complex.\allowbreak{}finrank\_real\_complex} de mathlib) et $[\mathbf H : \mathbf R] = 4 = 1 \cdot 2^2$ (\texttt{schur\_quaternion}, par \texttt{Quaternion.\allowbreak{}finrank\_eq\_four}). Les centres, $\mathbf C$ pour $\mathbf C$ et $\mathbf R$ pour $\mathbf H$, sont pris comme connus et ne sont pas formalisés, non plus que la relation $[D : k] = m e^2$. Les noms sont ceux de \href{\leanblobmain/Redaction582.lean}{\texttt{Redaction582.\allowbreak{}lean}}.

Ne sont pas démontrés ici : la décomposition de $k' \otimes_k S$ et l'interprétation des multiplicités par l'indice de Schur ; le critère de réalisabilité sur $\mathbf R$ d'un module semi-simple complexe ; la règle qui lit les trois types sur le caractère d'une représentation d'un groupe fini, par l'indicateur de Frobenius--Schur ; les autres demandes du n\textsuperscript{o}\,3 sur les groupes de Grothendieck, la trace et l'extension des scalaires.

\section[Rédaction Bourbaki n\textsuperscript{o}\,596 : toute classe est une différence]{Rédaction Bourbaki n\textsuperscript{o}\,596 : dans un groupe de Grothendieck, toute classe est une différence\protect\verifie{Redaction596}}\label{sec:bourbaki-596}

La rédaction n\textsuperscript{o}\,596 des archives Bourbaki (Archives Henri Poincaré, \href{https://archives-bourbaki.ahp-numerique.fr/items/show/1142}{item 1142}) est une rédaction du chapitre~VIII d'\emph{Algèbre}, sur les modules et anneaux semi-simples, que sa notice donne à Jean-Louis Verdier et Hyman Bass et date de septembre 1973 à août 1974. Sa seconde partie, datée de mai 1974 sur sa table des matières, s'ouvre sur un §~10, \q{Groupes et anneaux de Grothendieck}. Des corrections à l'encre, d'une main non identifiée, reprennent le texte tapé. Seul le §~10 est transcrit, et sa transcription reste locale pour l'instant ; on n'en cite ici que son titre. Le rapport de Serre sur le même chapitre, deux ans plus tôt, demandait ce paragraphe (section~\ref{sec:bourbaki-582}).

\emph{Le groupe de Grothendieck.} Soit $A$ un anneau. La rédaction part d'un ensemble additif $\underline C$ de classes de $A$-modules : il contient la classe du module nul et, avec les classes de $E$ et de $E'$, celle de $E \oplus E'$. Une fonction additive sur $\underline C$, à valeurs dans un groupe abélien, est une fonction $\varphi$ des classes telle que $\varphi(E) = \varphi(E') + \varphi(E'')$ pour toute suite exacte $0 \to E' \to E \to E'' \to 0$ de modules de type $\underline C$. La longueur, sur les modules de longueur finie, et la dimension, sur les espaces vectoriels de dimension finie, en sont des exemples. Le groupe $K(\underline C)$ est le quotient du groupe abélien libre sur $\underline C$ par les relations $[E] = [E'] + [E'']$ des suites exactes. La classe $[E]$ est la fonction additive universelle (proposition~1) : $K(\underline C)$ est l'endroit où l'on peut soustraire des modules, et où une suite exacte compte comme une somme.

\emph{La proposition~3\,a).} Elle affirme que tout élément $x$ de $K(\underline C)$ s'écrit $[E] - [F]$, pour deux modules $E$ et $F$ de type $\underline C$. Ainsi un entier relatif est la différence de deux entiers naturels. La démonstration de la page tient en deux lignes. Par construction, $x$ est une combinaison à coefficients entiers de classes, donc s'écrit $\sum_i [E_i] - \sum_j [F_j]$. On regroupe chaque côté en une somme directe, $E = \bigoplus_i E_i$ et $F = \bigoplus_j F_j$, puisque $[E_1 \oplus E_2] = [E_1] + [E_2]$ ; cette égalité vient de la suite exacte scindée $0 \to E_1 \to E_1 \oplus E_2 \to E_2 \to 0$.

L'argument n'utilise que deux faits : l'application $\mathrm{cl}(E) \mapsto [E]$ est additive pour la somme directe, et son image engendre $K(\underline C)$. Il vaut donc dans la généralité suivante.

\begin{encadre}
\begin{enonce}[Rédaction n\textsuperscript{o}\,596, §~10, proposition~3\,a), forme abstraite de l'édition]
Soient $M$ un monoïde commutatif, $G$ un groupe abélien et $f : M \to G$ un morphisme de monoïdes dont l'image engendre $G$. Tout élément de $G$ s'écrit $f(a) - f(b)$, avec $a, b \in M$.
\end{enonce}
\end{encadre}

\begin{proof}
L'ensemble des différences $f(a) - f(b)$ contient l'image de $f$ (prendre $b = 0$) et l'élément neutre. Il est stable par addition, car $(f(a) - f(b)) + (f(c) - f(d)) = f(a + c) - f(b + d)$, et par passage à l'opposé, en échangeant $a$ et $b$. C'est donc un sous-groupe qui contient l'image de $f$, et il est égal à $G$.
\end{proof}

\noindent La proposition~3\,a) est le cas où $M$ est l'ensemble $\underline C$ muni de la somme directe, de neutre la classe du module nul, et $G = K(\underline C)$. L'additivité de $f$ est l'égalité $[E_1 \oplus E_2] = [E_1] + [E_2]$, et l'image engendre par construction.

\begin{remarque}
L'énoncé abstrait couvre deux groupes distincts. Le groupe complété du monoïde $(\underline C, \oplus)$ ne connaît que les suites exactes scindées ; $K(\underline C)$ en est le quotient par les relations de toutes les suites exactes. Pour les groupes abéliens finis, la suite $0 \to \mathbf Z/2 \to \mathbf Z/4 \to \mathbf Z/2 \to 0$ donne $[\mathbf Z/4] = [\mathbf Z/2 \oplus \mathbf Z/2]$ dans $K(\underline C)$. Dans le groupe complété, les deux classes restent distinctes : le monoïde des classes de groupes abéliens finis est simplifiable, par le théorème de structure, et s'injecte dans son groupe complété. La proposition~3\,b) de la rédaction décrit justement, par des suites exactes, quand deux modules ont la même classe dans $K(\underline C)$.
\end{remarque}

\noindent\emph{Vérification en Lean.} L'énoncé encadré est \texttt{diff\_of\_closure} : pour un morphisme additif $f : M \to G$ dont l'image engendre $G$, tout élément de $G$ est un $f(a) - f(b)$. La preuve est la récurrence de la démonstration ci-dessus, sur les éléments du sous-groupe engendré. Le cas du groupe complété d'un monoïde commutatif, \texttt{GrothendieckAddGroup} dans mathlib, est \texttt{grothendieckAddGroup\_diff}. Les noms sont ceux de \href{\leanblobmain/Redaction596.lean}{\texttt{Redaction596.\allowbreak{}lean}}. Les modules, leurs classes, les suites exactes et le groupe $K(\underline C)$ lui-même ne sont pas formalisés : le passage de l'énoncé abstrait à la proposition~3\,a) est celui qu'on vient d'écrire, non vérifié.

Ne sont pas démontrés ici : la proposition~1, propriété universelle de $K(\underline C)$ ; la proposition~2, fonctorialité par restriction et extension des scalaires ; la proposition~3\,b) ; la suite du §~10, conséquences des théorèmes de Krull--Remak--Schmidt et de Jordan--Hölder, résolutions et structures multiplicatives.

\section[Rédaction Bourbaki \redac{596} : trois hypothèses manquantes]{Rédaction Bourbaki \redac{596} : trois énoncés auxquels manque une hypothèse\protect\verifie{Redaction596Characters}}\label{sec:bourbaki-596-hypotheses}

La rédaction \redac{596} (section~\ref{sec:bourbaki-596}) a paru, longtemps après, comme chapitre~VIII d'\emph{Algèbre} (Springer, 2012). Lue en entier, ses 289 pages sont pour l'essentiel justes et classiques : Krull--Schmidt, densité, Morita, Wedderburn, radical, algèbres absolument semi-simples, groupe de Brauer, représentations des groupes finis. Trois énoncés y sont faux tels qu'ils sont tapés, chacun parce qu'il omet une hypothèse que sa démonstration utilise. La transcription reste locale ; on n'en cite aucun passage.

\emph{§~7, corollaire de la proposition~10 (p.~111 du fac-similé).} Soient $K$ un corps, $A_0$, $A_1$ deux $K$-algèbres et $S_i$ un $A_i$-module simple. Si $K$ est algébriquement clos, $S_0 \otimes_K S_1$ serait un $(A_0 \otimes_K A_1)$-module simple. La proposition~10 ne traite que de modules de dimension finie, et la preuve du corollaire invoque le théorème de Burnside, qui demande la dimension finie ; le corollaire ne la demande pas.

\emph{§~12, n\textsuperscript{o}\,1, remarque après la formule~(2) (p.~249).} Pour un monoïde $G$ et un corps $k$, l'application $\mathrm{Tr}' : k \otimes_{\mathbf Z} R_k(G) \to FC(G, k)$, qui envoie la classe d'un module sur son caractère, serait un isomorphisme quand $k$ est algébriquement clos ; la référence est laissée en blanc. Il faut que $k[G]$ soit semi-simple, par exemple $G$ fini d'ordre premier à la caractéristique.

\emph{§~12, n\textsuperscript{o}\,2, proposition~4\,(d) (p.~258).} Sur un corps algébriquement clos, un $G$-module $E$ de dimension finie serait simple si et seulement si $\langle E, E\rangle = \dim_k \operatorname{End}_G(E) = 1$. La preuve tire de $\operatorname{End}_G(E) \cong k$ que $E$ est semi-simple, ce qui ne suit pas.

\begin{encadre}
\begin{contrexemple}[aux trois énoncés de la rédaction \redac{596}, tels qu'ils sont tapés]
\begin{enumerate}
\item[a)] Soit $L = \mathbf C(t)$, simple comme module sur lui-même. Le $(L \otimes_{\mathbf C} L)$-module $L \otimes_{\mathbf C} L$ n'est pas simple : le noyau de la multiplication $L \otimes L \to L$ est un idéal propre, qui contient $t \otimes 1 - 1 \otimes t \neq 0$.
\item[b)] Soient $G = \mathbf Z/2$, d'élément non trivial $g$, et $k$ algébriquement clos de caractéristique~$2$. Tout caractère vérifie $\chi(g) = \chi(1)$ ; la fonction centrale qui vaut $1$ en $1$ et $0$ en $g$ n'est donc pas combinaison de caractères, et $\mathrm{Tr}'$ n'est pas surjectif.
\item[c)] Soit $k$ une clôture algébrique de $\mathbf F_3$, et $G$ le groupe des matrices triangulaires supérieures inversibles à coefficients dans $\mathbf F_3$, d'ordre~$12$, engendré par $A = \operatorname{diag}(1, -1)$ et $N = \left(\begin{smallmatrix} 1 & 1 \\ 0 & 1 \end{smallmatrix}\right)$, agissant sur $E = k^2$. Seules les homothéties commutent à $A$ et à $N$, donc $\langle E, E\rangle = 1$ ; mais la droite $k e_1$ est stable, et $E$ n'est pas simple.
\end{enumerate}
\end{contrexemple}
\end{encadre}

\begin{proof}
a) Pour $x \in L$ hors de $\mathbf C$, $1$ n'est pas dans la droite $\mathbf C x$ ; il existe une forme linéaire $f$ sur $L$ nulle en $x$ et non nulle en $1$. L'application $\mathrm{id} \otimes f : L \otimes L \to L$ envoie $x \otimes 1$ sur $f(1)\,x \neq 0$ et $1 \otimes x$ sur $0$. Le noyau de la multiplication contient donc $x \otimes 1 - 1 \otimes x \neq 0$, et il ne contient pas $1$.

b) Si $\rho(g)^2 = 1$, alors $(\rho(g) - 1)^2 = \rho(g)^2 - 2\rho(g) + 1 = 0$ en caractéristique~$2$ ; l'endomorphisme nilpotent $\rho(g) - 1$ est de trace nulle, d'où $\chi(g) = \chi(1)$. Les fonctions $f$ telles que $f(g) = f(1)$ forment un sous-espace qui contient les caractères et pas la fonction indiquée.

c) Si $M$ commute à $A$, ses coefficients hors diagonale vérifient $2m_{12} = 2m_{21} = 0$, donc sont nuls puisque $2 \neq 0$ ; s'il commute à $N$, ses deux coefficients diagonaux sont égaux. Enfin $A e_1 = N e_1 = e_1$.
\end{proof}

\noindent Dans b) et c), la caractéristique divise l'ordre du groupe. Quand $k[G]$ est semi-simple (théorème de Maschke), les deux énoncés du §~12 sont justes ; c'est sans doute le cas que la page a en vue.

\begin{sloppypar}
\noindent\emph{Vérification en Lean.} a) est \texttt{tmul\_one\_ne\_one\_tmul} et \texttt{not\_isSimpleModule\_tensor}, pour toute extension $L \neq K$, puis \texttt{cor\_prop10\_fails} pour $\mathbf C(t)$ ; b) est \texttt{character\_eq\_of\_sq\_eq\_one} et \texttt{single\_one\_not\_mem\_span\_characters}, pour toute famille de représentations de dimension finie ; c) est \texttt{eq\_scalar\_of\_commute}, \texttt{A\_mulVec\_e1}, \texttt{N\_mulVec\_e1} et \texttt{prop4d\_fails}, dans \href{\leanblobmain/Redaction596Characters.lean}{\texttt{Redaction596\allowbreak{}Characters.\allowbreak{}lean}}. La finitude du groupe engendré par $A$ et $N$, et le fait que $R_k(G)$ s'identifie au groupe libre sur les modules simples, ne sont pas formalisés.
\end{sloppypar}

\emph{Autres écarts relevés.} À l'exemple~1 du §~4, n\textsuperscript{o}\,3 (p.~54), les sous-$M_n(A^{\circ})$-modules de $A_s^n$ correspondent aux idéaux à droite de $A$ : le module est artinien ou noethérien quand $A_d$ l'est, et non $A_s$ comme il est tapé ; un anneau artinien à droite qui ne l'est pas à gauche, comme $\left(\begin{smallmatrix} \mathbf Q & \mathbf R \\ 0 & \mathbf R \end{smallmatrix}\right)$, sépare les deux conditions. Cet écart est formalisé à la section~\ref{sec:bourbaki-596-modules}, avec quatre autres. La lecture a aussi relevé des fautes dans des démonstrations dont l'énoncé reste juste (un signe dans les produits croisés du §~9, une formule de Newton au §~12) ; elles ne sont pas reprises ici.

\litterature{les énoncés justes sont classiques : Curtis et Reiner, \emph{Representation Theory of Finite Groups and Associative Algebras} (1962) ; Serre, \emph{Représentations linéaires des groupes finis} (1967). Le texte imprimé de 2012 n'a pas été consulté : on ne sait pas si les trois énoncés y sont corrigés.}

\section[Rédaction Bourbaki \redac{387} : l'indice des puissances $n$-ièmes]{Rédaction Bourbaki \redac{387} : une formule pour l'indice des puissances $n$-ièmes des unités, fausse pour $\mathbf Q_2$\protect\verifie{Redaction387}}\label{sec:bourbaki-387}

La rédaction \redac{387}, \q{Algèbre commutative. Chapitre VI. Valuations}, des archives Bourbaki (Archives Henri Poincaré, \ahp{942}) est datée de septembre 1962 ; c'est une refonte par Dieudonné du chapitre qui parut en 1964. Son §~9 traite des corps localement compacts. L'exercice~3 de ce paragraphe (p.~67 du fac-similé) donne, pour un tel corps $K$ et le groupe $U$ des unités de son anneau d'entiers, une formule pour l'indice $[U : U^n]$ du sous-groupe des puissances $n$-ièmes. La formule y est tapée comme le produit $\operatorname{mod}_K(n)\cdot \operatorname{Card}({}_nU)$ du module de $n$ et du nombre de racines $n$-ièmes de l'unité. Sa transcription reste locale pour l'instant ; on n'en cite ici aucun passage.

\emph{La formule juste} est le quotient : $[U : U^n] = \operatorname{Card}({}_nU)/\operatorname{mod}_K(n)$ (pour $\mathbf Q_p$, Serre, \emph{Cours d'arithmétique}, chap.~II ; en général, Neukirch, \emph{Algebraic Number Theory}, chap.~II, §~5). Pour $K = \mathbf Q_2$ et $n = 2$, le module de $2$ est $|2|_2 = 1/2$, et les racines carrées de l'unité sont $\pm 1$. Le quotient vaut $2/(1/2) = 4$, le produit $(1/2) \cdot 2 = 1$.

\begin{encadre}
\begin{contrexemple}[à l'exercice~3 du §~9 de la rédaction \redac{387}, tel qu'il est tapé]
Pour $K = \mathbf Q_2$ et $n = 2$, la formule tapée donne $[U : U^2] = 1$, c'est-à-dire que toute unité de $\mathbf Z_2$ serait un carré. Or $-1$ n'est pas un carré dans $\mathbf Z_2$.
\end{contrexemple}
\end{encadre}

\begin{proof}
$|2|_2 = 1/2$, et l'équation $x^2 = 1$ n'a que les solutions $\pm 1$ dans l'anneau intègre $\mathbf Z_2$, qui sont distinctes. Si $x^2 = -1$ dans $\mathbf Z_2$, la réduction modulo~$4$ donnerait $\bar x^2 = 3$ dans $\mathbf Z/4$, où les carrés sont $0$ et $1$.
\end{proof}

\begin{sloppypar}
\noindent\emph{Vérification en Lean.} \texttt{norm\_two} donne $\|2\| = 1/2$ dans $\mathbf Q_2$, \texttt{sq\_eq\_one\_iff} et \texttt{one\_ne\_neg\_one} les deux racines carrées de l'unité, \texttt{not\_isSquare\_neg\_one} le fait que $-1$ n'est pas un carré (par réduction modulo~$4$), et \texttt{typed\_formula\_fails} le contre-exemple, dans \href{\leanblobmain/Redaction387.lean}{\texttt{Redaction387.\allowbreak{}lean}}. L'indice exact $4$ n'est pas formalisé.
\end{sloppypar}

\litterature{la formule juste est classique ; la forme tapée ne l'est pas. L'exercice imprimé (\emph{Algèbre commutative}, chap.~VI, 1964) n'a pas été consulté : on ne sait pas si l'erreur y subsiste.}

\section[Rédaction Bourbaki \redac{397} : divisible et injectif sur un anneau principal]{Rédaction Bourbaki \redac{397} : sur un anneau principal, divisible équivaut à injectif\protect\verifie{Redaction397}}\label{sec:bourbaki-397}

La rédaction \redac{397}, \q{Compléments au chapitre VII d'Algèbre. Dualité des modules de longueur finie sur un anneau principal}, des archives Bourbaki (Archives Henri Poincaré, \ahp{948}) est datée d'avril 1963. Elle définit les modules injectifs, démontre le critère de Baer (sa proposition~2) et en tire un corollaire : sur un anneau principal, un module est injectif si et seulement s'il est divisible ; puis un \q{corollaire du corollaire} : tout quotient d'un module injectif est injectif. Sa transcription reste locale pour l'instant ; on n'en cite ici aucun passage.

\emph{L'idée.} Sur un anneau principal $R$, un idéal est engendré par un seul élément $a$. Prolonger à $R$ une application linéaire $g$ définie sur l'idéal $(a)$, c'est choisir l'image $q$ de $1$, sous la seule contrainte $a\,q = g(a)$, ce qui revient à diviser $g(a)$ par $a$ dans $Q$. Le critère de Baer ramène l'injectivité à ces prolongements depuis les idéaux.

\begin{encadre}
\begin{enonce}[Bourbaki, rédaction \redac{397}, corollaire de la proposition~2 et son corollaire]
Soient $R$ un anneau commutatif et $Q$ un $R$-module.
\begin{enumerate}
\item[a)] Si $R$ est principal et si $Q$ est divisible (pour tout $r \neq 0$, l'application $q \mapsto rq$ de $Q$ dans $Q$ est surjective), alors $Q$ est injectif.
\item[b)] Si $R$ est intègre et si $Q$ est injectif, alors $Q$ est divisible.
\item[c)] Si $R$ est principal et intègre, tout quotient d'un $R$-module injectif est injectif.
\end{enumerate}
\end{enonce}
\end{encadre}

\begin{proof}
a) Par le critère de Baer, il suffit de prolonger à $R$ toute application linéaire $g$ définie sur un idéal $(a)$. Si $a = 0$, on prend l'application nulle. Sinon, on choisit $q$ tel que $a q = g(a)$ ; l'application $r \mapsto r q$ prolonge $g$, car $g(c a) = c\,g(a) = (c a)\, q$.

b) Soient $r \neq 0$ et $q \in Q$. Comme $R$ est intègre, $c \mapsto c r$ est une bijection de $R$ sur l'idéal $(r)$, et $c r \mapsto c q$ est une application linéaire de $(r)$ dans $Q$. Elle se prolonge à $R$ ; si $x$ est l'image de $1$, alors $r x$ est l'image de $r$, c'est-à-dire $q$.

c) Si $Q$ est injectif, il est divisible par b), et un quotient d'un module divisible est divisible ; il est donc injectif par a).
\end{proof}

\noindent Le point a) n'utilise pas l'intégrité : il vaut sur tout anneau commutatif principal.

\begin{sloppypar}
\noindent\emph{Vérification en Lean.} a) est \texttt{baer\_of\_divisible} et \texttt{injective\_of\_divisible}, b) \texttt{divisible\_of\_injective}, l'équivalence \texttt{injective\_iff\_divisible}, et c) \texttt{injective\_quotient}, dans \href{\leanblobmain/Redaction397.lean}{\texttt{Redaction397.\allowbreak{}lean}}. Le critère de Baer de mathlib demande que $R$ et $Q$ soient dans le même univers, et la formalisation garde cette restriction.
\end{sloppypar}

\litterature{retrouvée. Le critère est de Baer (1940) ; l'équivalence entre divisible et injectif sur un anneau principal est classique (Lam, \emph{Lectures on Modules and Rings}, §~3). mathlib n'en avait que le cas de $\mathbf Z$ (\texttt{Module.Baer.of\_divisible}) ; le cas général est proposé pour mathlib, et l'édition ne le donne pas pour nouveau.}

\section[Rédaction Bourbaki \redac{355} : premiers associés et extension plate]{Rédaction Bourbaki \redac{355} : les premiers associés d'un module plat réduit modulo un premier\protect\verifie{Redaction355}}\label{sec:bourbaki-355}

La rédaction \redac{355} des archives Bourbaki (Archives Henri Poincaré, \ahp{913}), datée de janvier 1961, est le second état d'un texte sur les idéaux premiers associés et l'extension des scalaires, dont le contenu se retrouve au §~2 du chapitre~IV d'\emph{Algèbre commutative} (1961). Le premier état (rédaction 352) traitait l'extension $A \to B$ ; la préface tapée du second dit que le rédacteur doit à Grothendieck l'idée de remplacer l'anneau $B$ par un $B$-module $F$ plat sur $A$. Sa transcription reste locale pour l'instant ; on n'en cite ici aucun passage.

\emph{L'énoncé.} Soient $\pi : A \to B$ un homomorphisme d'anneaux commutatifs, $F$ un $B$-module plat sur $A$, et $\mathfrak p$ un idéal premier de $A$. La proposition~1 dit, sans hypothèse noethérienne, que tout idéal premier associé au $B$-module $F/\mathfrak pF$ est au-dessus de $\mathfrak p$. C'est le premier pas de la proposition~2, qui calcule les premiers associés de $E \otimes_A F$.

\emph{L'intuition.} Le module $F/\mathfrak pF$ est $(A/\mathfrak p) \otimes_A F$, plat sur l'anneau intègre $A/\mathfrak p$, donc sans torsion : un élément de $A$ hors de $\mathfrak p$ n'y tue rien. Un premier associé, qui est fait d'éléments qui tuent quelque chose, ne contient donc pas d'autre élément de $A$ que ceux de $\mathfrak p$, et il les contient tous, puisque $\mathfrak p$ tue $F/\mathfrak pF$.

\begin{encadre}
\begin{enonce}[Bourbaki, rédaction \redac{355}, proposition~1]
Soient $\pi : A \to B$ un homomorphisme d'anneaux commutatifs, $F$ un $B$-module plat sur $A$, $\mathfrak p$ un idéal premier de $A$, et $\mathfrak P$ un idéal premier de $B$ associé au $B$-module $F/\mathfrak pF$, c'est-à-dire premier et égal à la racine de l'annulateur d'un élément. Alors $\pi^{-1}(\mathfrak P) = \mathfrak p$.
\end{enonce}
\end{encadre}

\begin{proof}
Soit $\bar x \in F/\mathfrak pF$ avec $\mathfrak P = \sqrt{\operatorname{Ann}_B(\bar x)}$. Si $a \in \mathfrak p$, alors $\pi(a)\bar x = a\bar x = 0$, d'où $a \in \pi^{-1}(\mathfrak P)$. Soit réciproquement $a \in \pi^{-1}(\mathfrak P)$ : il existe $n$ tel que $a^n \bar x = 0$. Si $a \notin \mathfrak p$, alors $a^n \notin \mathfrak p$, et la multiplication par $a^n$ est injective sur $A/\mathfrak p$, qui est intègre ; elle le reste sur $(A/\mathfrak p) \otimes_A F = F/\mathfrak pF$, puisque $F$ est plat. Donc $\bar x = 0$, son annulateur est $B$ et $\mathfrak P = B$, ce qui est exclu.
\end{proof}

\noindent La définition des premiers associés par la racine de l'annulateur, celle de mathlib, coïncide avec la définition usuelle pour un anneau $B$ noethérien ; la rédaction ne suppose rien, et la preuve vaut pour les deux.

\begin{sloppypar}
\noindent\emph{Vérification en Lean.} Le pas de platitude est \texttt{mem\_smul\_top\_of\_smul\_mem}, et la proposition \texttt{comap\_eq\_of\_isAssociatedPrime}, énoncée pour tout sous-$B$-module $N$ de $F$ dont le sous-$A$-module sous-jacent est $\mathfrak pF$ ; \texttt{smulTop\_restrictScalars} montre que $(\mathfrak pB)F$ en est un, dans \href{\leanblobmain/Redaction355.lean}{\texttt{Redaction355.\allowbreak{}lean}}. La proposition~2 n'est pas formalisée : l'égalité, pour $A$ noethérien, demande la décomposition primaire des modules.
\end{sloppypar}

\litterature{retrouvé : Bourbaki, \emph{Algèbre commutative}, chap.~IV, §~2 (numéro du paragraphe non vérifié sur le volume) ; Matsumura, \emph{Commutative Ring Theory}, théorème~23.2 ; EGA~IV$_2$, 3.3.1. mathlib a les premiers associés mais pas leur comportement par extension plate ; l'énoncé est proposé pour mathlib, et l'édition ne le donne pas pour nouveau. L'intérêt de la rédaction est historique : la préface situe en janvier 1961, entre les deux états, la généralisation due à Grothendieck.}



\section*{Déclaration sur l'usage de l'intelligence artificielle}\addcontentsline{toc}{section}{Déclaration sur l'usage de l'intelligence artificielle}

Les énoncés démontrés ici sont lus sur des transcriptions des rédactions produites par des modèles de langage d'Anthropic (Claude) sous des règles éditoriales écrites par l'auteur ; ces transcriptions n'ont pas été relues en entier par une personne. Les preuves formelles ont été écrites en Lean~4 par Claude Opus~5.5, sous la direction de l'auteur, et Lean en vérifie mécaniquement la correction contre la bibliothèque mathlib. Le même modèle en a tiré les démonstrations rédigées de ce document. L'auteur a relu le choix des énoncés, les corrections apportées aux pages et le texte, et en assume l'entière responsabilité. Lean garantit qu'un énoncé tient sous les hypothèses écrites. Il ne garantit ni que l'énoncé rend fidèlement la page, ni qu'il est nouveau.

\begin{thebibliography}{99}\small
\bibitem{briand2004} E.~Briand, « \textenglish{When is the algebra of multisymmetric polynomials generated by the elementary multisymmetric polynomials?} », \emph{Beiträge zur Algebra und Geometrie} 45 (2004), n\textsuperscript{o}\,2, 353--368.
\bibitem{bbi2001} D.~Burago, Yu.~Burago et S.~Ivanov, \emph{A Course in Metric Geometry}, Graduate Studies in Mathematics 33, American Mathematical Society, Providence, 2001 ; doi:10.1090/gsm/033.
\bibitem{freyd1964} P.~Freyd, \emph{Abelian Categories. An Introduction to the Theory of Functors}, Harper \& Row, New York, 1964.
\bibitem{hartshorne1977} R.~Hartshorne, \emph{Algebraic Geometry}, Graduate Texts in Mathematics 52, Springer, New York, 1977 ; doi:10.1007/978-1-4757-3849-0.
\bibitem{johnstone2002} P.~T.~Johnstone, \emph{Sketches of an Elephant. A Topos Theory Compendium}, vol.~1, Oxford Logic Guides 43, Oxford University Press, 2002.
\bibitem{lean4} L.~de Moura et S.~Ullrich, « The Lean~4 theorem prover and programming language », dans \emph{Automated Deduction -- CADE~28}, Lecture Notes in Computer Science 12699, Springer, 2021, 625--635 ; doi:10.1007/978-3-030-79876-5\_37.
\bibitem{mathlib2020} The mathlib Community, « The Lean mathematical library », dans \emph{Proceedings of the 9th ACM SIGPLAN International Conference on Certified Programs and Proofs (CPP~2020)}, ACM, 2020, 367--381 ; doi:10.1145/3372885.3373824.
\bibitem{mitchell1965} B.~Mitchell, \emph{Theory of Categories}, Pure and Applied Mathematics 17, Academic Press, New York, 1965.
\bibitem{mo423444} Question 423444, \emph{MathOverflow}, 2022, \url{https://mathoverflow.net/questions/423444}.
\bibitem{rydh2007} D.~Rydh, « A minimal set of generators for the ring of multisymmetric functions », \emph{Annales de l'Institut Fourier} 57 (2007), n\textsuperscript{o}\,6, 1741--1769 ; doi:10.5802/aif.2312 ; arXiv:0710.0470.
\bibitem{sga4-1} M.~Artin, A.~Grothendieck et J.-L.~Verdier, \emph{Théorie des topos et cohomologie étale des schémas (SGA~4)}, tome~1, Lecture Notes in Mathematics 269, Springer, Berlin, 1972 ; doi:10.1007/BFb0081551.
\bibitem{spanier1966} E.~H.~Spanier, \emph{Algebraic Topology}, McGraw-Hill, New York, 1966.
\bibitem{stacks} The Stacks Project Authors, \emph{The Stacks Project}, \url{https://stacks.math.columbia.edu}.
\bibitem{mo354568} Question 354568, \emph{MathOverflow}, 2020, \url{https://mathoverflow.net/questions/354568}.
\bibitem{neukirch1999} J.~Neukirch, \emph{Algebraic Number Theory}, Grundlehren der mathematischen Wissenschaften 322, Springer, Berlin, 1999 ; doi:10.1007/978-3-662-03983-0.
\bibitem{raoult2019} J.-C.~Raoult, « \textenglish{Semi-cartesian squares and the snake lemma} », \emph{International Journal of Mathematics and Computer Science} 14 (2019), n\textsuperscript{o}\,1, 121--133 ; arXiv:1712.01637.
\bibitem{serre1970} J.-P.~Serre, \emph{Cours d'arithmétique}, Presses Universitaires de France, Paris, 1970.
\end{thebibliography}

\end{document}
